% oss-plantilla: oss-2026.11
% Preámbulo corporativo fijo. No editar en sd-NN.tex: exportar_latex.py verificar lo compara byte a byte.
% Para cambiar el formato, editar esta plantilla u oss.sty y subir la versión de la primera línea.
\PassOptionsToPackage{unicode}{hyperref}
\PassOptionsToPackage{hyphens}{url}
\documentclass[11pt]{article}
\usepackage{amsmath,amssymb}
\usepackage{iftex}
\ifPDFTeX
  % Prism y otros editores en la nube: pdfLaTeX con Helvetica como familia por defecto.
  \usepackage[T1]{fontenc}
  \usepackage[utf8]{inputenc}
  \usepackage{textcomp}
  \usepackage[scaled=0.92]{helvet}
  \renewcommand{\familydefault}{\sfdefault}
  \DeclareUnicodeCharacter{2264}{\ensuremath{\leq}}
  \DeclareUnicodeCharacter{2265}{\ensuremath{\geq}}
  \DeclareUnicodeCharacter{2260}{\ensuremath{\neq}}
  \DeclareUnicodeCharacter{2248}{\ensuremath{\approx}}
  \DeclareUnicodeCharacter{2190}{\ensuremath{\leftarrow}}
  \DeclareUnicodeCharacter{2192}{\ensuremath{\rightarrow}}
  \DeclareUnicodeCharacter{25CF}{\ensuremath{\bullet}}
  \DeclareUnicodeCharacter{2713}{\ensuremath{\checkmark}}
\else
  % XeLaTeX/LuaLaTeX (compilación oficial local): DejaVu Sans.
  \usepackage{fontspec}
  \defaultfontfeatures{Ligatures=TeX,Scale=MatchLowercase}
  \setmainfont{DejaVu Sans}
  \setsansfont{DejaVu Sans}
  \setmonofont{DejaVu Sans Mono}
\fi
\usepackage[letterpaper,margin=20mm]{geometry}
\usepackage{microtype}
\usepackage{parskip}
\usepackage{xcolor}
\usepackage{longtable,booktabs,array}
\usepackage{calc}
\usepackage{etoolbox}
\makeatletter
\patchcmd\longtable{\par}{\if@noskipsec\mbox{}\fi\par}{}{}
\makeatother
\usepackage{footnotehyper}
\makesavenoteenv{longtable}
\usepackage{graphicx}
\makeatletter
\newsavebox\pandoc@box
\newcommand*\pandocbounded[1]{%
  \sbox\pandoc@box{#1}%
  \Gscale@div\@tempa{\textheight}{\dimexpr\ht\pandoc@box+\dp\pandoc@box\relax}%
  \Gscale@div\@tempb{\linewidth}{\wd\pandoc@box}%
  \ifdim\@tempb\p@<\@tempa\p@\let\@tempa\@tempb\fi
  \ifdim\@tempa\p@<\p@\scalebox{\@tempa}{\usebox\pandoc@box}%
  \else\usebox{\pandoc@box}%
  \fi}
\def\maxwidth{\ifdim\Gin@nat@width>\linewidth\linewidth\else\Gin@nat@width\fi}
\def\maxheight{\ifdim\Gin@nat@height>\textheight\textheight\else\Gin@nat@height\fi}
\makeatother
\setkeys{Gin}{width=\maxwidth,height=\maxheight,keepaspectratio}
\makeatletter
\def\fps@figure{htbp}
\makeatother
\usepackage[normalem]{ulem}
\providecommand{\st}[1]{\sout{#1}}
\usepackage{fancyvrb}
\setlength{\emergencystretch}{3em}
\providecommand{\tightlist}{%
  \setlength{\itemsep}{0pt}\setlength{\parskip}{0pt}}
\providecommand{\passthrough}[1]{#1}
\setcounter{secnumdepth}{-\maxdimen}
\ifPDFTeX
  \usepackage[spanish]{babel}
\else
  \usepackage[bidi=default]{babel}
  \babelprovide[main,import]{spanish}
\fi
\let\LanguageShortHands\languageshorthands
\def\languageshorthands#1{}
\usepackage{bookmark}
\usepackage{xurl}
\urlstyle{same}
\hypersetup{pdflang={es-CL},pdfcreator={Only Simple Solutions - exportar_latex.py}}
\usepackage{oss}
\begin{document}


\ossCover{Metodologías}{Subdocumento 6}

\renewcommand{\contentsname}{Índice general}
\renewcommand{\listfigurename}{Lista de figuras}
\renewcommand{\listtablename}{Lista de tablas}
\renewcommand{\tablename}{Tabla}
\renewcommand{\figurename}{Figura}
\renewcommand{\thefigure}{6.\arabic{figure}}
\renewcommand{\thetable}{6.\arabic{table}}
\setcounter{tocdepth}{4}
\makeatletter
\renewcommand\paragraph{\@startsection{paragraph}{4}{\z@}{2.2ex \@plus 1ex \@minus .2ex}{0.6ex}{\normalfont\normalsize\bfseries}}
\renewcommand\subparagraph{\@startsection{subparagraph}{5}{\z@}{1.8ex \@plus .6ex \@minus .2ex}{0.4ex}{\normalfont\normalsize\bfseries\itshape}}
\renewcommand*\l@figure{\@dottedtocline{1}{1.5em}{2.9em}}
\renewcommand*\l@table{\@dottedtocline{1}{1.5em}{2.9em}}

\makeatother
\tableofcontents
\clearpage
\listoffigures
\clearpage
\listoftables
\clearpage

\section{Capítulo 6 · Introducción a las Metodologías}\label{capitulo-6-introduccion-a-las-metodologias}

Este subdocumento define cómo Only Simple Solutions dirigirá el proyecto
y organizará la construcción, verificación y entrega de la solución de
Multitiendas Ancoa S.A. La gestión del apartado 6.1 integra compromisos
contractuales y coordinación de frentes; el desarrollo del 6.2 produce
incrementos dentro del alcance y de sus condiciones de calidad y
aceptación. Los formularios T-9 y T-10 recogerán el detalle operativo en archivos separados, y el Anexo A reúne las planillas de gestión de interesados. En este subdocumento, las tablas y figuras se numeran 6.n; en su anexo, A.x.y; y en los formularios, por apartado (por ejemplo, T-9.4.2). Toda referencia a otro documento indica su nombre.

La metodología utiliza el gobierno corporativo del Subdocumento 1, el
diagnóstico e interesados del Subdocumento 2, el alcance del Subdocumento 3 y la arquitectura del
Subdocumento 4. Sus decisiones alimentan el cronograma y recursos del Subdocumento 7, los riesgos
del Subdocumento 8, las pruebas del Subdocumento 9 y la organización del Subdocumento 12. La Tabla 6.1 muestra
la correspondencia de control que se empleará para relacionar gestión y
construcción con esos productos.

\begin{longtable}[]{@{}
  >{\raggedright\arraybackslash}p{(\columnwidth - 6\tabcolsep) * \real{0.2500}}
  >{\raggedright\arraybackslash}p{(\columnwidth - 6\tabcolsep) * \real{0.2500}}
  >{\raggedright\arraybackslash}p{(\columnwidth - 6\tabcolsep) * \real{0.2500}}
  >{\raggedright\arraybackslash}p{(\columnwidth - 6\tabcolsep) * \real{0.2500}}@{}}
\caption{Correspondencia de las metodologías y sus evidencias.}\label{tab:6-1}\\
\toprule\noalign{}
\begin{minipage}[b]{\linewidth}\raggedright
\textbf{Ámbito}
\end{minipage} & \begin{minipage}[b]{\linewidth}\raggedright
\textbf{Apartado}
\end{minipage} & \begin{minipage}[b]{\linewidth}\raggedright
\textbf{Control asociado}
\end{minipage} & \begin{minipage}[b]{\linewidth}\raggedright
\textbf{Evidencia de ejecución}
\end{minipage} \\
\midrule\noalign{}
\endfirsthead
\toprule\noalign{}
\begin{minipage}[b]{\linewidth}\raggedright
\textbf{Ámbito}
\end{minipage} & \begin{minipage}[b]{\linewidth}\raggedright
\textbf{Apartado}
\end{minipage} & \begin{minipage}[b]{\linewidth}\raggedright
\textbf{Control asociado}
\end{minipage} & \begin{minipage}[b]{\linewidth}\raggedright
\textbf{Evidencia de ejecución}
\end{minipage} \\
\midrule\noalign{}
\endhead
\bottomrule\noalign{}
\multicolumn{4}{@{}l}{\rule{0pt}{3.2ex}\textbf{Fuente:} Elaboración propia.} \\
\endlastfoot
Dirección y cambios & 6.1 & Gobierno, líneas base y control de cambios &
Actas, líneas base y decisiones \\
Construcción y entrega & 6.2 & DevSecOps, seguridad, pruebas y
aceptación & Versiones, pruebas y conformidades \\
Convivencia e implantación & 6.1 y 6.2 & Marchas blancas, pasos a
producción y convivencia & Conciliación, preparación y retorno \\
Operación y conocimiento & 6.1.3.4, 6.1.7.1 y 6.2.7 & Continuidad, métricas y transferencia & Servicio, métricas y transferencia \\
\end{longtable}

La correspondencia obliga a que una decisión de metodología tenga un
producto comprobable y un responsable. Los registros de trabajo y
pruebas respaldarán la verificación; la aceptación contractual seguirá
correspondiendo a la contraparte y al procedimiento exigidos.

Este capítulo describe las metodologías que regirán la dirección del
proyecto y el desarrollo de software para Multitiendas Ancoa S.A. La
propuesta debe coordinar la modernización progresiva de un ecosistema de
nueve plataformas y catorce interfaces heredadas, combinando reemplazo,
incorporación, integración y coexistencia según el sistema, sin
interrumpir la operación de las 22 tiendas ni de los dos centros de
distribución. A ello se suma la obligación de mantener una separación
controlada y auditable entre el negocio retail y la filial emisora.

La ejecución se desarrollará dentro de un contrato de 56 meses, con
producción de la Etapa 1 desde el mes 16, producción de la Etapa 2 e
inicio de la operación contractual desde el mes 21 y continuidad hasta
el mes 56. La planificación deberá absorber las cinco ventanas de
congelamiento comercial y la fecha variable del evento anual, sin
trasladar los hitos comprometidos (Bases Administrativas, 2026, art. 17;
Bases Técnicas del Caso 09, 2026, núm. 13.2).

Only Simple Solutions utilizará un enfoque híbrido. La dirección del
proyecto aplicará prácticas de PMBOK para mantener líneas base,
trazabilidad, control de cambios, gestión de interesados,
comunicaciones, adquisiciones y dependencias externas. La construcción
de software se organizará de forma incremental con Scrum, mientras que
Kanban se utilizará para trabajos de flujo continuo, habilitación,
plataforma, soporte e incidencias. La planificación detallada se
revisará conjuntamente cada cuatro sprints, como cadencia de
coordinación; el calendario integrado relacionará esas ocho semanas con
los meses e hitos contractuales. Esta revisión no crea nuevos
compromisos de alcance ni modifica fechas contractuales.

La metodología de desarrollo incorporará DevSecOps, integración y
entrega continuas, infraestructura como código y automatización de
pruebas. Cada incremento deberá producir evidencia suficiente para
demostrar el cumplimiento de los requisitos aplicables y deberá
integrarse al procedimiento formal de pruebas, aceptación y paso a
producción definido para el proyecto.

\subsection{6.1 Metodología de Gestión de Proyectos}\label{6-1-metodologia-de-gestion-de-proyectos}

La gestión combinará control formal e iteración. Las líneas base de
alcance y cronograma, los hitos, los criterios de aceptación, la gestión
de riesgos, los cambios y la relación contractual se administrarán
mediante mecanismos formales. En paralelo, los equipos podrán refinar el
detalle funcional y técnico mediante ciclos cortos, siempre dentro de
los límites aprobados y sin alterar las obligaciones de continuidad,
seguridad, segregación de datos y fechas comprometidas.

\subsubsection{6.1.1. Enfoque de gestión y adaptación del PMBOK}\label{6-1-1-enfoque-de-gestion-y-adaptacion-del-pmbok}

Only Simple Solutions aplicará una metodología híbrida basada en la Guía del PMBOK, sexta edición, complementada con Scrum para la construcción incremental de software y Kanban para la gestión del trabajo de habilitación y soporte. Esta combinación permitirá mantener el control formal de los compromisos contractuales y adaptar el detalle de ejecución a los resultados de las pruebas y la participación de los usuarios. Su aplicación responde a la necesidad de coordinar plataformas, proveedores e infraestructura mientras Ancoa mantiene su operación y la separación entre el negocio retail y el emisor de crédito fiscalizado.

La oficina de gestión, dirigida por el Jefe de Proyecto, planificará y
controlará formalmente el acta de constitución, el plan para la
dirección del proyecto, el alcance, la EDT, el cronograma, la validación
y control del alcance, el seguimiento del desempeño, el control
integrado de cambios y el cierre de fases. La gestión de adquisiciones y
recursos se coordinará con las dependencias técnicas y los plazos de
suministro, mientras que riesgos, calidad, comunicaciones e interesados
se administrarán como partes integradas del mismo plan de dirección.

La adaptación del PMBOK conservará los procesos necesarios para
controlar los compromisos de esta solución. La oficina de gestión
relacionará cada proceso con sus entradas, responsable, producto y
evidencia. Los planes subsidiarios se revisarán conjuntamente cuando
cambie una dependencia; la metodología no se limitará a nombrar áreas de
conocimiento (PMI, 2017).

La gestión del riesgo seguirá ISO 31000 y mantendrá un registro vivo,
revisado en cada Comité de Proyecto. Cada riesgo tendrá responsable,
disparador, respuesta y plazo; las respuestas que afecten una línea base
se tramitarán mediante control integrado de cambios. El Plan de Riesgos
del Subdocumento 8 y el Formulario T-16 desarrollarán su detalle (Bases
Técnicas Transversales, 2026, RT-19.04).

La construcción de servicios e integraciones se organizará con Scrum en
sprints de dos semanas. Las iteraciones permitirán validar reglas de
negocio con usuarios, comprobar contratos de integración y detectar
tempranamente problemas de datos, seguridad o funcionamiento. El product
owner ordenará el backlog dentro del alcance aprobado; los developers
planificarán y ejecutarán el trabajo del sprint; y el scrum master
facilitará la aplicación del marco y la resolución de impedimentos. Cada
incremento deberá cumplir la definición de terminado y aportar evidencia
vinculada con los requisitos y entregables. La revisión del sprint no
sustituirá la aceptación contractual ni la autorización de paso a
producción.

Las solicitudes de ambientes, tareas de plataforma e incidencias se
gestionarán mediante Kanban, porque su llegada y prioridad pueden variar
durante la ejecución. Cada flujo tendrá estados, criterios de entrada y
salida, responsables y límites de trabajo en curso según la capacidad
disponible. El responsable del flujo revisará bloqueos, antigüedad del
trabajo y tiempos de resolución. El líder de operación coordinará la
atención de incidencias que afecten al servicio, mientras el jefe de
proyecto resolverá su impacto sobre los demás frentes. Este mecanismo
complementará la programación integrada y permitirá atender necesidades
variables sin incorporar automáticamente nuevos compromisos de alcance.

Los responsables del flujo definirán qué constituye una solicitud,
cuándo empieza y termina, sus estados, políticas de prioridad y control
del trabajo en curso. También establecerán una expectativa de servicio
con plazo y probabilidad, inicialmente estimada y revisada con datos de
tiempo de ciclo. Este pronóstico de gestión no sustituirá los tiempos
contractuales de atención de incidentes (Kanban Guides, 2025).

Para establecer la línea base de alcance, el jefe de proyecto
consolidará el enunciado de alcance, la EDT y su diccionario, con
entregables, exclusiones y criterios de aceptación. Para establecer la
línea base de cronograma, integrará actividades, duraciones,
dependencias, recursos e hitos, comprobando que la programación permita
cumplir las dos etapas y sus solapamientos. El líder funcional y las
contrapartes usuarias verificarán el alcance; los líderes técnicos
comprobarán dependencias y capacidad; y el Comité Ejecutivo formalizará
las decisiones que excedan las atribuciones operativas y requieran su
aprobación. Las versiones aprobadas y sus antecedentes se conservarán en
el espacio colaborativo del proyecto.

La actualización del avance y del pronóstico mantendrá visibles las
diferencias respecto de las líneas base. Los ajustes de detalle del
backlog podrán acordarse cuando no alteren requisitos, entregables,
calidad ni fechas comprometidas. Toda modificación de las líneas base
requerirá una solicitud formal, análisis de impacto y aprobación previa
del Comité Ejecutivo; las modificaciones contractuales exigirán además
acuerdo escrito de ambas partes. El procedimiento se desarrolla en
6.1.3. Ninguna actualización podrá modificar el cronograma obligatorio,
con producción de la Etapa 1 desde el mes 16, producción de la Etapa 2 e
inicio de la operación contractual en el mes 21, y término del contrato
en el mes 56. Las desviaciones se atenderán mediante planes de
recuperación que preserven esos plazos y las obligaciones de seguridad y
separación de datos.

\subsubsection{6.1.2. Gobiernos, roles y decisiones}\label{6-1-2-gobiernos-roles-y-decisiones}

El jefe de proyecto de Only Simple Solutions dirigirá la ejecución, integrará los frentes y coordinará sus dependencias. Tendrá dedicación exclusiva durante la implementación y facultades para comprometer a la empresa en materias de ejecución. Por parte de Ancoa, el administrador del contrato supervisará el cumplimiento contractual, aprobará los estados de pago y gestionará modificaciones y sanciones; la contraparte técnica validará entregables, aprobará avances técnicos, coordinará pruebas y emitirá las conformidades de los hitos. El patrocinador representará las prioridades estratégicas del cliente en el Comité Ejecutivo. Estas atribuciones se mantendrán diferenciadas.

La estructura de Only Simple Solutions considerará, como mínimo, jefe de
proyecto, arquitecto de solución, encargado de seguridad de la
información, líder de datos, líder de desarrollo, líder funcional, líder
de calidad y pruebas, líder de integración, líder de operación/SRE y
líder de implantación y gestión del cambio. Las responsabilidades ágiles
de product owner y scrum master complementarán esta estructura, pero no
otorgarán autoridad para aprobar cambios contractuales, aceptar hitos ni
autorizar pasos a producción.

La asignación de personas y dedicaciones cubrirá esos roles durante los
períodos exigidos. El líder de desarrollo y el líder funcional tendrán
dedicación completa durante implementación; datos e integración serán
permanentes en esa fase, el líder de operación/SRE participará desde el mes 6, en que el hito H3 entrega los ambientes, y el líder de implantación desde el mes 8. La composición por
equipo, las suplencias y cualquier acumulación de funciones se
documentarán en el Subdocumento 12, el Formulario T-8 y la nivelación
T-15, sin contabilizar dos veces la capacidad de una persona (Bases
Técnicas Transversales, 2026, núm. 19.2).

El Comité Ejecutivo sesionará mensualmente con el patrocinador del
cliente, la gerencia de Only Simple Solutions y el administrador del
contrato. Resolverá decisiones estratégicas, asuntos escalados, cambios
de alcance y riesgos mayores. El Comité de Proyecto sesionará
quincenalmente con la contraparte técnica y el jefe de proyecto;
revisará cronograma, entregables, desviaciones, riesgos y decisiones
operativas. Los cambios formales requerirán aprobación previa del Comité
Ejecutivo y acuerdo escrito de ambas partes cuando modifiquen el
contrato.

El Comité de Arquitectura sesionará mensualmente con el arquitecto de
solución, el encargado de seguridad y los referentes técnicos de Ancoa
para aprobar decisiones de arquitectura y revisar deuda y riesgos
técnicos. El Comité de Operación sesionará mensualmente desde el mes 13,
con el líder de operación, la mesa de servicio y la contraparte técnica,
para revisar niveles de servicio, incidentes, problemas y mejora
continua. Los equipos de ambas partes realizarán una reunión semanal de
seguimiento de impedimentos y compromisos. Las actividades de Scrum y
Kanban complementarán estas instancias sin sustituirlas.

La contraparte técnica suscribirá las actas de aceptación de los hitos,
con el entregable, evidencia objetiva, trazabilidad y observaciones
resueltas. El cliente dispondrá de diez días hábiles para revisar y
pronunciarse, y Only Simple Solutions dispondrá de diez días hábiles
para subsanar las observaciones. Para cada paso de fase, el jefe de
proyecto reunirá las evaluaciones de calidad, seguridad, datos,
implantación y operación; el Comité de Proyecto revisará la preparación
y la contraparte técnica formalizará la conformidad exigida. El cierre
de una marcha blanca y la habilitación del paso a producción requerirán
el cumplimiento de las condiciones de cierre de la marcha blanca y su
acta de aceptación. Un comité interno ni un product owner podrán
dispensarlas (Bases Administrativas, 2026, arts. 17.3 y 18).

Los líderes de desarrollo, funcional, integración, datos, calidad y
pruebas, seguridad, operación e implantación prepararán las evaluaciones
de su disciplina. Cada producto tendrá un product owner identificado,
responsable de ordenar su backlog dentro del mandato del cliente. Los
developers decidirán cómo realizar el trabajo del sprint y el scrum
master facilitará el marco y la resolución de impedimentos. Las
asignaciones y suplencias se documentarán y se mantendrán consistentes
con la estructura del equipo y la nivelación de recursos; los roles
ágiles no otorgarán autoridad para modificar el contrato, aceptar hitos
ni autorizar pasos a producción.

Las decisiones sobre datos y obligaciones del emisor tendrán
responsables explícitos. El líder de datos evaluará autoridad,
integridad y tratamiento de la información; el encargado de seguridad
verificará accesos y segregación; y la contraparte de cumplimiento
designada por Ancoa evaluará las obligaciones del emisor fiscalizado.
Toda decisión que afecte cruces entre ambos negocios seguirá una ruta única, coherente con el Subdocumento 4, que exige aprobación jurídica de cada cruce, y con el Subdocumento 5, que asigna el dominio Frontera a Contraloría y al Negocio Financiero en conjunto. Primero, el líder de datos preparará la ficha de intercambio de X-01 con la finalidad, la base de licitud y los campos mínimos del cruce. Segundo, el encargado de seguridad y la contraparte de cumplimiento designada por Ancoa emitirán sus evaluaciones documentadas. Tercero, el área jurídica de Ancoa validará la base de licitud, y Contraloría y el Negocio Financiero aprobarán en conjunto la ficha; esta aprobación es la única que habilita el cruce. Cuarto, el Comité de Arquitectura verificará que la política de X-01 aplica exactamente la ficha aprobada y registrará la decisión de arquitectura. El Comité Ejecutivo será informado de cada ficha aprobada; si un cruce afecta el alcance o una línea base, se tramitará además por control integrado de cambios. Ante un desacuerdo, el cruce permanecerá denegado y el asunto se escalará al Comité Ejecutivo, que podrá redefinir o repriorizar la necesidad, pero no sustituir la validación jurídica ni la aprobación conjunta. Ninguna prioridad comercial autorizará eliminar la separación entre retail y el emisor. Como mecanismo de resolución, las
decisiones conjuntas requerirán representantes autorizados de ambas
partes y responsables de la materia afectada. El acuerdo será explícito
y quedará en acta; si falta un responsable necesario o existe
desacuerdo, el asunto se mantendrá abierto y se escalará, sin aprobación
por silencio. Los bloqueos operativos se tratarán en el Comité de
Proyecto, las controversias técnicas en el Comité de Arquitectura y los
asuntos contractuales o fuera de sus atribuciones en el Comité
Ejecutivo. Las amenazas a seguridad, integridad de datos o continuidad
se comunicarán inmediatamente. Las actas se emitirán dentro de los dos
días hábiles siguientes, con acuerdos, responsables y plazos.

\subsubsection{6.1.3. Gestión de integración y control de cambios}\label{6-1-3-gestion-de-integracion-y-control-de-cambios}

Only Simple Solutions aplicará los siete procesos de gestión de la integración para mantener la coherencia entre los objetivos de Ancoa, los compromisos aprobados y el trabajo de los equipos. El jefe de proyecto coordinará el desarrollo, infraestructura híbrida, integraciones, datos, seguridad, implantación y operación mediante el plan de dirección integrado. Las decisiones de cada frente se evaluarán considerando sus efectos sobre los demás, preservando la continuidad del servicio y la separación entre retail y el emisor de crédito fiscalizado.

La integración de la gestión vinculará los entregables de la EDT con sus
responsables, dependencias, criterios de aceptación y evidencias. Scrum
y Kanban organizarán la ejecución dentro de este marco; sus ajustes no
modificarán por sí solos el alcance, las líneas base ni los compromisos
contractuales.

\paragraph{6.1.3.1 Desarrollar el acta de constitución del proyecto}\label{6-1-3-1-desarrollar-el-acta-de-constitucion-del-proyecto}

Al inicio del contrato, el jefe de proyecto preparará el acta de constitución con el patrocinador y los representantes autorizados de Ancoa. El acta formalizará el propósito del proyecto, los objetivos, el alcance general, los principales entregables, las restricciones, los riesgos iniciales y las responsabilidades de dirección. También establecerá las facultades del jefe de proyecto y los mecanismos para solicitar decisiones que excedan su autoridad.

Su preparación utilizará los antecedentes de la contratación y del
alcance comprometido. Los supuestos, restricciones y dependencias se
incorporarán al registro correspondiente, identificando su responsable y
el trabajo que condicionan. El acta no sustituirá el contrato ni
permitirá omitir obligaciones; proporcionará una referencia común para
iniciar la planificación y coordinar a ambas partes.

\paragraph{6.1.3.2 Desarrollar el plan para la dirección del proyecto}\label{6-1-3-2-desarrollar-el-plan-para-la-direccion-del-proyecto}

El jefe de proyecto consolidará el plan de dirección integrado con las aportaciones de los responsables técnicos y funcionales. Este reunirá los planes de alcance, cronograma, recursos, calidad, riesgos, comunicaciones, adquisiciones e interesados, junto con las líneas base de alcance, cronograma y costos. Establecerá cómo se ejecutará, supervisará y cerrará el trabajo, y cómo se tramitarán los cambios.

La programación relacionará las dependencias entre entregables. La
construcción de una integración requerirá contratos y responsabilidades
sobre datos definidos; sus pruebas necesitarán ambientes y controles
habilitados; y su implantación dependerá de evidencia de calidad,
preparación de usuarios y capacidad de soporte. El plan identificará el
resultado que habilita cada trabajo, el responsable de entregarlo y la
fecha en que debe estar disponible.

El arquitecto de solución comprobará la coherencia entre los componentes
de nube y locales; el líder de integración coordinará interfaces y
dependencias entre plataformas; el líder de datos verificará migración,
conciliación y autoridad sobre la información; y los responsables de
seguridad, calidad, implantación y operación aportarán sus condiciones
de preparación y aceptación. El jefe de proyecto resolverá la
coordinación de recursos y elevará los desacuerdos a las instancias de
gobierno correspondientes.

Las versiones aprobadas se conservarán en el espacio colaborativo, con
identificación de fecha, responsable y decisión de aprobación. Cada
equipo trabajará sobre la versión vigente. El registro de avance y el
pronóstico se actualizarán sin reemplazar automáticamente las líneas
base, para mantener visibles las desviaciones respecto de los
compromisos originales.

\paragraph{6.1.3.3 Dirigir y gestionar el trabajo del proyecto}\label{6-1-3-3-dirigir-y-gestionar-el-trabajo-del-proyecto}

El jefe de proyecto coordinará la ejecución del plan aprobado y de los cambios autorizados. Los equipos de desarrollo construirán incrementos mediante Scrum; los responsables de plataforma y soporte gestionarán su flujo mediante Kanban; y las adquisiciones, instalaciones, migraciones y pasos a producción seguirán la programación integrada y sus condiciones de autorización.

La reunión semanal de seguimiento comprobará las dependencias entre
frentes. Cada responsable comunicará los resultados obtenidos, las
evidencias disponibles, los impedimentos y los compromisos siguientes.
Cuando una dependencia no esté resuelta, se registrará su efecto y se
acordará una acción con responsable y fecha. El inicio del trabajo
dependiente requerirá verificar sus condiciones de entrada, evitando
avanzar sobre diseños, datos o ambientes que todavía no permitan obtener
un resultado válido.

Durante los frentes simultáneos, el jefe de proyecto distinguirá la
capacidad destinada a construcción de la reservada para marcha blanca y
atención de producción. Una incidencia que requiera participación del
equipo de desarrollo se coordinará con el líder de operación y se
evaluará por su efecto sobre el servicio y los objetivos de ejecución.
La redistribución de trabajo no supondrá una autorización para reducir
controles o desplazar hitos contractuales.

\paragraph{6.1.3.4 Gestionar el conocimiento del proyecto}\label{6-1-3-4-gestionar-el-conocimiento-del-proyecto}

Only Simple Solutions mantendrá el conocimiento necesario para que las decisiones y procedimientos puedan comprenderse y utilizarse por los equipos que construyen, implantan y operan la solución. Las decisiones técnicas se documentarán en el catálogo de decisiones de arquitectura, y los aprendizajes de pruebas, migraciones, incidencias y retrospectivas se incorporarán al registro de lecciones aprendidas.

Cada lección identificará el contexto, el problema observado, la
respuesta aplicada y el resultado. El responsable del frente propondrá
cómo utilizarla en el trabajo siguiente. Por ejemplo, una dificultad
detectada en un ensayo de migración deberá traducirse en ajustes del
procedimiento, verificaciones adicionales o preparación de los
responsables antes del siguiente ensayo.

El conocimiento se transferirá mediante revisión de documentación,
sesiones entre equipos y ejecución acompañada de procedimientos. Los
responsables de operación comprobarán que pueden utilizar los
antecedentes recibidos. El acceso se limitará según la función y el
dominio de información, evitando que la transferencia de conocimiento
exponga datos o expedientes del emisor a personas no autorizadas.

\paragraph{6.1.3.5 Monitorear y controlar el trabajo del proyecto}\label{6-1-3-5-monitorear-y-controlar-el-trabajo-del-proyecto}

El jefe de proyecto consolidará la información de desempeño y la comparará con las líneas base, los criterios de aceptación y los compromisos de cada frente. El seguimiento considerará entregables terminados y verificados, dependencias abiertas, resultados de pruebas, riesgos, incidencias y decisiones pendientes. La existencia de tareas ejecutadas no se considerará suficiente para declarar cumplido un entregable.

La medición mediante valor ganado relacionará el trabajo verificado de
la EDT con su presupuesto y calendario de control. El valor planificado
(PV) será el valor presupuestado del trabajo previsto hasta la fecha de
corte; el valor ganado (EV), el valor presupuestado del trabajo
efectivamente realizado y comprobado; y el costo real (AC), el costo
incurrido por ese trabajo. El jefe de proyecto definirá previamente la
regla de reconocimiento de cada paquete, apoyada en productos y
evidencias, para evitar porcentajes subjetivos de avance (Prasad et al.,
2006).

Los índices se calcularán como SPI = EV/PV y CPI = EV/AC. Si el
denominador es cero, se registrará que el índice no es calculable y su
causa. El reconocimiento de EV para control no equivaldrá por sí solo a
la aceptación contractual ni al derecho a pago. Las métricas de Scrum y
Kanban complementarán el seguimiento de capacidad, bloqueos y entrega
sin sustituir el valor ganado.

El informe mensual presentará los índices, su fecha de corte, evidencia
y explicación de desviaciones, con trazabilidad hacia los paquetes de
trabajo. Los antecedentes financieros se custodiarán y comunicarán por
los canales contractuales autorizados; la Oferta Técnica no divulgará
precios ni montos que permitan inferir la Oferta Económica (Bases
Técnicas Transversales, 2026, RT-19.06 y RT-19.07).

Ante una desviación, el jefe de proyecto analizará su causa y efecto con
los responsables implicados. Las acciones correctivas buscarán recuperar
el desempeño; las preventivas reducirán la posibilidad de
incumplimientos futuros; y las reparaciones de defectos corregirán
resultados no conformes. Toda desviación superior al 10 \% en un hito se
comunicará dentro de los cinco días hábiles siguientes a su detección,
acompañada de un plan de recuperación. Si una acción modifica
compromisos o documentos sujetos a aprobación, se tramitará por control
integrado de cambios antes de ejecutarse.

\paragraph{6.1.3.6 Realizar el control integrado de cambios}\label{6-1-3-6-realizar-el-control-integrado-de-cambios}

El control integrado de cambios evaluará cada modificación como parte del proyecto completo. Una necesidad funcional, un hallazgo técnico, una incidencia o una observación del cliente podrá originar una solicitud, pero su registro no autorizará la ejecución.

Registro de la solicitud. El solicitante describirá el motivo, el
resultado esperado y los entregables afectados. El jefe de proyecto
asignará un identificador y registrará la fecha, el origen, los
antecedentes y el estado de evaluación. La solicitud se vinculará con
los requisitos, observaciones, riesgos o incidencias que la justifican.

Análisis de impacto. El jefe de proyecto coordinará la evaluación de
alcance, plazo, costo, riesgo y niveles de servicio. Los líderes
correspondientes analizarán además arquitectura, interfaces, datos,
seguridad, calidad, recursos, capacitación y operación. Se identificarán
las dependencias afectadas, las alternativas y las consecuencias de
aprobar o rechazar la solicitud. Los cambios que involucren datos u
obligaciones del emisor incorporarán la evaluación del líder de datos,
del encargado de seguridad y de la contraparte de cumplimiento de Ancoa.

Decisión. El Comité de Proyecto revisará los antecedentes y preparará
una recomendación. El Comité Ejecutivo aprobará o rechazará el cambio
antes de su ejecución. Cuando exista una modificación contractual, se
requerirá acuerdo escrito de ambas partes y se respetarán los límites
establecidos en el contrato. No se admitirán modificaciones que alteren
el cronograma obligatorio, el objeto principal, la naturaleza del
contrato o las garantías mínimas.

Las modificaciones contractuales respetarán el límite acumulado del 20
\% del valor original del contrato y las demás condiciones contractuales
aplicables. El registro de cambios identificará el efecto acumulado y la
autorización, sin publicar importes de oferta en este capítulo (Bases
Administrativas, 2026, art. 72).

Actualización y comunicación. Una vez autorizado el cambio, el jefe de
proyecto actualizará las líneas base y los documentos afectados,
conservando las versiones anteriores y el fundamento de la modificación.
Los responsables ajustarán requisitos, decisiones de arquitectura,
backlogs, procedimientos y pruebas según corresponda. La decisión se
comunicará a los equipos y representantes afectados, indicando qué
cambia, desde cuándo y qué acciones deben realizar.

Ejecución y cierre. El responsable designado ejecutará el cambio
autorizado y aportará evidencia de su resultado. El jefe de proyecto
comprobará que se realizó conforme a la decisión y que se actualizaron
los registros relacionados. La aceptación del entregable afectado
seguirá correspondiendo a la contraparte técnica. Las solicitudes
rechazadas conservarán su justificación y se comunicarán al solicitante.

El refinamiento del backlog podrá dividir tareas, aclarar detalles o
ajustar el orden de construcción cuando mantenga los requisitos,
entregables, controles y fechas aprobados. Si el ajuste afecta una línea
base o existe incertidumbre sobre sus consecuencias, se someterá a
evaluación formal. La atención urgente de una incidencia utilizará los
procedimientos de respuesta y reversión autorizados, sin convertir la
urgencia en aprobación de nuevos compromisos.

\paragraph{6.1.3.7 Cerrar el proyecto o fase}\label{6-1-3-7-cerrar-el-proyecto-o-fase}

El jefe de proyecto coordinará el cierre de cada fase comprobando los entregables, la resolución de observaciones, las evidencias de verificación y las actas de aceptación. La contraparte técnica formalizará la aceptación conforme a sus atribuciones. El cierre incluirá la actualización de documentación, el registro de lecciones y la transferencia de los resultados al frente que continúa el trabajo.

En cada marcha blanca se comprobarán, de forma copulativa, las
condiciones que habilitan producción. No existirán incidentes críticos
ni altos abiertos atribuibles a la solución; el volumen real
comprometido y los indicadores de disponibilidad y respuesta se habrán
sostenido durante al menos las últimas cuatro semanas; la conciliación
no tendrá diferencias no explicadas; el personal del cliente habrá
recibido la capacitación y certificación exigidas; y la contraparte
técnica habrá suscrito el acta de aceptación. Una revisión de sprint o
declaración interna de término no podrá sustituirlas (Bases
Administrativas, 2026, art. 17.3).

Si al cierre del período no se cumplen esas condiciones, Only Simple
Solutions aplicará la recuperación y extensión de marcha blanca que
exige el contrato, a su cargo y sin desplazar las fases siguientes ni
dispensar las responsabilidades por atraso. El plan de implantación
deberá prever recursos, dependencias y holguras compatibles con este
procedimiento.

El cierre de implementación del mes 21 se distinguirá del cierre
contractual del mes 56. El primero transferirá ambos alcances a
operación con sus procedimientos, documentación y responsabilidades; el
segundo verificará las obligaciones finales de transferencia y
reversibilidad. Ningún cierre extinguirá responsabilidades que deban
mantenerse durante el período contractual siguiente.

\subsubsection{6.1.4. Gestión de interesados}\label{6-1-4-gestion-de-interesados}

En el proyecto de Multitiendas Ancoa, la gestión de interesados tiene un rol clave para asegurar la viabilidad de la solución. El proyecto modificará las formas de trabajo en 22 tiendas, dos centros de distribución y la filial emisora, y su resultado deberá ser aceptado por áreas cuyos objetivos no siempre coinciden. Por ello, Only Simple Solutions tratará a los interesados como un componente del plan que se identifica, se analiza, se involucra y se monitorea durante todo el contrato.

Se aplicarán los cuatro procesos del área de gestión de interesados del
PMBOK (PMI, 2017), que comprenden identificar a los interesados,
planificar su involucramiento, gestionar su participación y monitorear
su involucramiento.

En esta sección se usan cuatro términos con un significado fijo.

\begin{itemize}
\item
  \textbf{Interesado:} persona, área u organización que puede afectar,
  verse afectada o percibirse afectada por una decisión, actividad o
  resultado del proyecto (PMI, 2017).
\item
  \textbf{Grupo de interés:} conjunto de interesados reunidos según la
  naturaleza de su participación. Organiza el análisis y la gestión,
  pero no es un interesado en sí.
\item
  \textbf{Actor:} interesado que interactúa directamente con los
  procesos afectados y, por tanto, con la solución que los apoyará o
  transformará.
\item
  \textbf{Interesado no operativo:} interesado que decide, financia,
  controla o fiscaliza sin operar la solución.
\end{itemize}

Se conservan las siete agrupaciones identificadas en el Subdocumento 2.
En esta metodología, el grupo Soporte y control operacional amplía el
análisis de soporte tecnológico y control operacional con la función de
coordinación de personas, sin presumir un organigrama confirmado.

\begin{itemize}
\item
  \textbf{Dirección y control:} Directorio, Gerencia General y
  Contraloría. Resuelven las decisiones mayores, supervisan los hitos y
  verifican los controles de su ámbito.
\item
  \textbf{Propiedad:} grupo familiar controlador y fondos de inversión.
  Condicionan las decisiones estratégicas a través del Directorio.
\item
  \textbf{Negocios y operación:} Gerencia Comercial, Gerencia de Canales
  Digitales, Gerencia de Logística y Negocio Financiero. Aportan
  prioridades y reglas de negocio y validan los resultados dentro de sus
  atribuciones.
\item
  \textbf{Soporte y control operacional:} Tecnologías de la Información
  (TI), Prevención de Pérdidas y Gestión de Personas. Habilitan la
  integración y la operación, controlan el inventario y sostienen la
  dotación que debe adoptar la solución.
\item
  \textbf{Operación de tienda:} jefaturas de tienda y departamento,
  vendedores y cajeros, y personal de reposición y bodega. Usarán la
  solución todos los días.
\item
  \textbf{Clientes y terceros operacionales:} clientes, titulares de
  tarjeta, vendedores de marketplace, proveedores tecnológicos,
  proveedores logísticos, repositores externos y administradores de
  centros comerciales. Participan en los flujos que les corresponden.
\item
  \textbf{Organismos externos:} Comisión para el Mercado Financiero
  (CMF), Servicio Nacional del Consumidor (SERNAC), Servicio de
  Impuestos Internos (SII) y Agencia de Protección de Datos Personales.
  Fiscalizan y reciben evidencia a través de Ancoa.
\end{itemize}

Los cuatro procesos no se ejecutarán una sola vez. Only Simple Solutions
los repetirá al cierre de cada intervalo de planificación, al pasar de la
Etapa 1 a la Etapa 2 y antes de cada marcha blanca y paso a producción,
porque en cada uno de esos momentos cambian los interesados que deben
participar y el nivel de compromiso que el proyecto necesita de ellos.

\paragraph{6.1.4.1. Identificación y clasificación de interesados}\label{6-1-4-1-identificacion-y-clasificacion-de-interesados}

Identificar a los interesados es el proceso de reconocer de forma periódica a quienes pueden afectar o verse afectados por el proyecto, y de documentar sus intereses, su participación, sus interdependencias, su influencia y su impacto potencial en el éxito (PMI, 2017).

\textbf{Criterios de identificación y clasificación}

Una persona, área u organización se registrará como interesado cuando
cumpla al menos uno de los siguientes cuatro criterios.

\begin{enumerate}
\def\labelenumi{\arabic{enumi}.}
\item
  Operará la solución en su trabajo diario.
\item
  Decide, financia o aprueba el proyecto o alguno de sus entregables.
\item
  Su forma de trabajo, sus ingresos o sus metas cambiarán por efecto de
  la solución.
\item
  Fiscaliza una materia que la solución toca, como existencia, precio,
  crédito, datos personales o documentos tributarios.
\end{enumerate}

Cada interesado registrado se clasifica, además, según cuatro criterios.

\begin{itemize}
\item
  \textbf{Relación con Ancoa:} interno, si está bajo la autoridad
  organizacional de la compañía, o externo, si participa por una
  relación comercial, contractual o legal.
\item
  \textbf{Forma de participación:} actor, si opera la solución, o
  interesado no operativo, si decide, controla o fiscaliza sin operarla.
\item
  \textbf{Contexto:} retail, emisor, frontera retail--emisor (participa
  en ambos negocios o controla el cruce entre ellos) o plataforma
  compartida (sostiene la infraestructura que usan ambos); los dueños de
  la compañía se clasifican en el contexto propiedad. Este criterio
  responde a la operación de Ancoa, porque Ancoa es a la vez una tienda
  y un emisor de crédito fiscalizado; ningún cruce de datos entre ambos
  ámbitos se habilitará sin una finalidad definida, una base de licitud
  y los controles que correspondan (Ley N.º 21.719, 2024).
\item
  \textbf{Poder e interés:} capacidad de condicionar las decisiones del
  proyecto y grado en que la solución cambia su trabajo o sus
  obligaciones. Se analiza más adelante en esta sección.
\end{itemize}

\textbf{Herramientas y técnicas}

Una vez adjudicado el contrato, Only Simple Solutions completará la
identificación con cinco técnicas de recopilación de datos recomendadas
por el PMBOK (PMI, 2017), adaptadas a la operación de Ancoa.

\begin{itemize}
\item
  \textbf{Entrevistas:} conversaciones individuales con cada gerencia,
  con Contraloría, con la contraparte técnica de TI y con Gestión de
  Personas, para precisar expectativas, prioridades, criterios de éxito
  y preocupaciones. Retomarán a las personas consideradas en el
  levantamiento inicial de Ancoa.
\item
  \textbf{Encuestas o cuestionarios:} consultas breves al personal de
  tienda, adaptadas a sus turnos y respondidas por medios accesibles
  durante sus turnos, considerando que 3.820 vendedores y cajeros
  comparten 640 terminales.
\item
  \textbf{Talleres colaborativos:} sesiones con jefaturas de tienda,
  Logística, TI y el Negocio Financiero para validar los procesos
  críticos, como la venta con crédito, la reserva omnicanal, el cambio
  de precio y el conteo cíclico.
\item
  \textbf{Análisis de documentación organizacional:} revisión de
  organigramas, procesos, contratos con terceros e informes de auditoría
  y fiscalización, para identificar responsables y contrastar lo formal
  con lo que ocurre en la práctica.
\item
  \textbf{Observación directa:} visitas a una de las tres tiendas del
  piloto de punto de venta y al centro de distribución principal, para
  detectar tiempos muertos, errores y actores que no aparecen en las
  entrevistas.
\end{itemize}

El jefe de proyecto y el líder de implantación y gestión del cambio de
Only Simple Solutions aplicarán, además, juicio de expertos para revisar
los resultados con la contraparte técnica. Las entrevistas y talleres se
programarán fuera de las ventanas de congelamiento comercial y de los
períodos de mayor carga de las tiendas, en coordinación con Ancoa.

\textbf{Ajustes al ecosistema de interesados}

Al aplicar los criterios anteriores sobre el análisis de actores, se
identificaron cuatro ajustes al ecosistema.

\begin{itemize}
\item
  \textbf{Se incorpora una función de coordinación de personas.} Ancoa
  designará el interlocutor que coordine dotación, inducción y
  capacitación con las gerencias. Su incorporación al modelo de
  interesados no presume un cargo u organigrama ya confirmado ni
  autoridad unilateral para modificar comisiones. La carga formativa
  anual se dimensionará considerando 3.820 vendedores y cajeros, una
  rotación anual de 62 \% y cerca de 1.900 incorporaciones de temporada.
  La planificación de capacitación controlará los dobles conteos de
  personal temporal o reingresos (Bases Técnicas del Caso 09, 2026,
  núm. 2.4).
\item
  \textbf{Se mantienen diferenciados los proveedores tecnológicos y
  logísticos.} Los seis proveedores de las plataformas actuales deberán
  colaborar en el levantamiento de las interfaces; no se presume que uno
  de ellos conozca el mapa completo. Los proveedores logísticos y de
  última milla participan en los contratos y eventos del pedido. La
  estrategia de gestión distinguirá ambas dependencias.
\item
  \textbf{Se considera la Agencia de Protección de Datos Personales.} El
  modelo incorpora la autoridad creada por la Ley N.º 21.719 para las
  obligaciones de protección y tratamiento de información que sean
  aplicables durante el contrato. Su participación será fiscalizadora y
  se canalizará por los representantes autorizados de Ancoa; no se le
  asigna aprobación de decisiones internas del proyecto.
\item
  \textbf{Se precisan los registros organizacionales.} El Directorio y
  la Gerencia General se registran por separado por sus distintas
  atribuciones. Marketing se coordinará con Canales Digitales y
  Comercial según la materia, sin presumir dependencia jerárquica de una
  de ellas. Ancoa confirmará los representantes y sus facultades durante
  el levantamiento.
\end{itemize}

Con estos ajustes, la propuesta identifica un universo analítico de 26
interesados de Ancoa. Este modelo no equivale a una nómina de personas
designadas ni incluye a todo el equipo de Only Simple Solutions. Los
roles de gobierno definidos en la sección 6.1.2 corresponden a personas
de estos interesados. El patrocinador pertenece a la Gerencia General,
la contraparte técnica a TI y la contraparte de cumplimiento a
Contraloría y al Negocio Financiero.

\textbf{Ecosistema de interesados}

La Figura 6.1 ubica a los 26 interesados en sus siete grupos. El color
del borde indica el contexto, el tipo de borde distingue a los actores
de los interesados no operativos y el código de cada uno será su
identificador en el registro de interesados.

\begin{center}
\begin{minipage}{\linewidth}
\centering
\ossFigCaption{Ecosistema de actores e interesados de Multitiendas Ancoa.}\label{fig:6-1}

\includegraphics[width=\linewidth,height=0.66\textheight,keepaspectratio]{figuras/recepcion-20261025/112122353c05.png}

\raggedright\textbf{Fuente:} Elaboración propia a partir del Subdocumento 2, sección 2.4.
\end{minipage}
\end{center}

La figura distingue 13 interesados internos y 13 externos bajo la
convención de considerar Propiedad externa a la estructura operativa. La
adopción se concentra en quienes ejecutan los procesos de tienda y en
los terceros que intervienen en ellos. La frontera Retail--Emisor
involucra a dirección, Contraloría y las áreas responsables de ambos
negocios; cada cruce seguirá la ruta única de aprobación descrita en la sección 6.1.2. Las comunicaciones con fiscalizadores se realizarán por
Ancoa. Propiedad se relacionará con el proyecto por medio del
Directorio, sin crear una vía de aprobación paralela.

\textbf{Registro de interesados}

La información obtenida se documentará en el registro de interesados,
conservado en el PMIS definido en 6.1.5.1. La Tabla 6.2 sintetiza cinco
registros del universo de interesados para mostrar su función de
gestión. El registro completo identificará a I-01 a I-26; los nombres,
contactos, atribuciones y expectativas de los representantes se
registrarán durante el levantamiento del mes 1.

\begin{longtable}[]{@{}
  >{\raggedright\arraybackslash}p{(\columnwidth - 8\tabcolsep) * \real{0.2000}}
  >{\raggedright\arraybackslash}p{(\columnwidth - 8\tabcolsep) * \real{0.2000}}
  >{\raggedright\arraybackslash}p{(\columnwidth - 8\tabcolsep) * \real{0.2000}}
  >{\raggedright\arraybackslash}p{(\columnwidth - 8\tabcolsep) * \real{0.2000}}
  >{\raggedright\arraybackslash}p{(\columnwidth - 8\tabcolsep) * \real{0.2000}}@{}}
\caption{Síntesis del registro de interesados.}\label{tab:6-2}\\
\toprule\noalign{}
\begin{minipage}[b]{\linewidth}\raggedright
\textbf{ID}
\end{minipage} & \begin{minipage}[b]{\linewidth}\raggedright
\textbf{Interesado}
\end{minipage} & \begin{minipage}[b]{\linewidth}\raggedright
\textbf{Contexto}
\end{minipage} & \begin{minipage}[b]{\linewidth}\raggedright
\textbf{Función de gestión}
\end{minipage} & \begin{minipage}[b]{\linewidth}\raggedright
\textbf{Canal de validación}
\end{minipage} \\
\midrule\noalign{}
\endfirsthead
\toprule\noalign{}
\begin{minipage}[b]{\linewidth}\raggedright
\textbf{ID}
\end{minipage} & \begin{minipage}[b]{\linewidth}\raggedright
\textbf{Interesado}
\end{minipage} & \begin{minipage}[b]{\linewidth}\raggedright
\textbf{Contexto}
\end{minipage} & \begin{minipage}[b]{\linewidth}\raggedright
\textbf{Función de gestión}
\end{minipage} & \begin{minipage}[b]{\linewidth}\raggedright
\textbf{Canal de validación}
\end{minipage} \\
\midrule\noalign{}
\endhead
\bottomrule\noalign{}
\multicolumn{5}{@{}l}{\rule{0pt}{3.2ex}\textbf{Fuente:} Elaboración propia a partir de PMI (2017).} \\
\endlastfoot
I-03 & Contraloría & Frontera & Control de cumplimiento & Contraparte de
cumplimiento \\
I-12 & Coordinación de personas & Plataforma compartida & Formación e
inducción & Representante designado por Ancoa \\
I-13 & Jefaturas de tienda & Retail & Adopción y prueba en sala &
Jefaturas y usuarios clave \\
I-19 & Proveedores tecnológicos & Plataforma compartida & Dependencias e
interfaces & Gestiones contractuales de Ancoa \\
I-26 & Agencia de Protección de Datos Personales & Organismo externo &
Fiscalización aplicable & Representantes autorizados de Ancoa \\
\end{longtable}

La síntesis muestra responsabilidades y canales. El registro completo se define una sola vez, con tres bloques del PMBOK
(PMI, 2017), y esa misma definición rige la planilla del Anexo A.1 y la
propuesta inicial del Formulario T-9.

\begin{itemize}
\item
  \textbf{Identificación:} código (I-01 a I-26), interesado, grupo de
  interés y rol en el proyecto; representante titular y suplente, cargo y
  área, correo y teléfono; atribuciones confirmadas por Ancoa y canal de
  validación.
\item
  \textbf{Evaluación:} requerimientos y expectativas, influencia, interés,
  dirección de la influencia y actitud observada.
\item
  \textbf{Clasificación:} relación con Ancoa, forma de participación,
  contexto y cuadrante de poder e interés, según los cuatro criterios
  descritos al inicio de esta sección.
\end{itemize}

Los niveles actual y deseado de involucramiento no forman parte del
registro; se gestionan en la matriz de evaluación del involucramiento
(sección 6.1.4.2 y Anexo A.2). El Formulario T-9 contiene la propuesta
inicial de Only Simple Solutions, preparada con la información disponible
antes de la adjudicación: identificación y clasificación de los 26
interesados y niveles deseados propuestos. Las planillas del Anexo A son
el formato en blanco en que esa propuesta se validará y completará con
los representantes durante el levantamiento del mes 1, y en que se
mantendrán los registros vigentes durante todo el contrato. Las
evaluaciones sensibles de actitud tendrán acceso restringido al jefe de
proyecto, al líder de implantación y al patrocinador de Ancoa, conforme a
los permisos del plan.

\textbf{Análisis de poder e interés}

Con el registro completo se elaborará la matriz de poder e interés. El
poder es la capacidad del interesado para aprobar, detener o condicionar
decisiones del proyecto; el interés es el grado en que la solución
cambia su trabajo, sus resultados o las obligaciones que fiscaliza. La
Figura 6.2 muestra la ubicación de los 26 interesados, relacionada con
el análisis del Subdocumento 2.

\begin{center}
\begin{minipage}{\linewidth}
\centering
\ossFigCaption{Matriz de poder e interés de los interesados.}\label{fig:6-2}

\includegraphics[width=\linewidth,height=0.66\textheight,keepaspectratio]{figuras/recepcion-20261025/0e54198d8fb4.png}

\raggedright\textbf{Fuente:} Elaboración propia a partir de la Figura 6.1 y de PMI (2017).
\end{minipage}
\end{center}

Doce interesados quedan en el cuadrante gestionar de cerca,
correspondientes a la dirección, las cuatro gerencias de negocio, los
tres integrantes de Soporte y control operacional, la CMF y el SERNAC.
Las dos autoridades se ubican allí porque Ancoa mantiene con ellas un
plan de remediación y un procedimiento abierto; aun así, Only Simple
Solutions no se comunicará directamente con ellas, sino que entregará la
evidencia a la gerencia de Ancoa que tiene esa relación (sección
6.1.5.1). La función de coordinación de personas se ubica en este
cuadrante por su influencia sobre la disponibilidad para capacitación.
Su interlocutor coordinará dotación, inducción y capacitación, sin
atribución para decidir unilateralmente reglas de comisión.

Dos ubicaciones requieren una precisión. Las jefaturas de tienda figuran
en mantener informados por su cargo, pero su influencia práctica es
mayor. Sin su apoyo, la solución no se usa en sala, por lo que recibirán
además el rol de participación descrito en la sección 6.1.4.2. Los
proveedores tecnológicos tienen poco poder formal, pero conocen las
catorce interfaces actuales; su colaboración se asegurará también por
vía contractual, a través de Ancoa y de la gestión de adquisiciones.

Only Simple Solutions registrará también la dirección de la influencia
de cada interesado (PMI, 2017), ascendente hacia el Directorio y la
Gerencia General, descendente hacia las jefaturas y el personal de
tienda, lateral entre gerencias y con la filial, y hacia afuera con
clientes, terceros y autoridades. Esta dirección define quién comunica
cada mensaje. Un cambio en tienda lo anunciará la gerencia responsable a
través de cada jefatura, no Only Simple Solutions.

\textbf{Salidas del proceso}

El proceso producirá el registro de interesados y actualizaciones del
registro de supuestos y del registro de riesgos. Si aparece un
interesado nuevo, como un proveedor que se incorpora o una tienda que
cambia de administrador, se actualizarán el registro y la matriz; si el
cambio afecta una línea base o un entregable, se tramitará mediante el
control integrado de cambios (sección 6.1.3.6).

\paragraph{6.1.4.2. Planificar el involucramiento de los interesados}\label{6-1-4-2-planificar-el-involucramiento-de-los-interesados}

Planificar el involucramiento es el proceso de desarrollar enfoques para involucrar a los interesados, a partir de sus necesidades, expectativas, intereses y posible impacto en el proyecto (PMI, 2017). Su resultado será la estrategia de involucramiento de los interesados, que formará parte del plan para la dirección del proyecto y se aprobará en el hito H1. El proceso usará como entradas el registro de interesados, el plan de gestión de las comunicaciones (sección 6.1.5), el registro de riesgos, los supuestos del Subdocumento 2, el calendario de hitos, las olas de implantación en tiendas y las cinco ventanas de congelamiento comercial.

\textbf{Matriz de evaluación del involucramiento}

La herramienta central del proceso es la matriz de evaluación del
involucramiento, que compara el nivel actual de compromiso de cada
interesado con el nivel que el proyecto necesita para tener éxito (PMI,
2017). La brecha entre ambos determina las acciones de involucramiento,
y su cierre se mide en el monitoreo. La matriz usa cinco niveles.

\begin{itemize}
\item
  \textbf{Desconocedor:} no conoce el proyecto ni los impactos que puede
  generar.
\item
  \textbf{Reticente:} conoce el proyecto y sus implicancias, pero se
  resiste al cambio.
\item
  \textbf{Neutral:} conoce el proyecto, pero no manifiesta apoyo ni
  oposición.
\item
  \textbf{De apoyo:} comprende el alcance y los impactos del proyecto y
  colabora con sus actividades y resultados.
\item
  \textbf{Líder:} está plenamente informado y comprometido, y promueve
  activamente el éxito del proyecto dentro de la organización.
\end{itemize}

La Tabla 6.3 presenta cinco ejemplos de acciones y niveles deseados. El
nivel actual se determinará mediante entrevistas, encuestas y talleres
del levantamiento. Los niveles deseados se acordarán con Ancoa al
aprobar la estrategia de involucramiento, con sus acciones, responsables
y momentos de comprobación.

\begin{longtable}[]{@{}
  >{\raggedright\arraybackslash}p{(\columnwidth - 8\tabcolsep) * \real{0.2000}}
  >{\raggedright\arraybackslash}p{(\columnwidth - 8\tabcolsep) * \real{0.2000}}
  >{\raggedright\arraybackslash}p{(\columnwidth - 8\tabcolsep) * \real{0.2000}}
  >{\raggedright\arraybackslash}p{(\columnwidth - 8\tabcolsep) * \real{0.2000}}
  >{\raggedright\arraybackslash}p{(\columnwidth - 8\tabcolsep) * \real{0.2000}}@{}}
\caption{Matriz de evaluación del involucramiento.}\label{tab:6-3}\\
\toprule\noalign{}
\begin{minipage}[b]{\linewidth}\raggedright
\textbf{ID}
\end{minipage} & \begin{minipage}[b]{\linewidth}\raggedright
\textbf{Interesado}
\end{minipage} & \begin{minipage}[b]{\linewidth}\raggedright
\textbf{Nivel actual}
\end{minipage} & \begin{minipage}[b]{\linewidth}\raggedright
\textbf{Nivel deseado}
\end{minipage} & \begin{minipage}[b]{\linewidth}\raggedright
\textbf{Acción de involucramiento}
\end{minipage} \\
\midrule\noalign{}
\endfirsthead
\toprule\noalign{}
\begin{minipage}[b]{\linewidth}\raggedright
\textbf{ID}
\end{minipage} & \begin{minipage}[b]{\linewidth}\raggedright
\textbf{Interesado}
\end{minipage} & \begin{minipage}[b]{\linewidth}\raggedright
\textbf{Nivel actual}
\end{minipage} & \begin{minipage}[b]{\linewidth}\raggedright
\textbf{Nivel deseado}
\end{minipage} & \begin{minipage}[b]{\linewidth}\raggedright
\textbf{Acción de involucramiento}
\end{minipage} \\
\midrule\noalign{}
\endhead
\bottomrule\noalign{}
\multicolumn{5}{@{}l}{\rule{0pt}{3.2ex}\textbf{Fuente:} Elaboración propia a partir de PMI (2017).} \\
\endlastfoot
I-02 & Gerencia General & Se determina en levantamiento & Líder &
Resolver prioridades y escalamiento \\
I-06 & Comercial & Se determina en levantamiento & De apoyo & Acordar
reglas de precio y comisiones \\
I-12 & Coordinación de personas & Se determina en levantamiento & De
apoyo & Programar formación por perfil y turno \\
I-13 & Jefaturas de tienda & Se determina en levantamiento & De apoyo &
Participar en pilotos y adopción \\
I-19 & Proveedores tecnológicos & Se determina en levantamiento & De
apoyo & Verificar interfaces y transición \\
\end{longtable}

Las entrevistas del levantamiento inicial permiten anticipar dónde
estarán las brechas más críticas. La jefatura de tienda entrevistada
señala que el despacho desde tienda perjudica los resultados de su
local, y el vendedor advierte que dejaría de ofrecer la tarjeta si el
proceso se alarga; ambos grupos podrían partir como reticentes. La
Gerencia Comercial y el Negocio Financiero apoyan el proyecto, pero
objetan reglas que les afectan, como el control de los cambios de precio
y la secuencia que deja el crédito para después. Estas señales
orientarán las primeras entrevistas y talleres, y la matriz se
actualizará al cierre de cada intervalo de planificación.

\textbf{Fase de análisis de los interesados}

Para diseñar la estrategia, Only Simple Solutions aplicará tres fases
complementarias, de análisis, diseño y validación. En la fase de
análisis se usarán el registro de interesados, la matriz de poder e
interés y la matriz de evaluación del involucramiento, y se evaluarán
tres aspectos de cada interesado.

\begin{itemize}
\item
  \textbf{Nivel de influencia:} capacidad para aprobar, detener o
  condicionar el trabajo, según su posición en la matriz de poder e
  interés y sus atribuciones formales.
\item
  \textbf{Tipo de interés:} estratégico, cuando se relaciona con
  resultados de negocio o decisiones corporativas; operativo, cuando el
  proyecto modifica el trabajo cotidiano; técnico, cuando se relaciona
  con integración, datos u operación de la plataforma; o regulatorio,
  cuando se relaciona con una obligación fiscalizada.
\item
  \textbf{Brecha de involucramiento:} diferencia entre el nivel actual y
  el deseado registrados en la matriz de la Tabla 6.3, que indica
  quiénes requieren acciones específicas.
\end{itemize}

Este análisis define el enfoque de gestión de cada cuadrante, que la
Tabla 6.4 resume. La ubicación de cada interesado se revisará durante el
proyecto y no se considerará permanente.

\begin{longtable}[]{@{}
  >{\raggedright\arraybackslash}p{(\columnwidth - 6\tabcolsep) * \real{0.2500}}
  >{\raggedright\arraybackslash}p{(\columnwidth - 6\tabcolsep) * \real{0.2500}}
  >{\raggedright\arraybackslash}p{(\columnwidth - 6\tabcolsep) * \real{0.2500}}
  >{\raggedright\arraybackslash}p{(\columnwidth - 6\tabcolsep) * \real{0.2500}}@{}}
\caption{Enfoque de gestión por cuadrante de poder e interés.}\label{tab:6-4}\\
\toprule\noalign{}
\begin{minipage}[b]{\linewidth}\raggedright
\textbf{Cuadrante}
\end{minipage} & \begin{minipage}[b]{\linewidth}\raggedright
\textbf{Interesados}
\end{minipage} & \begin{minipage}[b]{\linewidth}\raggedright
\textbf{Enfoque de gestión}
\end{minipage} & \begin{minipage}[b]{\linewidth}\raggedright
\textbf{Frecuencia de contacto}
\end{minipage} \\
\midrule\noalign{}
\endfirsthead
\toprule\noalign{}
\begin{minipage}[b]{\linewidth}\raggedright
\textbf{Cuadrante}
\end{minipage} & \begin{minipage}[b]{\linewidth}\raggedright
\textbf{Interesados}
\end{minipage} & \begin{minipage}[b]{\linewidth}\raggedright
\textbf{Enfoque de gestión}
\end{minipage} & \begin{minipage}[b]{\linewidth}\raggedright
\textbf{Frecuencia de contacto}
\end{minipage} \\
\midrule\noalign{}
\endhead
\bottomrule\noalign{}
\multicolumn{4}{@{}l}{\rule{0pt}{3.2ex}\textbf{Fuente:} Elaboración propia a partir de PMI (2017).} \\
\endlastfoot
Gestionar de cerca & I-01 a I-03, I-06 a I-12, I-23 e I-24 & Decisión y
validación interna; autoridades mediante evidencia por Ancoa & Semanal a
mensual según instancia; autoridades a requerimiento \\
Mantener satisfechos & I-04, I-05, I-25 e I-26 & Información resumida
por hito, canalizada por el Directorio. Evidencia de cumplimiento a
través de Ancoa & Por hito o a requerimiento \\
Mantener informados & I-13 a I-20 & Información oportuna, capacitación,
pilotos y canales de retroalimentación. Terceros coordinados a través de
Ancoa & Mensual; semanal en pilotos y marcha blanca \\
Monitorear & I-21 e I-22 & Accesos, obras y enlaces en centros
comerciales & Por ola de
implantación \\
\end{longtable}

El enfoque concentra el esfuerzo donde está el riesgo. Los diez
interesados internos de gestión cercana participan en decisiones y
validaciones dentro de sus atribuciones. La CMF y el SERNAC, también
incluidos en ese cuadrante, ejercen influencia fiscalizadora y reciben
antecedentes por medio de Ancoa; no integran el gobierno del proyecto.
Los ocho interesados de mantener informados participan en la operación
y en los flujos con terceros. Los otros dos cuadrantes requieren poco
tiempo, pero un incumplimiento ante una autoridad o una obra no
coordinada en un centro comercial puede detener una puesta en marcha.

\textbf{Fase de diseño de estrategias de involucramiento}

En la segunda fase se diseñarán estrategias diferenciadas según el
perfil de cada interesado, organizadas en cuatro dimensiones.

\begin{itemize}
\item
  \textbf{Comunicación:} cada grupo tendrá canales formales definidos en
  la matriz de comunicaciones (sección 6.1.5.2). El Directorio, la
  Gerencia General y las gerencias recibirán la información en el Comité
  Ejecutivo mensual; los dueños de producto revisarán los incrementos en
  las revisiones de sprint; y TI, Prevención de Pérdidas y Contraloría
  participarán en las validaciones de su ámbito. Las jefaturas recibirán
  un boletín mensual, semanal durante la marcha blanca, y el personal de
  tienda conocerá el propósito del cambio mediante mensajes breves por
  canales disponibles y demostraciones guiadas. Con los externos, la
  comunicación se hará siempre a través de Ancoa.
\item
  \textbf{Capacitación y formación:} la formación se diseñará para una
  dotación con 62 \% de rotación anual y cerca de 1.900 incorporaciones
  de temporada. Tendrá módulos breves, de no más de 30 minutos, práctica
  en un ambiente de entrenamiento, ayuda en pantalla y un curso de
  inducción permanente que Gestión de Personas incorporará a la
  inducción de cada nuevo vendedor o cajero. Los dueños de producto y
  las jefaturas se prepararán para revisar incrementos y criterios de
  aceptación, el personal de TI recibirá formación técnica continua para
  operar la solución, y los repositores externos se capacitarán a través
  de su empleador, con material entregado por Only Simple Solutions. El
  plan cubrirá personas usuarias finales y avanzadas, administradoras,
  equipo técnico y soporte de niveles 1 y 2; combinará formación
  presencial por sitio, sesiones sincrónicas, autoformación y
  acompañamiento en puesto. El material será editable, en español y de
  propiedad de Ancoa. Durante Operación se capacitará al personal nuevo
  y se realizarán al menos dos jornadas anuales de actualización, sin
  costo adicional; la mentoría técnica continuará al menos seis meses
  después de la producción de Etapa 2 (Bases Técnicas Transversales,
  2026, RT-22.01 a RT-22.08).
\item
  \textbf{Participación y co-creación:} Ancoa designará un product owner
  para cada dominio, con precio y catálogo en la Gerencia Comercial,
  inventario y abastecimiento en Logística, pedidos en Canales Digitales
  y crédito en el Negocio Financiero. Cada área nombrará usuarios clave
  que participarán en las revisiones de sprint y en las pruebas de
  aceptación. Prevención de Pérdidas participará en las pruebas de
  inventario, y Contraloría y el Negocio Financiero en la validación de
  la separación de datos. Las jefaturas, vendedores y cajeros de las
  tiendas del piloto probarán la solución antes de cada ola, y se
  realizarán talleres con vendedores de marketplace, proveedores y
  transportistas, a través de Ancoa, para validar sus integraciones.
  Only Simple Solutions se compromete a que cada una de las 22 tiendas cuente con un agente de cambio activo antes del hito H6. La designación corresponde a la jefatura de cada tienda y es una dependencia de Ancoa; Only Simple Solutions propondrá el perfil, capacitará a los agentes y escalará al Comité de Proyecto toda designación pendiente.
\item
  \textbf{Reconocimiento y refuerzo positivo:} Only Simple Solutions y
  Ancoa destacarán a las tiendas, agentes de cambio y usuarios piloto
  que aporten mejoras o alcancen antes sus metas de adopción, y los
  aportes de los usuarios clave quedarán registrados en las notas de
  cada versión. Para el personal de venta, Only Simple Solutions
  entregará los datos que permitan a la Gerencia Comercial y a Gestión
  de Personas evaluar ajustes a la regla de comisión, de modo que los
  nuevos pasos de la venta no reduzcan el ingreso del vendedor; la
  decisión corresponde a Ancoa.
\end{itemize}

\textbf{Fase de validación de la estrategia}

En la tercera fase, las estrategias se presentarán a los representantes
de Ancoa de cada grupo de interés para su validación y ajuste. La
estrategia de involucramiento se aprobará junto con el plan para la
dirección del proyecto en el hito H1 y contendrá la matriz de evaluación
validada, las estrategias por dimensión, el calendario por ola y las
dedicaciones comprometidas por Ancoa. En la misma instancia se confirmarán las metas de los indicadores de la Tabla 6.5, que Only Simple Solutions propone desde ya; cualquier ajuste quedará fundamentado en el acta. Desde entonces, todo
cambio en las dedicaciones o en los niveles deseados se revisará en el
Comité de Proyecto y, si afecta una línea base, se tramitará como
solicitud de cambio.

\paragraph{6.1.4.3. Gestionar la participación de los interesados}\label{6-1-4-3-gestionar-la-participacion-de-los-interesados}

Gestionar la participación es el proceso de comunicarse y trabajar con los interesados para satisfacer necesidades y expectativas, abordar incidentes y sostener su involucramiento (PMI, 2017). Only Simple Solutions iniciará la interacción en el mes 1 para levantar y coordinar necesidades; aplicará la estrategia formal desde su aprobación en H1. La gestión integra las instancias del proyecto y la inspección de incrementos, y se organizará en las cuatro líneas de trabajo siguientes.

\textbf{1. Ejecución de las estrategias de comunicación}

La participación ocurrirá en las instancias de gobierno de la sección
6.1.2 y en las ceremonias ágiles de la sección 6.2. Cada instancia tiene
un propósito distinto para los interesados.

\begin{itemize}
\item
  \textbf{Comité Ejecutivo, mensual:} el patrocinador, la gerencia de
  Only Simple Solutions y el administrador del contrato resuelven
  decisiones mayores y conflictos escalados. Las gerencias de negocio
  aportan antecedentes cuando corresponda.
\item
  \textbf{Comité de Proyecto, quincenal:} la contraparte técnica y el
  jefe de proyecto revisan avance, cambios y estado del involucramiento,
  con aportes de los dueños de producto.
\item
  \textbf{Revisión de sprint, cada dos semanas:} los usuarios clave
  prueban lo construido y entregan retroalimentación.
\item
  \textbf{Comité de Arquitectura, mensual:} TI y el encargado de
  seguridad revisan las decisiones técnicas que afectan la operación
  futura.
\item
  \textbf{Comité de Operación, mensual desde el mes 13:} el líder de
  operación, la mesa de servicio y la contraparte técnica revisan
  preparación, marcha blanca y, desde el mes 16, el alcance en
  producción, con las áreas usuarias pertinentes.
\end{itemize}

En las tiendas, la participación se canalizará a través de la jefatura y
del agente de cambio, que recogerán las dudas del personal y las
llevarán a la reunión semanal de seguimiento de su ola. Los acuerdos,
solicitudes y decisiones de todas las instancias quedarán registrados en
Jira, para asegurar su trazabilidad.

\textbf{2. Participación y co-creación}

La retroalimentación de los interesados se transformará en mejoras
concretas mediante talleres funcionales con usuarios clave, pruebas de
aceptación con jefaturas y usuarios de las tiendas del piloto, y
talleres de integración con terceros a través de Ancoa. Las
observaciones se registrarán en plantillas de control y el product owner
que corresponda las incorporará al backlog, dentro del alcance aprobado.

\textbf{3. Gestión de expectativas y resolución de conflictos}

Durante la ejecución surgirán diferencias entre lo planificado y lo que
esperan los interesados. Only Simple Solutions aplicará un protocolo de
cuatro pasos.

\begin{enumerate}
\def\labelenumi{\arabic{enumi}.}
\item
  Registro de la observación o del conflicto en Jira, con su responsable
  de Ancoa y la fecha en que se necesita la decisión.
\item
  Análisis de su impacto en alcance, plazo, costo y riesgo, coordinado
  por el jefe de proyecto.
\item
  Reunión de conciliación con los interesados afectados, facilitada por
  el scrum master o por el líder de implantación y gestión del cambio.
\item
  Comunicación del acuerdo y, si modifica una línea base, tramitación
  mediante el control integrado de cambios (sección 6.1.3.6).
\end{enumerate}

El análisis de actores del Subdocumento 2 anticipa siete tensiones entre
interesados. Only Simple Solutions no las resolverá por su cuenta. Las
hará visibles, preparará alternativas con datos y llevará la decisión a
la instancia que corresponda.

\begin{enumerate}
\def\labelenumi{\arabic{enumi}.}
\item
  \textbf{Rapidez del crédito frente a evidencia del consentimiento
  (I-09 e I-03).} La evaluación y la entrega de la información
  precontractual se diseñarán como un solo flujo, que ambos dueños
  validarán antes de su paso a producción.
\item
  \textbf{Comisión del vendedor frente a los pasos de cumplimiento (I-14
  e I-03).} El flujo de venta con crédito se probará en las tiendas del
  piloto midiendo su duración, y los resultados se entregarán a la
  Gerencia Comercial y a Gestión de Personas para revisar la regla de
  comisión.
\item
  \textbf{Existencia publicada frente a existencia confiable (I-07 e
  I-08).} La regla de publicación la definirán ambas gerencias con los
  datos de exactitud de inventario de cada ola; si no hay acuerdo,
  decidirá el Comité de Proyecto.
\item
  \textbf{Velocidad del cambio de precio frente al precio exhibido (I-06
  e I-13).} Los cambios de precio se programarán en oleadas coordinadas
  con el reetiquetado de las tiendas, porque el SERNAC fiscaliza que el
  precio cobrado sea igual al exhibido.
\item
  \textbf{Despacho desde tienda frente a la venta en sala (I-07, I-13 e
  I-14).} La regla de asignación de unidades y el reconocimiento de la
  venta a la tienda que despacha los decidirá la Gerencia Comercial, con
  datos de ambos canales.
\item
  \textbf{Uso comercial de datos financieros (I-06, I-03 e I-09).}
  Ningún cruce de datos entre retail y emisor se habilitará sin finalidad definida y base de licitud (Ley N.º 21.719, 2024), ni fuera de la ruta única de la sección 6.1.2: validación jurídica de Ancoa y aprobación conjunta de Contraloría y del Negocio Financiero.
\item
  \textbf{Prioridad del inventario frente al plazo de 2029 (I-02 e
  I-09).} La preparación de la migración de la cartera se planificará
  desde la Etapa 1 para no comprometer el fin del soporte de la
  plataforma de crédito; la secuencia la resolverá el Comité Ejecutivo.
\end{enumerate}

Para estas conversaciones, el equipo aplicará las habilidades
interpersonales que indica el PMBOK, como gestión de conflictos,
negociación, conciencia cultural y conciencia política (PMI, 2017). Si
un conflicto no se resuelve, se escalará en tres niveles. Primero, el
product owner con el jefe de proyecto, en un plazo de cinco días
hábiles; luego, el Comité de Proyecto en su sesión siguiente; y
finalmente, el Comité Ejecutivo, para conflictos entre gerencias o que
afecten un hito.

\textbf{4. Integración con la gestión del cambio organizacional}

Para el grupo Operación de tienda, la participación se complementará con
el modelo ADKAR, que ordena el cambio individual en cinco elementos,
conciencia, deseo, conocimiento, habilidad y refuerzo (Prosci, s.f.). El
líder de implantación y gestión del cambio de Only Simple Solutions lo
aplicará junto con Gestión de Personas y los agentes de cambio. La
conciencia y el deseo se trabajarán desde las ocho semanas previas a
cada ola, y el conocimiento y la habilidad desde las cuatro previas; el
refuerzo acompañará la marcha blanca y el mes siguiente. La programación
por turnos y dedicaciones se justificará en el plan de implantación. El
logro de cada elemento se comprobará por comprensión, participación y
ejecución práctica, sin darlo por alcanzado únicamente porque transcurra
ese tiempo. La secuencia comenzará en las tiendas del piloto, y sus
lecciones ajustarán el plan antes de la primera ola.

\textbf{Verificación antes de cada marcha blanca y paso a producción}

Antes de cada marcha blanca y paso a producción, el jefe de proyecto y
la contraparte técnica comprobarán cinco condiciones de participación.

\begin{enumerate}
\def\labelenumi{\arabic{enumi}.}
\item
  Están representadas las áreas realmente afectadas, con representantes
  que tienen las atribuciones necesarias.
\item
  Los dueños de producto están designados y cumplen su dedicación, y no
  hay decisiones de Ancoa pendientes fuera de plazo.
\item
  La estrategia no vulnera la separación entre retail y emisor; las
  materias de crédito y datos cuentan con las validaciones de datos,
  seguridad y cumplimiento.
\item
  Las tiendas de la ola tienen su agente de cambio y su personal
  capacitado.
\item
  Los terceros afectados, como proveedores, repositores y
  administradores de centros comerciales, recibieron aviso con
  anticipación a través de Ancoa.
\end{enumerate}

Si una condición no se cumple, la situación se registrará como
impedimento o riesgo, y el jefe de proyecto acordará una alternativa con
la contraparte técnica; si afecta una decisión crítica, la elevará a la
instancia de gobierno que corresponda. Una sesión con usuarios o una
aprobación comercial no podrá sustituir las validaciones de datos,
seguridad y cumplimiento.

\textbf{Salidas del proceso}

El proceso producirá solicitudes de cambio y actualizaciones del
registro de interesados, del registro de incidentes, de la estrategia de
involucramiento y del plan de comunicaciones. Lo aprendido en cada ola,
en especial lo que funcionó o falló con las tiendas reticentes, se
registrará en el registro de lecciones aprendidas y se aplicará a la ola
siguiente.

\paragraph{6.1.4.4. Monitorear el involucramiento de los interesados}\label{6-1-4-4-monitorear-el-involucramiento-de-los-interesados}

Monitorear el involucramiento es el proceso de supervisar las relaciones con los interesados y adaptar las estrategias para involucrarlos (PMI, 2017). Será un proceso iterativo y basado en evidencia, orientado a comprobar si cada interesado se mantiene en el nivel requerido y a detectar a tiempo cambios en su influencia, disponibilidad o disposición frente al proyecto.

El jefe de proyecto actualizará la matriz de evaluación del
involucramiento al cierre de cada intervalo de planificación, con
información de comités, revisiones, talleres, pruebas y encuestas.
Además, hará una revisión extraordinaria antes de cada marcha blanca y
paso a producción, al inicio de la Etapa 2 y cuando cambie un
representante o sus responsabilidades. Los resultados se revisarán en el
Comité de Proyecto. Para cada grupo de interés se registrarán, como
mínimo, cuatro aspectos.

\begin{itemize}
\item
  Nivel de participación en las actividades previstas, como comités,
  revisiones, validaciones y capacitaciones.
\item
  Cumplimiento de los compromisos asumidos.
\item
  Retroalimentación cualitativa sobre la comunicación y la implantación.
\item
  Evolución de su nivel de involucramiento respecto del nivel deseado.
\end{itemize}

Estos antecedentes se consolidarán en los indicadores de la Tabla 6.5. Sus metas son la propuesta de Only Simple Solutions, se confirman en H1 y son umbrales internos de gestión y se
gestionarán sin reducir obligaciones contractuales ni condiciones de
aceptación.

\begin{longtable}[]{@{}
  >{\raggedright\arraybackslash}p{(\columnwidth - 8\tabcolsep) * \real{0.2000}}
  >{\raggedright\arraybackslash}p{(\columnwidth - 8\tabcolsep) * \real{0.2000}}
  >{\raggedright\arraybackslash}p{(\columnwidth - 8\tabcolsep) * \real{0.2000}}
  >{\raggedright\arraybackslash}p{(\columnwidth - 8\tabcolsep) * \real{0.2000}}
  >{\raggedright\arraybackslash}p{(\columnwidth - 8\tabcolsep) * \real{0.2000}}@{}}
\caption{Indicadores de gestión del involucramiento.}\label{tab:6-5}\\
\toprule\noalign{}
\begin{minipage}[b]{\linewidth}\raggedright
\textbf{Indicador}
\end{minipage} & \begin{minipage}[b]{\linewidth}\raggedright
\textbf{Cálculo}
\end{minipage} & \begin{minipage}[b]{\linewidth}\raggedright
\textbf{Fuente de información}
\end{minipage} & \begin{minipage}[b]{\linewidth}\raggedright
\textbf{Frecuencia}
\end{minipage} & \begin{minipage}[b]{\linewidth}\raggedright
\textbf{Meta}
\end{minipage} \\
\midrule\noalign{}
\endfirsthead
\toprule\noalign{}
\begin{minipage}[b]{\linewidth}\raggedright
\textbf{Indicador}
\end{minipage} & \begin{minipage}[b]{\linewidth}\raggedright
\textbf{Cálculo}
\end{minipage} & \begin{minipage}[b]{\linewidth}\raggedright
\textbf{Fuente de información}
\end{minipage} & \begin{minipage}[b]{\linewidth}\raggedright
\textbf{Frecuencia}
\end{minipage} & \begin{minipage}[b]{\linewidth}\raggedright
\textbf{Meta}
\end{minipage} \\
\midrule\noalign{}
\endhead
\bottomrule\noalign{}
\multicolumn{5}{@{}l}{\rule{0pt}{3.2ex}\textbf{Fuente:} Elaboración propia a partir de PMI (2017).} \\
\endlastfoot
Participación efectiva & Convocados que asistieron / convocados & Actas
y listas de asistencia & Mensual & ≥ 85 \% \\
Cumplimiento de compromisos & Cumplidos en plazo / compromisos
vencidos del período & Registro de compromisos en Jira & Mensual & ≥ 90
\% \\
Decisiones de Ancoa a tiempo & Resueltas en plazo / decisiones vencidas
del período & Registro de decisiones en Jira & Mensual & ≥ 90
\% \\
Satisfacción de los interesados & Promedio de encuesta breve, escala de
1 a 5 & Encuesta posterior a revisiones y pilotos & Por intervalo de planificación & ≥
4,0 \\
Resolución de conflictos & Días hábiles desde el registro hasta la
resolución efectiva & Registro de incidentes en Jira & Mensual & ≤ 10
días hábiles \\
Cobertura de agentes de cambio & Tiendas con agente de cambio activo /
22 tiendas & Registro de agentes de cambio & Por intervalo de planificación & 22 de 22 antes de H6 (designación por la jefatura de cada tienda) \\
Personal capacitado por ola & Personal que aprobó la capacitación /
dotación de la ola & Plataforma de capacitación & Antes de cada ola & ≥
95 \% \\
Interesados internos críticos en nivel deseado & Internos de gestión
cercana en nivel deseado / internos de ese cuadrante & Matriz de
evaluación del involucramiento & Por intervalo de planificación & 100 \% antes de cada
paso a producción \\
\end{longtable}

Los tres primeros indicadores miden la participación de quienes deciden
y advierten temprano una baja de compromiso. Los indicadores de agentes
de cambio y de personal capacitado miden la preparación de quienes
ejecutan el cambio en sala, que es donde las entrevistas anticipan las
brechas más grandes. El último cierra el ciclo. Verifica que los
interesados de gestión cercana hayan llegado al nivel que la estrategia
comprometió.

Las tasas de cumplimiento de compromisos y de decisiones a tiempo
utilizarán, en numerador y denominador, la misma cohorte de elementos
cuyo plazo venció en el período de medición. Los elementos con plazo
futuro no se contarán como atrasados ni se mezclarán con esa cohorte.

El escalamiento de conflictos se registrará por separado de su
resolución. Un asunto escalado continuará abierto hasta comprobar el
acuerdo y su ejecución. Si no existen casos o población aplicable en un
período, el responsable registrará esa situación en vez de producir un
porcentaje sin denominador. Las metas son objetivos internos; el
cumplimiento de una obligación contractual se verificará en cada evento.
Para autoridades externas se controlarán requerimientos, antecedentes y
plazos, sin medirlas como personas que deban adoptar una actitud de
apoyo al proveedor.

\textbf{Umbrales y acciones de mejora}

Only Simple Solutions activará una acción de mejora cuando un indicador
quede bajo su meta en dos mediciones seguidas, o cuando un interesado de
gestión cercana retroceda un nivel en la matriz. Según la causa, las
acciones podrán incluir las siguientes.

Una condición bloqueante de preparación, seguridad o aceptación se
atenderá desde su detección, sin esperar dos mediciones. El umbral
interno de capacitación del 95 \% no sustituye las condiciones por
perfil del plan aprobado ni la evaluación y certificación de las personas
administradoras y del equipo técnico de Ancoa.

\begin{itemize}
\item
  Reforzar la comunicación en los canales menos utilizados o cambiar el
  canal de un grupo.
\item
  Reprogramar talleres o capacitaciones con baja asistencia, en horarios
  de menor afluencia.
\item
  Convocar reuniones extraordinarias de coordinación en las áreas
  críticas.
\item
  Solicitar a la gerencia respectiva el cambio o la incorporación de
  representantes, si un product owner no cumple su dedicación.
\item
  Activar un plan de recuperación y escalar la preparación insuficiente
  de una ola. Los ajustes utilizarán las holguras aprobadas sin
  desplazar las fechas contractuales ni dispensar capacitación,
  certificación o condiciones de aceptación.
\end{itemize}

Los ajustes se aprobarán en el Comité de Proyecto y quedarán en la
versión vigente de la estrategia de involucramiento. Si un ajuste
modifica una línea base, un entregable o un compromiso contractual, se
tramitará mediante el control integrado de cambios (sección 6.1.3.6). Si
un interesado crítico interno permanece bajo el nivel requerido durante
dos intervalos de planificación seguidos, el asunto se escalará al Comité Ejecutivo.

\textbf{Salidas del proceso}

El proceso producirá información de desempeño del trabajo, que se
incluirá en el informe mensual al Comité Ejecutivo, además de
solicitudes de cambio y actualizaciones del registro de interesados, de
la estrategia de involucramiento y del plan de comunicaciones (PMI,
2017). El monitoreo alimentará también el registro de lecciones
aprendidas, con las prácticas que favorecieron la participación y las
que generaron resistencia. Al cierre de cada etapa, en los meses 16 y
21, los resultados del involucramiento se incorporarán al informe de
cierre correspondiente. Desde el mes 21, el Comité de Operación
continuará el seguimiento de los usuarios hasta el término del contrato
en el mes 56.

\subsubsection{6.1.5. Gestión de comunicaciones y seguimiento}\label{6-1-5-gestion-de-comunicaciones-y-seguimiento}

En Multitiendas Ancoa, la comunicación forma parte del control operacional del proyecto. Un precio mal informado puede transformarse en una diferencia observable por el cliente, un incidente en caja que no se escala a tiempo puede convertirse en una venta perdida y una decisión pendiente puede bloquear el trabajo de desarrollo o implantación. Por ello, la gestión de las comunicaciones asegurará que cada interesado reciba la información que necesita, en el momento adecuado, por un canal acorde con su forma de trabajo y con un registro que permita rastrear acuerdos, decisiones y compromisos.

Se aplicarán los tres procesos del área de gestión de las comunicaciones
del PMBOK, planificar la gestión de las comunicaciones, gestionar las
comunicaciones y monitorearlas. Esta gestión se apoyará en los siete
grupos de interés identificados en el Subdocumento 2, en la estrategia
de involucramiento y en las cadencias de gobierno definidas para el
proyecto.

El seguimiento distinguirá las comunicaciones periódicas de las
comunicaciones extraordinarias. Las primeras mantendrán informados a los
equipos y a las instancias de gobierno sobre el avance y los
compromisos. Las segundas permitirán notificar oportunamente las
desviaciones relevantes de los hitos y los incidentes críticos de
seguridad, mediante responsables, canales y plazos definidos, sin
depender de la siguiente reunión programada.

\paragraph{6.1.5.1. Planificar la gestión de las comunicaciones}\label{6-1-5-1-planificar-la-gestion-de-las-comunicaciones}

La planificación definirá qué información necesita cada interesado, quién la emite, por qué medio, con qué frecuencia y con qué reglas. Su resultado será el plan de gestión de las comunicaciones, que formará parte del plan para la dirección del proyecto, se elaborará en el mes 1 y se aprobará en el hito H1.

Como referencia del universo analítico se calcula el máximo teórico de
pares de comunicación mediante $n(n-1)/2$. Con 26 interesados, $26 \times 25/2$
= 325 pares potenciales. El valor no representa conversaciones
efectivas, personas nominadas ni toda la coordinación del equipo OSS;
ayuda a justificar canales y niveles definidos, en vez de una
comunicación indiscriminada entre todos. El plan se apoya en las tres
decisiones siguientes.

\begin{enumerate}
\def\labelenumi{\arabic{enumi}.}
\item
  Fuente única de datos de gestión. Jira será el registro operativo de
  tareas, responsables, plazos, incidencias y compromisos, y Microsoft
  SharePoint será el espacio colaborativo y repositorio de documentación
  vigente, accesible para los representantes autorizados de Ancoa. En
  conjunto constituirán el sistema de información para la dirección del
  proyecto (PMIS). Slack se utilizará para la coordinación cotidiana
  entre los equipos, pero los acuerdos que afecten decisiones,
  responsabilidades, plazos o entregables deberán formalizarse en Jira o
  SharePoint, según corresponda. Los informes se elaborarán a partir de
  información trazable de estas fuentes, con consolidación y revisión
  del jefe de proyecto.
\item
  Comunicación por niveles. Cada instancia de gobierno recibe la
  información de su nivel y escala solo lo que supera su umbral.
\item
  Comunicación externa a través de Ancoa. Los terceros, los
  administradores de centros comerciales y las autoridades se relacionan
  con el proyecto por medio de la gerencia de Ancoa que tiene la
  relación con ellos. Con la Comisión para el Mercado Financiero (CMF),
  esa gerencia es la filial emisora.
\end{enumerate}

El espacio colaborativo de SharePoint organizará la documentación del
proyecto, los entregables, las actas, el registro de riesgos y el
registro de cambios, con identificación de responsable, fecha de
actualización y versión vigente. Cada líder de frente actualizará la
documentación bajo su responsabilidad cuando se produzcan cambios, y el
jefe de proyecto verificará semanalmente que los registros y documentos
estén completos y actualizados. Los accesos se asignarán conforme a las
funciones de cada participante y a la clasificación de la información,
restringiendo los antecedentes confidenciales y los correspondientes al
negocio financiero a las personas autorizadas.

La Figura 6.3 presenta el plan como un mapa conceptual. La figura
relaciona los insumos, ocho componentes de decisión y los productos que
genera. Cada componente identifica una pregunta que el plan debe
resolver.

\begin{center}
\begin{minipage}{\linewidth}
\centering
\ossFigCaption{Mapa conceptual del plan de gestión de las comunicaciones.}\label{fig:6-3}

\includegraphics[width=\linewidth,height=0.66\textheight,keepaspectratio]{figuras/recepcion-20261025/1c5944e6ac44.png}

\raggedright\textbf{Fuente:} Elaboración propia a partir de PMI (2017).
\end{minipage}
\end{center}

Los componentes se agrupan en tres tipos. Los de contenido (qué
información, quién la emite, a quién se dirige y cuándo) se definen con
el registro de interesados, la estrategia de involucramiento y las
cadencias de gobierno. Los de forma de entrega (método y canal) se
adaptan a cada grupo. Por ejemplo, el personal de tienda recibirá
mensajes breves por canales disponibles, porque trabaja por turnos y
comparte terminales, y no se programarán despliegues o cambios
productivos dentro de las ventanas de congelamiento comercial, sin
introducir excepciones generales a las ventanas de congelamiento. La
comunicación y coordinación permanecerán activas durante dichos
períodos. Los componentes de reglas (restricciones y escalamiento)
responden a dos condiciones propias de la operación de Ancoa. La
información del emisor no se comparte con el retail sin autorización, y
toda desviación que supere su umbral debe llegar a tiempo a quien puede
resolverla.

El plan combinará canales formales, como informes, actas y comités, con
canales ágiles, como tableros en línea y reuniones breves, para
equilibrar la trazabilidad que exige el contrato con la interacción
continua que necesitan los equipos.

Jira mantendrá un tablero de estado del proyecto, disponible para
consulta de los representantes autorizados de Ancoa en cualquier
momento. El tablero mostrará el avance de los hitos y entregables, las
actividades pendientes, las desviaciones, los riesgos, las incidencias y
los compromisos. El jefe de proyecto supervisará su actualización y la
correspondencia entre los indicadores publicados y sus registros de
origen. Los datos operativos se revisarán diariamente, mientras que los
indicadores consolidados de desempeño se publicarán con su fecha de
corte para distinguir la información vigente de aquella correspondiente
al último cierre mensual.

\paragraph{6.1.5.2 Matriz de comunicaciones}\label{6-1-5-2-matriz-de-comunicaciones}

La matriz de comunicaciones lleva el plan a la práctica. Indica qué información se comunicará, quién la emite, quién la recibe y por qué canal. La Tabla 6.6 presenta las comunicaciones esenciales del proyecto; la frecuencia de cada una sigue las cadencias de gobierno, y la matriz se ajustará en el mes 1 con los representantes de cada grupo de interés.

Se diferenciarán las comunicaciones programadas de aquellas que deban
emitirse por la detección de una desviación relevante o un incidente
crítico de seguridad. Para estas últimas se mantendrán contactos
titulares y suplentes, así como evidencia de la fecha de detección,
emisión y recepción de los avisos.

\begin{longtable}[]{@{}
  >{\raggedright\arraybackslash}p{(\columnwidth - 6\tabcolsep) * \real{0.2500}}
  >{\raggedright\arraybackslash}p{(\columnwidth - 6\tabcolsep) * \real{0.2500}}
  >{\raggedright\arraybackslash}p{(\columnwidth - 6\tabcolsep) * \real{0.2500}}
  >{\raggedright\arraybackslash}p{(\columnwidth - 6\tabcolsep) * \real{0.2500}}@{}}
\caption{Síntesis de comunicaciones y plazos.}\label{tab:6-6}\\
\toprule\noalign{}
\begin{minipage}[b]{\linewidth}\raggedright
\textbf{Información}
\end{minipage} & \begin{minipage}[b]{\linewidth}\raggedright
\textbf{Responsable OSS}
\end{minipage} & \begin{minipage}[b]{\linewidth}\raggedright
\textbf{Destinatario}
\end{minipage} & \begin{minipage}[b]{\linewidth}\raggedright
\textbf{Canal y oportunidad}
\end{minipage} \\
\midrule\noalign{}
\endfirsthead
\toprule\noalign{}
\begin{minipage}[b]{\linewidth}\raggedright
\textbf{Información}
\end{minipage} & \begin{minipage}[b]{\linewidth}\raggedright
\textbf{Responsable OSS}
\end{minipage} & \begin{minipage}[b]{\linewidth}\raggedright
\textbf{Destinatario}
\end{minipage} & \begin{minipage}[b]{\linewidth}\raggedright
\textbf{Canal y oportunidad}
\end{minipage} \\
\midrule\noalign{}
\endhead
\bottomrule\noalign{}
\multicolumn{4}{@{}l}{\rule{0pt}{3.2ex}\textbf{Fuente:} Elaboración propia.} \\
\endlastfoot
Desempeño, valor ganado y riesgos & Jefe de proyecto & Contrapartes y
Comité Ejecutivo & Informe y sesión mensual \\
Acuerdos de los cuatro comités & Jefe de proyecto & Participantes
autorizados & Acta dentro de dos días hábiles \\
Desviación de hito superior al 10 \% & Jefe de proyecto & Administrador
y contraparte técnica & Aviso y recuperación dentro de cinco días
hábiles \\
Incidente crítico de seguridad & Encargado de seguridad & Contrapartes
de seguridad y técnica & Aviso dentro de dos horas desde detección \\
Brecha de seguridad o datos & Encargado de seguridad & Contrapartes
autorizadas & Informe preliminar dentro de 24 horas \\
Causa raíz de brecha & Encargado de seguridad & Contrapartes autorizadas
& Informe dentro de cinco días hábiles \\
Preparación de tiendas y formación & Líder de implantación & Jefaturas,
agentes y usuarios & Boletín, sesiones y ayuda por ola \\
Evidencia regulatoria y frontera & Responsables de datos y cumplimiento
con seguridad & Filial o gerencia de Ancoa & Entrega por procedimiento
aplicable \\
Migración de cartera & Líder de datos y líder funcional & Filial emisora
y responsables de atención & Plan de comunicación a titulares por ola \\
\end{longtable}

La matriz cubre a los siete grupos de interés. Dirección y control y
Propiedad reciben información consolidada; Negocios y operación y
Soporte y control operacional participan en la construcción y
validación; Operación de tienda recibe información y formación por
medios compatibles con sus turnos; y Clientes y terceros operacionales y
Organismos externos se relacionan mediante los responsables de Ancoa.
Ninguna comunicación de Only Simple Solutions sustituirá la
representación institucional de la filial ante la autoridad.

Las comunicaciones de ejecución complementarán esta síntesis. El jefe de
proyecto comunicará avance, cambios e impedimentos a los equipos de
ambas partes en el seguimiento semanal y al Comité de Proyecto
quincenal. El arquitecto y el líder de implantación coordinarán notas de
interfaz y programas de trabajos con los terceros afectados. Los
paquetes de evidencia de la frontera conservarán firma, integridad y
trazabilidad. Los equipos de desarrollo revisarán incrementos con dueños
de producto y usuarios clave al cierre del sprint y coordinarán el
trabajo diariamente; el arquitecto comunicará las decisiones aprobadas
en el repositorio y en el comité mensual. La marcha blanca tendrá sala
de control con seguimiento diario. La formación dispondrá de cápsulas,
práctica y ayuda accesible para personal de tienda. Los cambios de
interfaz y trabajos físicos se comunicarán mediante nota técnica y
programa de trabajos a los terceros afectados por medio de Ancoa;
Propiedad recibirá resúmenes por hito a través del Directorio.

La migración de los 620.000 clientes con saldo incorporará un plan de
comunicación por ola, preparado con la filial y los responsables de
atención. Precisará destinatarios, contenido, calendario y canal;
explicará los efectos para el titular y preservará la información
precontractual y evidencia exigidas. Las comunicaciones de crédito se
validarán con personas usuarias reales por su claridad, además del
  control jurídico y de seguridad (Bases Técnicas del Caso 09, 2026,
  núm. 13.3 y RT-13.12).

La entrega y actualización de las actas estará a cargo del jefe de
proyecto, con apoyo del responsable designado en cada comité. Para el
Comité Ejecutivo, el Comité de Proyecto, el Comité de Arquitectura y el
Comité de Operación se registrarán acuerdos, responsables y plazos, y
las actas se levantarán y distribuirán dentro de los dos días hábiles
siguientes a cada sesión. La publicación en SharePoint y la asociación
de los compromisos con Jira permitirán mantener su trazabilidad hasta el
cierre.

\paragraph{6.1.5.3 Gestión de las comunicaciones}\label{6-1-5-3-gestion-de-las-comunicaciones}

Durante la ejecución, la gestión de las comunicaciones pondrá en marcha el plan y la matriz. Consistirá en generar, distribuir, almacenar y recuperar la información del proyecto de forma oportuna, y en asegurar que cada mensaje llegue a su destinatario por el canal acordado y deje un registro verificable. Para mantener la coherencia entre lo planificado, lo ejecutado y sus resultados, la gestión se estructurará en tres áreas, técnica, organizacional y administrativa. Cada una tendrá un responsable, mecanismos propios y un registro que respalda lo comunicado.

\emph{\textbf{6.1.5.3.1 Área técnica}}

El área técnica coordinará la información entre los equipos de
desarrollo, plataforma, datos y calidad de Only Simple Solutions y las
áreas funcionales y técnicas de Ancoa. La coordinación estará a cargo
del arquitecto de solución y los responsables de desarrollo, con
participación de los ingenieros de solución/desarrollo y los líderes de
datos, calidad y operación según las dependencias involucradas. Su
propósito es que los equipos que construyen en paralelo trabajen con la
misma información y que TI de Ancoa conozca cada decisión que la afecta.

Los developers de cada equipo realizarán el Daily Scrum de 15 minutos
para inspeccionar el avance hacia el objetivo del sprint y adaptar su
plan. Los bloqueos que requieran coordinación externa se tratarán en el
foro correspondiente. Las dependencias entre equipos o frentes se
revisarán en una coordinación técnica apoyada por un tablero de
dependencias, con la frecuencia necesaria para evitar bloqueos. Las
decisiones técnicas que requieran aprobación arquitectónica se elevarán
al Comité de Arquitectura; las decisiones de detalle dentro de
lineamientos vigentes podrán resolverse por el arquitecto de solución y
los responsables técnicos y quedarán registradas y disponibles para
todos. Antes de construir una integración, el responsable técnico de la
integración, bajo la coordinación del arquitecto de solución, publicará
su contrato de interfaz para que los equipos y terceros trabajen sobre
la misma especificación.

Toda comunicación técnica quedará vinculada al elemento del backlog o a
la decisión que la originó. Además, el personal de TI de Ancoa
participará en las sesiones técnicas de plataforma y datos que
correspondan, de modo que el conocimiento se transfiera durante la
construcción y no únicamente al final del proyecto.

Jira concentrará el estado de las tareas, las incidencias, las
dependencias y los compromisos técnicos. Slack facilitará la
coordinación inmediata y la resolución de consultas entre los
integrantes del equipo, mientras que SharePoint conservará los
documentos técnicos aprobados y el registro de decisiones de
arquitectura. Las decisiones que afecten una integración, un entregable
o una línea base no quedarán respaldadas únicamente por una conversación
en Slack, sino que se formalizarán en el registro correspondiente.

\emph{\textbf{6.1.5.3.2 Área organizacional}}

El área organizacional coordinará la información entre las gerencias de
Ancoa, las instancias de gobierno, los equipos del proyecto y las 22
tiendas. La liderarán el jefe de proyecto y el líder de implantación y
gestión del cambio, y su propósito es que las decisiones bajen desde los
comités hasta las tiendas y que los problemas suban a tiempo hasta quien
puede resolverlos.

Las instancias de gobierno se desarrollarán conforme a las cadencias
establecidas. El Comité Ejecutivo sesionará mensualmente, el Comité de
Proyecto quincenalmente, el Comité de Arquitectura mensualmente y el
Comité de Operación mensualmente desde el mes 13. Todos estos comités
levantarán actas dentro de los dos días hábiles siguientes a cada
sesión, con las decisiones adoptadas, los acuerdos, sus responsables y
los plazos comprometidos.

La reunión semanal de seguimiento registrará impedimentos y compromisos
operativos. Los product owners revisarán los incrementos y realizarán la
validación funcional dentro de sus atribuciones, sin sustituir la
aceptación contractual, y las gerencias involucradas participarán en la
planificación de cada intervalo de planificación cuando corresponda.

El jefe de proyecto supervisará la publicación de las actas en
SharePoint y la incorporación de los compromisos a Jira, verificando que
cada acuerdo tenga un responsable y un plazo definido. Las decisiones
pendientes se mantendrán visibles en el tablero hasta su resolución, de
modo que una reunión concluida no implique el cierre de los asuntos que
permanezcan abiertos.

Only Simple Solutions se compromete a que cada una de las 22 tiendas cuente con un agente de cambio activo antes del hito H6. La designación corresponde a la jefatura de cada tienda y es una dependencia de Ancoa; Only Simple Solutions propondrá el perfil, capacitará a los agentes y escalará al Comité de Proyecto toda designación pendiente.
Cada mensaje dirigido a las tiendas indicará qué cambia, cuándo cambia y
a quién preguntar. En sentido inverso, los responsables de desarrollo,
el arquitecto de solución y el líder de implantación y gestión del
cambio escalarán de inmediato al jefe de proyecto los impedimentos o
incidentes que superen su umbral, sin esperar la sesión ordinaria.

\emph{\textbf{6.1.5.3.3 Área administrativa}}

El área administrativa coordinará los compromisos formales del proyecto
con Ancoa, los proveedores, los auditores y, a través de Ancoa, las
autoridades. La liderará el jefe de proyecto, con apoyo del encargado de
seguridad de la información, y su propósito es que cada compromiso quede
respaldado y pueda verificarse.

El jefe de proyecto entregará cada mes el informe de desempeño, con el
estado del cronograma y los hitos, avance físico y financiero, valor
ganado con sus índices de desempeño, entregables del período,
desviaciones, riesgos, incidencias, decisiones pendientes y compromisos
del período siguiente.

El informe mensual incorporará los índices de desempeño del cronograma
(SPI) y del costo (CPI), obtenidos mediante la técnica de valor ganado,
junto con la comparación entre los resultados del período y las líneas
base aprobadas. La medición se realizará sobre los paquetes de trabajo
de la EDT, utilizando el valor planificado, el valor ganado y el costo
real, conforme al procedimiento de monitoreo y control definido en la
sección 6.1.3.5. Los índices se presentarán con su período de medición y
las explicaciones de las desviaciones relevantes, evitando sustituir la
medición objetiva por declaraciones cualitativas de avance.

La elaboración del informe utilizará los registros de Jira, el
cronograma del proyecto, las evidencias documentales conservadas en
SharePoint y los antecedentes de costos autorizados. El jefe de proyecto
comprobará la consistencia de las fuentes antes de distribuir el informe
a los representantes autorizados de Ancoa. Los antecedentes financieros
de ejecución se comunicarán mediante los mecanismos contractuales
correspondientes, manteniendo la reserva de los montos de la Oferta
Económica.

Los compromisos con proveedores de nube, enlaces y equipos se
gestionarán por canales formales, con registros de recepción y
seguimiento de su desempeño, y cada hito se cerrará mediante el acta de
aceptación correspondiente.

Cuando se detecte una desviación superior al 10 \% en un hito respecto
de su línea base de control, el jefe de proyecto registrará la fecha de
detección, el hito afectado, la magnitud de la desviación, sus causas y
las consecuencias previstas sobre las dependencias del proyecto. Dentro
de los cinco días hábiles siguientes a esa detección, comunicará
formalmente la situación al administrador del contrato y a la
contraparte técnica de Ancoa, acompañando un plan de recuperación que
indique las medidas correctivas, los responsables, los plazos y los
mecanismos para verificar la recuperación del desempeño.

La comunicación de estas desviaciones será independiente del informe
mensual y de las reuniones ordinarias de los comités. Cuando se requiera
una decisión que exceda las atribuciones del jefe de proyecto, el asunto
se escalará al Comité Ejecutivo. Los planes de recuperación se
formularán respetando los plazos contractuales obligatorios y las
restricciones de operación; las modificaciones que requieran aprobación
formal seguirán el procedimiento de control integrado de cambios.

Cuando la filial o una gerencia de Ancoa deba responder a CMF, SERNAC,
SII u otra autoridad aplicable, el responsable funcional o de datos
preparará el contenido con la contraparte de cumplimiento. Seguridad
comprobará accesos, confidencialidad, integridad y trazabilidad. El jefe
de proyecto coordinará la entrega a los representantes de Ancoa que
respondan institucionalmente. Los documentos se conservarán en el
espacio colaborativo, vinculados a sus registros y con fecha, emisor y
receptor.

Ante la detección de un incidente de seguridad con indicios de
criticidad, el encargado de seguridad activará el escalamiento y
remitirá información preliminar a las contrapartes autorizadas. La
notificación de un incidente crítico se realizará en todo caso dentro de
las dos horas desde su detección, sin esperar la clasificación
definitiva. Los antecedentes se actualizarán a medida que avance la
investigación (Bases Técnicas Transversales, 2026, RT-11.18).

El aviso inicial identificará la fecha y hora de detección, los sistemas
o servicios potencialmente afectados, la severidad conocida, las medidas
de contención iniciadas y el responsable de coordinar la respuesta. La
comunicación se actualizará cuando existan nuevos antecedentes
relevantes. La convocatoria de la mesa de análisis de incidentes
definida en el gobierno interno de seguridad de Only Simple Solutions no
suspenderá el plazo de notificación ni constituirá una condición para
informar al cliente.

La clasificación definitiva, la investigación y las actividades de
recuperación seguirán el plan de respuesta a incidentes de seguridad.
Las comunicaciones relacionadas con información del negocio financiero
se realizarán con las contrapartes autorizadas de la filial emisora,
manteniendo la confidencialidad y la segregación exigidas para ese
ámbito.

Toda brecha de seguridad o de datos personales se notificará además
dentro de las 24 horas desde la detección, con informe preliminar; el
análisis de causa raíz se entregará dentro de los cinco días hábiles
siguientes. Estos controles son distintos del aviso crítico de dos horas
y se cumplirán conjuntamente cuando ambos sean aplicables (Bases
Técnicas Transversales, 2026, RT-11.19).

\emph{\textbf{Integración de las tres áreas}}

Las tres áreas no operan por separado. Conforman un sistema que conecta
los distintos niveles del proyecto. El área técnica sostiene el nivel
operativo, donde los equipos construyen y prueban; el área
organizacional conecta ese nivel con el táctico y el estratégico, donde
los comités deciden; y el área administrativa formaliza lo decidido en
informes, actas y contratos. La información fluye en ambos sentidos. Los
avances y los problemas suben desde los equipos y las tiendas hasta los
comités, y las decisiones y prioridades bajan hasta quienes deben
ejecutarlas. El PMIS del proyecto, conformado por Jira y Microsoft
SharePoint, actúa como punto de unión. Las tres áreas utilizan registros
comunes y consultan la misma información vigente.

Slack complementará este sistema como canal de comunicación cotidiana,
sin reemplazar los registros formales del proyecto. Las tareas,
incidencias y compromisos se gestionarán en Jira; la documentación, las
actas y las evidencias aprobadas se conservarán en SharePoint; y el
tablero de estado permitirá al cliente consultar la situación del
proyecto a partir de la información registrada y consolidada por sus
responsables.

\emph{\textbf{Flujo de procesos según el tipo de comunicación}}

Dentro de cualquiera de las tres áreas, cada mensaje seguirá uno de los
tres métodos de comunicación que reconoce el PMBOK, elegido según su
propósito (PMI, 2017). La Figura 6.4 muestra el flujo de procesos de
cada tipo. Primero se identifica la necesidad de comunicar y su
propósito, y luego el mensaje sigue la secuencia de pasos que
corresponde a su método.

\begin{center}
\begin{minipage}{\linewidth}
\centering
\ossFigCaption{Flujo de procesos según el tipo de comunicación.}\label{fig:6-4}

\includegraphics[width=\linewidth,height=0.66\textheight,keepaspectratio]{figuras/recepcion-20261025/9e84b9e7806f.png}

\raggedright\textbf{Fuente:} Elaboración propia a partir de PMI (2017).
\end{minipage}
\end{center}

\textbf{Comunicación interactiva.} Se usará para decidir, coordinar o
resolver asuntos mediante intercambio directo. La convocatoria de
comités incluirá agenda y datos; la sesión registrará decisiones,
responsables y plazos, y su acta se distribuirá dentro de dos días
hábiles. Las reuniones diarias, revisiones y talleres conservarán
registros de resultados y compromisos proporcionales a su finalidad, sin
imponerles automáticamente un acta de comité. Un asunto que continúe
abierto seguirá en el registro hasta comprobar su cierre.

En la rama interactiva de la Figura 6.4, el acta dentro de dos días
hábiles corresponde a las sesiones de los comités contractuales. Las
demás reuniones conservarán el registro de resultados y compromisos que
corresponda a su finalidad.

Slack se utilizará como apoyo a la comunicación interactiva entre los
equipos técnicos, particularmente para resolver consultas e impedimentos
cotidianos. Cuando una conversación produzca una decisión que afecte
actividades, responsabilidades, plazos o entregables, esta deberá
formalizarse en Jira o en el registro documental correspondiente de
SharePoint.

\textbf{Comunicación dirigida (push):} Se usará para hacer llegar
información planificada a destinatarios definidos, cuando no se requiere
una respuesta inmediata. Es el método del informe mensual de desempeño,
del boletín para las tiendas, de los avisos de cambio de interfaz y de
los paquetes de evidencia. Su proceso comienza con la generación del
contenido a partir del PMIS, sigue con su revisión y aprobación por el
responsable, continúa con la distribución por el canal definido en la
matriz y termina con la confirmación de recepción y su registro. Como
enviar no garantiza que el mensaje se lea, cada envío se verificará con
un acuse de lectura o con la confirmación del destinatario.

Este método también se aplicará a las notificaciones extraordinarias,
pero con los plazos específicos de cada situación. Los avisos sobre
desviaciones superiores al 10 \% y sobre incidentes críticos de
seguridad se emitirán sin esperar la siguiente comunicación mensual. La
fecha de detección, la emisión del aviso y la evidencia de notificación
quedarán registradas para permitir su verificación posterior.

\textbf{Comunicación disponible (pull):} Se usará para la información de
consulta frecuente o dirigida a un público amplio, que cada interesado
revisa cuando la necesita. Es el método del tablero del proyecto, del
repositorio de documentos, del registro de decisiones de arquitectura y
de las cápsulas de capacitación y ayuda en pantalla que usará el
personal nuevo de las tiendas. Su proceso comienza con la publicación
del contenido, sigue con un aviso a los interesados cuando hay una
actualización relevante, continúa con la consulta por demanda y termina
con la medición de su uso y la recolección de opiniones, que permiten
mantener el contenido vigente.

Jira y SharePoint serán los medios principales de consulta del proyecto.
El tablero de Jira mostrará el estado operativo y los indicadores de
seguimiento, mientras que SharePoint mantendrá los documentos vigentes y
las evidencias asociadas. Ambos estarán disponibles para los
representantes autorizados de Ancoa, con permisos según sus
responsabilidades y con identificación de la fecha de actualización de
la información consultada.

Los tres métodos se complementan. Por ejemplo, una modificación del
flujo de preparación de pedidos en tienda se revisa funcionalmente con
los usuarios clave durante la revisión de sprint (interactiva), se
comunica a las jefaturas antes de su habilitación (dirigida) y se
documenta en una cápsula o guía que el personal puede consultar durante
su turno (disponible). En las comunicaciones externas, los tres métodos
se aplicarán por medio de las contrapartes competentes de Ancoa.

Se verificará adicionalmente la disponibilidad de consulta del tablero
para los representantes autorizados del cliente y la vigencia de sus
registros. Las comunicaciones extraordinarias tendrán como criterio de
cumplimiento la notificación dentro de los plazos establecidos, de cinco
días hábiles desde la detección de una desviación superior al 10 \% en
un hito y dos horas desde la detección de un incidente crítico de
seguridad. Los resultados de estas verificaciones se registrarán
individualmente y cualquier incumplimiento activará una acción
correctiva.

\paragraph{6.1.5.4 Monitoreo y control de las comunicaciones}\label{6-1-5-4-monitoreo-y-control-de-las-comunicaciones}

El monitoreo verificará que la información llegue a quien la necesita, que se entienda y que produzca las decisiones o acciones esperadas, no solo que se envíe. Será un proceso continuo, que combinará el control documental del PMBOK con la revisión ágil al cierre de cada intervalo de planificación, y permitirá ajustar el plan cuando un canal, un mensaje o una frecuencia no funcione.

\emph{\textbf{Metodología de ejecución del monitoreo}}

El monitoreo se ejecutará como un ciclo de seis fases, que se repite en
cada intervalo de planificación y se acorta durante la marcha blanca. La
Figura 6.5 muestra las seis fases del ciclo y resume a la derecha sus
cadencias de revisión.

\begin{center}
\begin{minipage}{\linewidth}
\centering
\ossFigCaption{Ciclo de monitoreo y control de las comunicaciones.}\label{fig:6-5}

\includegraphics[width=\linewidth,height=0.66\textheight,keepaspectratio]{figuras/recepcion-20261025/41d36fe83bda.png}

\raggedright\textbf{Fuente:} Elaboración propia a partir de PMI (2017).
\end{minipage}
\end{center}

\textbf{Fase 1. Recolectar.} Se reúnen listas de asistencia a comités,
revisiones y talleres, actas y su fecha de publicación, acuses de
lectura de mensajes dirigidos, estadísticas de uso del tablero y del
repositorio y resultados de la encuesta breve de eficacia aplicada por intervalo de planificación a los representantes pertinentes. Los datos se obtendrán
preferentemente del PMIS y de los registros de los responsables,
reduciendo la consolidación manual y manteniendo trazabilidad de su
origen. Se añadirán las fechas que acreditan cumplimiento. El registro
identificará sesión y distribución de actas; corte y entrega de
informes; actualización del tablero; detección y notificación de
desviaciones, incidentes y brechas; y recepción de informes de causa
raíz. Cada dato conservará responsable, fuente y evidencia en el PMIS.

\textbf{Fase 2. Medir.} Se comprobará por evento el cumplimiento de los
plazos obligatorios, con objetivo de cumplimiento del 100 \%. Por
separado se medirán las metas internas de eficacia, como lectura del
boletín por agentes de cambio de al menos 90 \%, participación interna
de al menos 85 \% y claridad de al menos 4 en una escala de 1 a 5. La
trazabilidad de acuerdos tendrá meta del 100 \%. Ninguno de estos
indicadores admitirá tolerancia frente a un aviso crítico tardío, acta
fuera de plazo o condición contractual incumplida.

\textbf{Fase 3. Analizar.} Cuando un indicador queda bajo su meta, el
responsable busca la causa, por ejemplo, un canal poco usado, un mensaje
poco claro, una reunión en un horario incompatible con los turnos de
tienda o una sobrecarga de comunicaciones. El análisis distingue si el
problema está en el contenido, en el canal, en la frecuencia o en el
destinatario, porque cada caso requiere una corrección distinta.

Cuando el problema corresponda al incumplimiento de un plazo
obligatorio, se revisarán la fecha de detección, la responsabilidad
asignada, el mecanismo de escalamiento y el tiempo transcurrido hasta la
emisión de la comunicación, identificando las medidas necesarias para
evitar su repetición.

\textbf{Fase 4. Decidir.} El jefe de proyecto presenta al Comité de
Proyecto las desviaciones y las acciones correctivas propuestas, como
reforzar un canal, cambiar el formato de un mensaje o reprogramar una
reunión. El comité las aprueba, les asigna un responsable y fija un
plazo.

Las comunicaciones sujetas a plazos obligatorios se emitirán por los
responsables designados sin esperar la aprobación de acciones
correctivas en una sesión de comité. Las decisiones del comité
permitirán ajustar posteriormente los mecanismos de comunicación y
resolver problemas que requieran coordinación entre ambas partes.

\textbf{Fase 5. Ajustar.} Se aplican las acciones aprobadas. Se
actualizan la matriz de comunicaciones, los canales o la frecuencia, y
se informa el cambio a los interesados afectados. Cuando el ajuste es
relevante, se registra como una nueva versión del plan de gestión de las
comunicaciones.

\textbf{Fase 6. Registrar.} Cada hallazgo se anota en una bitácora con
su fecha, canal, observación y acción correctiva, y lo aprendido pasa al
registro de lecciones aprendidas . Con este registro comienza el
siguiente ciclo.

La bitácora se mantendrá vinculada con los registros de Jira y las
evidencias documentales de SharePoint. Para cada incumplimiento se
conservarán sus antecedentes, el responsable de corrección, el plazo
comprometido y la evidencia que permita verificar su cierre.

La cadencia del monitoreo sigue la del gobierno del proyecto. El scrum
master revisa a diario la comunicación interna de su equipo; el jefe de
proyecto revisa cada semana las actas, los acuerdos y la bitácora; los
indicadores se incluyen cada mes en el informe de desempeño y se revisan
en el Comité de Proyecto; y al cierre de cada intervalo de planificación se aplica la
encuesta de eficacia y se ajusta el plan. En marcha blanca el ciclo se
acelera. La sala de control revisará a diario la comunicación con las
tiendas, porque un problema en caja se resuelve en horas o se convierte
en venta perdida.

La Tabla 6.7 distribuye las actividades de monitoreo entre los
responsables. Ancoa participa en la revisión de las comunicaciones con
terceros y en la validación de los resultados, porque la eficacia de una
comunicación la juzga quien la recibe.

\begin{longtable}[]{@{}
  >{\raggedright\arraybackslash}p{(\columnwidth - 8\tabcolsep) * \real{0.2000}}
  >{\raggedright\arraybackslash}p{(\columnwidth - 8\tabcolsep) * \real{0.2000}}
  >{\raggedright\arraybackslash}p{(\columnwidth - 8\tabcolsep) * \real{0.2000}}
  >{\raggedright\arraybackslash}p{(\columnwidth - 8\tabcolsep) * \real{0.2000}}
  >{\raggedright\arraybackslash}p{(\columnwidth - 8\tabcolsep) * \real{0.2000}}@{}}
\caption{Distribución de actividades de monitoreo.}\label{tab:6-7}\\
\toprule\noalign{}
\begin{minipage}[b]{\linewidth}\raggedright
\textbf{Actividad}
\end{minipage} & \begin{minipage}[b]{\linewidth}\raggedright
\textbf{Responsable operativo}
\end{minipage} & \begin{minipage}[b]{\linewidth}\raggedright
\textbf{Instancia de decisión}
\end{minipage} & \begin{minipage}[b]{\linewidth}\raggedright
\textbf{Participación de Ancoa}
\end{minipage} & \begin{minipage}[b]{\linewidth}\raggedright
\textbf{Evidencia}
\end{minipage} \\
\midrule\noalign{}
\endfirsthead
\toprule\noalign{}
\begin{minipage}[b]{\linewidth}\raggedright
\textbf{Actividad}
\end{minipage} & \begin{minipage}[b]{\linewidth}\raggedright
\textbf{Responsable operativo}
\end{minipage} & \begin{minipage}[b]{\linewidth}\raggedright
\textbf{Instancia de decisión}
\end{minipage} & \begin{minipage}[b]{\linewidth}\raggedright
\textbf{Participación de Ancoa}
\end{minipage} & \begin{minipage}[b]{\linewidth}\raggedright
\textbf{Evidencia}
\end{minipage} \\
\midrule\noalign{}
\endhead
\bottomrule\noalign{}
\multicolumn{5}{@{}l}{\rule{0pt}{3.2ex}\textbf{Fuente:} Elaboración propia.} \\
\endlastfoot
Coordinación del equipo & Developers y Scrum Master según función &
Equipo & Usuarios cuando corresponda & Tablero y resultados \\
Control de informes y actas & Jefe de proyecto & Comité de Proyecto &
Contraparte técnica & Fechas y documentos \\
Comunicación con tiendas & Líder de implantación & Comité de Proyecto &
Jefaturas y agentes & Recepción y seguimiento \\
Comunicación con terceros & Jefe de proyecto & Instancia competente &
Gerencia con relación contractual & Avisos y compromisos \\
Registro de hallazgos & Responsable del frente & Jefe de proyecto &
Contraparte afectada & Bitácora \\
Ajustes del plan & Jefe de proyecto & Comité de Proyecto &
Representantes autorizados & Decisión y versión \\
Verificación del tablero & Jefe de proyecto & Comité de Proyecto &
Consulta de contraparte & Fecha de corte \\
Desviaciones y recuperación & Jefe de proyecto & Gobierno contractual &
Administrador y contraparte & Plan y evidencia de cierre \\
Avisos de seguridad y brechas & Encargado de seguridad & Respuesta a
incidentes & Contrapartes autorizadas & Detección y notificación \\
\end{longtable}

El encargado de seguridad de la información tendrá la responsabilidad
específica de verificar los plazos y registros de notificación de
incidentes críticos de seguridad. Esta responsabilidad se ejercerá
conforme a la estructura de gobierno de seguridad de Only Simple
Solutions y se coordinará con el jefe de proyecto y las contrapartes
autorizadas de Ancoa.

Si una meta interna de eficacia permanece incumplida durante dos
revisiones consecutivas, el jefe de proyecto escalará el problema al
Comité Ejecutivo. Los incumplimientos individuales de plazos u
obligaciones se registrarán y atenderán desde su detección conforme a su
criticidad, sin esperar esas revisiones. Al cierre de cada etapa, el
jefe de proyecto entregará una evaluación de comunicaciones con
resultados y medidas aplicables a operación.

Como controles mínimos de trazabilidad y reporte, Ancoa dispondrá de un
tablero de estado actualizado y accesible durante el proyecto; las actas
de los comités se emitirán dentro de los dos días hábiles siguientes y
registrarán acuerdos, responsables y plazos; toda desviación superior al
10 \% en un hito se comunicará dentro de cinco días hábiles con su plan
de recuperación; y los incidentes críticos de seguridad se notificarán a
Ancoa dentro de las dos horas siguientes a su detección. El espacio
colaborativo mantendrá disponibles los entregables, actas, registros de
riesgos y cambios y la documentación vigente.

\subsubsection{6.1.6. Gestión de adquisiciones}\label{6-1-6-gestion-de-adquisiciones}

Only Simple Solutions gestionará las adquisiciones y las dependencias de terceros como parte de la planificación y ejecución general del proyecto, considerando desde el inicio aquellos servicios, componentes o trabajos cuya disponibilidad dependa de una organización distinta del equipo directo. Esto comprende los servicios de nube, seguridad y conectividad, infraestructura local, hardware y dispositivos cuya adquisición corresponda a Ancoa, habilitaciones físicas necesarias para la implantación, integraciones con proveedores tecnológicos existentes, servicios prestados por terceros operacionales y trabajos ejecutados por subcontratistas o aliados de Only Simple Solutions.

El propósito de esta gestión será evitar que una actividad crítica quede
condicionada a un componente, servicio o proveedor sin responsable,
fecha requerida o condición de aceptación conocida. Por ello, las
dependencias se administrarán desde su identificación y especificación
hasta su recepción, integración y cierre, vinculándolas con la EDT, el
cronograma, los riesgos, las pruebas y los entregables que puedan verse
afectados.

Para efectos de gestión se distinguirán tres situaciones. La primera
corresponde a los proveedores y aliados tecnológicos asociados a la
solución propuesta, entre ellos Microsoft para los servicios sobre
Azure, Fortinet en materias de seguridad y conectividad y Dell
Technologies, directamente o mediante integradores autorizados, para
componentes de infraestructura local. La segunda comprende a los
proveedores tecnológicos que actualmente prestan servicios a Ancoa y con
los cuales será necesario coordinar accesos, interfaces, ambientes,
migraciones o actividades de transición. Finalmente se encuentran los
terceros operacionales del ecosistema de Ancoa, como vendedores de
marketplace, operadores de última milla, empresas que mantienen personal
en tienda y administradores de centros comerciales. Only Simple
Solutions no administrará la operación de estos últimos, pero sí deberá
considerar las dependencias que generan para la solución.

Independientemente de quién realice la compra, contratación o ejecución
material, cada dependencia tendrá un responsable identificado dentro de
Only Simple Solutions.

\paragraph{6.1.6.1 Planificación de adquisiciones y dependencias externas}\label{6-1-6-1-planificacion-de-adquisiciones-y-dependencias-externas}

Las necesidades externas se identificarán durante la planificación de cada frente de trabajo y antes de comprometer actividades que requieran infraestructura, equipamiento, accesos, servicios, habilitaciones o intervención de terceros. Para cada dependencia se definirá el resultado esperado, el responsable dentro de Only Simple Solutions, la contraparte correspondiente, el paquete de trabajo o entregable afectado, la fecha en que debe estar disponible y las condiciones que deberán cumplirse antes de utilizarla.

La definición incluirá además la especificación técnica aplicable, los
criterios de recepción, las evidencias que permitan comprobar su
cumplimiento, los principales riesgos asociados y, cuando la criticidad
lo justifique, una alternativa de contingencia. De esta manera, una
dependencia externa no se tratará únicamente como una fecha dentro del
cronograma, sino como una condición verificable para continuar con el
trabajo posterior.

Estas dependencias quedarán registradas en Jira y vinculadas con la
actividad o entregable correspondiente. El registro permitirá mantener
visibles el responsable, la contraparte externa, el plazo requerido, el
estado, los impedimentos y los compromisos abiertos. Slack podrá
utilizarse para la coordinación cotidiana entre los equipos, pero
cualquier acuerdo que modifique fechas, responsabilidades, condiciones
de entrega o prioridades deberá quedar registrado en Jira para conservar
su trazabilidad.

El jefe de proyecto consolidará el estado de las dependencias externas y
verificará su coherencia con el cronograma integrado. Los líderes de
cada frente, por su parte, deberán identificar oportunamente las
necesidades asociadas a sus actividades y advertir cuando una
dependencia pueda comprometer un hito o afectar el inicio de una tarea
posterior. Esta revisión será especialmente relevante antes de las
etapas de integración, pruebas, implantación y paso a producción, en las
cuales no se iniciarán trabajos dependientes mientras sus condiciones de
entrada no hayan sido verificadas.

Cuando la adquisición sea responsabilidad de Ancoa, Only Simple
Solutions proporcionará con anticipación las cantidades requeridas,
capacidades, características técnicas, interfaces, condiciones de
compatibilidad, requisitos de instalación, soporte, continuidad,
administración y demás antecedentes necesarios para que el elemento
pueda incorporarse correctamente a la solución. Cuando la provisión
corresponda a Only Simple Solutions o a un tercero contratado bajo su
responsabilidad, la selección deberá mantenerse alineada con la
arquitectura, la seguridad, la operación, el cronograma y los criterios
de calidad definidos para el proyecto.

\paragraph{6.1.6.2 Selección y evaluación de adquisiciones y dependencias externas}\label{6-1-6-2-seleccion-y-evaluacion-de-adquisiciones-y-dependencias-externas}

La definición de componentes y servicios se realizará sobre criterios técnicos verificables. La disponibilidad comercial de un producto o la sola declaración de cumplimiento de su fabricante no serán suficientes para aceptar su incorporación al proyecto; deberá comprobarse además su adecuación a la arquitectura, la operación prevista y las condiciones particulares de Ancoa.

Esta evaluación será multidisciplinaria. El arquitecto de solución
verificará la compatibilidad arquitectónica, interoperabilidad,
escalabilidad, mantenibilidad y capacidad de evolución de los
componentes. El líder de integración coordinará contratos de interfaz y
dependencias con plataformas y proveedores existentes; el líder
funcional verificará las reglas de negocio y condiciones funcionales
aplicables. El encargado de seguridad de la información revisará los
controles de acceso, segregación, protección de datos, trazabilidad y
gestión de vulnerabilidades. El líder de datos participará cuando el
componente o tercero intervenga en migraciones, almacenamiento,
intercambio, transformación o conciliación de información.

A su vez, el líder de calidad y pruebas comprobará que exista una forma
objetiva de verificar el resultado recibido y que no se incorpore a un
entregable mientras mantenga observaciones bloqueantes. El líder de
operación/SRE revisará las condiciones de soporte, monitoreo,
continuidad, recuperación y mantenimiento, mientras que el líder de
implantación y gestión del cambio coordinará la disponibilidad de
equipos, la preparación de los sitios y las habilitaciones requeridas
para los despliegues.

\textbf{Servicios Microsoft Azure}

Los servicios de nube se implementarán sobre Microsoft Azure como parte
de la arquitectura híbrida de la solución. Su selección y configuración
considerará la función que cumplen dentro de la plataforma, la
disponibilidad esperada, recuperación, segmentación, observabilidad,
administración, automatización, respaldo, seguridad y capacidad de
reversibilidad.

Las configuraciones susceptibles de automatización se mantendrán
reproducibles y versionadas, junto con la documentación necesaria para
administrar, reconstruir o transferir la operación. Se privilegiarán
servicios administrados cuando permitan disminuir complejidad
operacional sin comprometer la interoperabilidad, la seguridad ni la
capacidad futura de migración.

También se identificarán aquellos componentes cuya utilización genere
una dependencia tecnológica significativa con el proveedor. En estos
casos se conservarán los datos, configuraciones, interfaces y
procedimientos necesarios para facilitar una eventual sustitución o
transición. La participación de Microsoft o de un integrador autorizado
no modificará la responsabilidad de Only Simple Solutions por la
correcta incorporación de los servicios a la solución.

\textbf{Seguridad y conectividad}

Los componentes de seguridad y conectividad de Fortinet que resulten
seleccionados, o sus equivalentes justificados, se incorporarán de
acuerdo con las necesidades de segmentación, protección perimetral,
interconexión segura, continuidad de los enlaces y administración de la
red híbrida.

El arquitecto de solución y el encargado de seguridad de la información
verificarán que las configuraciones implementadas mantengan la
separación necesaria entre ambientes y dominios y produzcan la
trazabilidad suficiente para su operación y auditoría. Los cambios
realizados durante la implantación o posteriormente en operación
quedarán controlados y documentados.

Cuando la provisión o instalación física corresponda a Ancoa o a un integrador contratado por este, según la Tabla 6.8, Only Simple Solutions
mantendrá la responsabilidad por la especificación técnica, la
coordinación de su incorporación y la verificación de que el resultado
obtenido se ajusta a la arquitectura prevista.

\textbf{Infraestructura local y hardware}

La selección de infraestructura local podrá evaluar componentes Dell
Technologies u opciones técnicamente equivalentes según el diseño
físico. La mención de fabricante no acredita por sí sola una
contratación o alianza; el Formulario T-11 recogerá marca, modelo,
soporte y especificación de los componentes seleccionados. Only Simple
Solutions determinará los componentes requeridos, sus cantidades,
capacidades, características, interfaces, condiciones de compatibilidad,
redundancia, administración, soporte y requisitos de instalación.

Ancoa adquirirá el hardware de terreno y los dispositivos operacionales
que el alcance asigna al cliente, como terminales, lectores, impresoras
y equipamiento de tiendas y centros de distribución. Only Simple
Solutions especificará, dimensionará, coordinará y verificará su
incorporación. La provisión de los componentes on-premise de la solución
y su licenciamiento permanecerá bajo responsabilidad de Only Simple
Solutions; el inventario de implementos distinguirá por sitio lo
  adquirido por Ancoa y lo provisto por Only Simple Solutions (Bases
  Administrativas, 2026, art. 14.2; Bases Técnicas del Caso 09, 2026,
  cap. 11).

La responsabilidad de Only Simple Solutions abarcará la elaboración de
las especificaciones, el dimensionamiento y definición de cantidades, la
verificación de compatibilidad, las condiciones de preparación e
instalación, la coordinación técnica de su incorporación, la recepción y
las pruebas correspondientes. Cuando aplique, también comprenderá su
configuración antes de la entrada en servicio.

Un equipo no se considerará habilitado únicamente por haber sido
recibido físicamente. Antes de integrarlo deberá comprobarse su
correspondencia con la especificación, configuración, conectividad,
incorporación al inventario, controles de seguridad y funcionamiento en
el entorno en que será utilizado.

\textbf{Obras y habilitaciones a cargo de Ancoa}

La obra civil de separación de instalaciones es responsabilidad de Ancoa, con especificación y coordinación del proponente. Las obras eléctricas, canalizaciones y cableado de tiendas y centros de distribución las ejecutará Ancoa según la especificación y el costeo de Only Simple Solutions. En DC-01 y DC-02, Only Simple Solutions proveerá la conectividad, la seguridad y las canalizaciones, cañerías o ductos requeridos para su habilitación; ninguna exclusión general trasladará esa obligación al cliente (Bases Técnicas Transversales, 2026, RT-06.06 y RT-06.33).

La Tabla 6.8 fija el reparto de obras y suministros por sitio. Usa los
códigos del glosario de la arquitectura: DC-01 (casa matriz) y DC-02
(sitio de respaldo) alojan los componentes on-premise de la solución;
SIT-01 corresponde a las tiendas; CD-01 y CD-02, a los centros de
distribución; y CLD-01 y CLD-02, a las regiones de Azure. Este reparto es
el mismo que rige el alcance del Subdocumento 3 (SP-04) y la arquitectura
física del Subdocumento 4 (Tabla 4.11).

\begin{longtable}[]{@{}
  >{\raggedright\arraybackslash}p{(\columnwidth - 10\tabcolsep) * \real{0.1600}}
  >{\raggedright\arraybackslash}p{(\columnwidth - 10\tabcolsep) * \real{0.2300}}
  >{\raggedright\arraybackslash}p{(\columnwidth - 10\tabcolsep) * \real{0.2100}}
  >{\raggedright\arraybackslash}p{(\columnwidth - 10\tabcolsep) * \real{0.2000}}
  >{\raggedright\arraybackslash}p{(\columnwidth - 10\tabcolsep) * \real{0.2000}}@{}}
\caption{Reparto de obras y suministros por sitio.}\label{tab:6-8-reparto}\\
\toprule\noalign{}
\begin{minipage}[b]{\linewidth}\raggedright
\textbf{Sitio}
\end{minipage} & \begin{minipage}[b]{\linewidth}\raggedright
\textbf{Obras civiles y eléctricas, canalizaciones y cableado}
\end{minipage} & \begin{minipage}[b]{\linewidth}\raggedright
\textbf{Equipamiento de red y seguridad}
\end{minipage} & \begin{minipage}[b]{\linewidth}\raggedright
\textbf{Hardware y dispositivos operacionales}
\end{minipage} & \begin{minipage}[b]{\linewidth}\raggedright
\textbf{Componentes de la solución}
\end{minipage} \\
\midrule\noalign{}
\endfirsthead
\toprule\noalign{}
\begin{minipage}[b]{\linewidth}\raggedright
\textbf{Sitio}
\end{minipage} & \begin{minipage}[b]{\linewidth}\raggedright
\textbf{Obras civiles y eléctricas, canalizaciones y cableado}
\end{minipage} & \begin{minipage}[b]{\linewidth}\raggedright
\textbf{Equipamiento de red y seguridad}
\end{minipage} & \begin{minipage}[b]{\linewidth}\raggedright
\textbf{Hardware y dispositivos operacionales}
\end{minipage} & \begin{minipage}[b]{\linewidth}\raggedright
\textbf{Componentes de la solución}
\end{minipage} \\
\midrule\noalign{}
\endhead
\bottomrule\noalign{}
\multicolumn{5}{@{}l}{\rule{0pt}{3.2ex}\textbf{Fuente:} Elaboración propia a partir de las Bases Técnicas Transversales (2026, RT-06.06 y RT-06.33).} \\
\endlastfoot
\textbf{DC-01 (casa matriz) y DC-02 (sitio de respaldo)} & Obra civil de separación de instalaciones: Ancoa, con especificación de Only Simple Solutions. Conectividad, canalizaciones, cañerías y ductos: Only Simple Solutions. & Only Simple Solutions provee, instala y configura. & No aplica. & Only Simple Solutions provee, licencia, instala y configura servidores, almacenamiento y módulos de seguridad. \\ \addlinespace
\textbf{SIT-01 (tiendas)} & Ancoa ejecuta, según la especificación y el costeo de Only Simple Solutions. & Ancoa adquiere e instala, directamente o con su integrador; Only Simple Solutions especifica, configura y verifica. & Ancoa adquiere cajas, terminales, lectores e impresoras; Only Simple Solutions especifica, dimensiona y verifica. & Only Simple Solutions provee, licencia y configura el nodo local de la solución. \\ \addlinespace
\textbf{CD-01 y CD-02 (centros de distribución)} & Ancoa ejecuta, según la especificación y el costeo de Only Simple Solutions. & Ancoa adquiere e instala, directamente o con su integrador; Only Simple Solutions especifica, configura y verifica. & Ancoa adquiere lectores, impresoras de etiquetas y equipamiento del centro; Only Simple Solutions especifica, dimensiona y verifica. & Only Simple Solutions provee, licencia y configura los componentes de la solución del sitio. \\ \addlinespace
\textbf{CLD-01 y CLD-02 (Azure)} & No aplica. & Only Simple Solutions configura los servicios de seguridad de red de la nube. & No aplica. & Only Simple Solutions provee y configura los servicios. \\
\end{longtable}

En todos los sitios, el inventario de implementos distinguirá lo
adquirido por Ancoa y lo provisto por Only Simple Solutions, con su
criterio de recepción.

No obstante, Only Simple Solutions deberá definir las condiciones
técnicas que cada sitio debe cumplir. Estas podrán incluir energía,
respaldo eléctrico, conectividad, montaje, espacio, climatización,
seguridad física y otras condiciones ambientales necesarias para operar
los componentes previstos.

Antes de iniciar una instalación o despliegue dependiente de estas
habilitaciones, el líder de implantación y gestión del cambio, con apoyo
del arquitecto de solución, verificará que el sitio se encuentre
efectivamente preparado. Una obra no será considerada disponible para el
cronograma solamente porque su ejecutor la haya informado como
terminada; deberá verificarse que satisface las condiciones técnicas que
permiten comenzar la actividad siguiente.

\textbf{Proveedores tecnológicos existentes de Ancoa}

La transición deberá coexistir durante distintos períodos con las nueve
plataformas y los seis proveedores actualmente presentes en el
ecosistema tecnológico de Ancoa. A ello se suma el levantamiento y
rediseño progresivo de las catorce interfaces punto a punto existentes,
por lo que la gestión de estos proveedores formará parte explícita del
control de dependencias del proyecto.

Para cada caso se identificarán el sistema y proveedor involucrado, el
responsable técnico interno, la contraparte de Ancoa, los accesos o
información necesarios, la interfaz afectada, los ambientes requeridos,
datos de prueba, mecanismos de autenticación, ventanas de intervención y
criterios que deberán satisfacerse antes de modificar o retirar una
integración existente.

Cuando Only Simple Solutions no mantenga una relación contractual
directa con el proveedor, Ancoa será responsable de efectuar las
gestiones administrativas o comerciales necesarias. Only Simple
Solutions conservará la responsabilidad por la dependencia técnica.
Deberá especificarla, preparar las actividades asociadas, efectuar
seguimiento y advertir oportunamente cualquier impedimento que pueda
afectar el proyecto.

Por esta razón, ninguna actividad crítica se mantendrá simplemente con
un estado de ``pendiente del proveedor''. Toda dependencia abierta
deberá contar con un responsable, una fecha requerida, un impacto
identificado y una acción de seguimiento.

\textbf{Terceros operacionales del ecosistema de Ancoa}

La solución también deberá interactuar con organizaciones que participan
en la operación de Ancoa sin formar parte del alcance directo de Only
Simple Solutions. En estos casos la gestión se concentrará en las
interfaces, condiciones técnicas y responsabilidades que puedan afectar
el funcionamiento de la solución, sin asumir la administración de la
operación propia del tercero.

Los vendedores del marketplace continuarán operando sus plataformas.
Only Simple Solutions gestionará las interfaces necesarias para el
intercambio de catálogo, pedidos, estados, devoluciones y demás procesos
contemplados, además de la información requerida para medir su
desempeño. Las empresas de última milla seguirán siendo responsables del
transporte de los pedidos; la solución deberá integrarse con ellas para
recibir y mantener la trazabilidad de los eventos asociados a la
entrega.

La relación con los administradores de centros comerciales permanecerá
bajo responsabilidad de Ancoa. Sin perjuicio de ello, la arquitectura
deberá considerar que parte de la conectividad o infraestructura de una
tienda puede depender de estos terceros, por lo que una indisponibilidad
de sus servicios no deberá impedir las funciones críticas previstas para
la operación local.

Cuando intervenga personal perteneciente a proveedores de Ancoa, como
los repositores externos, Only Simple Solutions podrá establecer los
requisitos técnicos, operacionales y de acceso necesarios para utilizar
la solución, pero su implementación deberá coordinarse mediante Ancoa y
la empresa empleadora correspondiente. No se asumirán facultades de
dirección laboral sobre este personal. Los accesos que se requieran
serán individualizados y se otorgarán de acuerdo con las funciones
efectivamente necesarias.

\textbf{Subcontratistas y aliados de Only Simple Solutions}

La participación de subcontratistas o especialistas externos se
realizará bajo responsabilidad de Only Simple Solutions y con las
autorizaciones que correspondan. Cada intervención deberá tener un
alcance delimitado, responsable interno, productos de trabajo, fechas,
criterios de recepción y condiciones de acceso claramente establecidos.

La subcontratación requerirá autorización previa y escrita de Ancoa y
respetará el límite del 40 \% del valor del contrato. La custodia u
operación de información crítica por terceros estará sujeta a
autorización expresa; cualquier acceso a datos exigirá declaración,
auditoría y acuerdo de tratamiento. Only Simple Solutions mantendrá su
responsabilidad contractual y no presentará como subcontratista a un
  proveedor propio de Ancoa (Bases Administrativas, 2026, art. 73).

Los productos elaborados por estos terceros estarán sujetos a los mismos
controles de seguridad, calidad, confidencialidad, trazabilidad y
documentación aplicados al equipo propio. La utilización de un
especialista externo no transferirá la responsabilidad del resultado
frente a Ancoa.

Cuando sea necesario otorgar acceso a ambientes o información, estos
accesos serán nominativos, sujetos a mínimo privilegio y limitados al
período estrictamente requerido. Al finalizar la participación del
tercero deberán revocarse las credenciales y verificarse la devolución o
eliminación de la información que corresponda.

\paragraph{6.1.6.3 Supervisión y control de dependencias externas}\label{6-1-6-3-supervision-y-control-de-dependencias-externas}

Las adquisiciones y dependencias externas se supervisarán dentro del mismo modelo de seguimiento utilizado para las demás actividades del proyecto. En Jira se mantendrá el estado de cada una, distinguiendo si se encuentra identificada, especificada, solicitada, en provisión o ejecución, recibida, en verificación, aceptada, cerrada o afectada por una desviación.

Las dependencias que puedan afectar actividades próximas serán revisadas
en la reunión semanal de seguimiento. Cada responsable técnico mantendrá
la coordinación cotidiana con su contraparte, mientras que el jefe de
proyecto concentrará aquellas situaciones que afecten a más de un
frente, comprometan hitos o requieran intervención de Ancoa.

Ante una desviación se deberá registrar qué condición no se cumplió,
cuáles son las actividades o entregables expuestos, qué acción se
requiere, quién debe ejecutarla y en qué fecha. Cuando exista una
alternativa de contingencia, esta también quedará registrada.

El escalamiento utilizará las instancias de gobierno ya definidas para
el proyecto. Las dependencias técnicas o arquitectónicas podrán ser
revisadas en el Comité de Arquitectura; aquellas que afecten el
cronograma, entregables o decisiones de ejecución podrán tratarse en el
Comité de Proyecto. Cuando la situación requiera una decisión
estratégica o exceda las atribuciones operativas, será elevada al Comité
Ejecutivo. Durante la etapa de operación, las dependencias permanentes
de servicios externos también podrán revisarse en el Comité de Operación
junto con los incidentes, problemas y niveles de servicio asociados.

Los cambios informados por fabricantes y proveedores se evaluarán antes
de ser incorporados a los ambientes del proyecto. Esto incluye nuevas
versiones, modificaciones de interfaz, cambios relevantes de
configuración, vulnerabilidades, discontinuidades o término de soporte.
Si una modificación afecta una línea base, un requisito, la
arquitectura, un entregable o un compromiso del proyecto, deberá
gestionarse mediante el proceso de control integrado de cambios.

\paragraph{6.1.6.4 Recepción técnica y criterios de aceptación}\label{6-1-6-4-recepcion-tecnica-y-criterios-de-aceptacion}

La recepción de una compra, la terminación de una obra o la declaración de cumplimiento emitida por un proveedor no implicarán automáticamente que la dependencia se encuentra satisfecha. Todo componente, servicio o trabajo externo deberá superar una verificación técnica antes de ser incorporado al proyecto.

La Tabla 6.9 sintetiza los criterios por tipo de dependencia. Estos se
definirán antes de solicitar, adquirir o contratar el resultado
correspondiente y se ajustarán a la naturaleza de cada dependencia. En
términos generales, la recepción deberá comprobar su correspondencia con
la especificación aprobada, funcionamiento esperado, condiciones de
seguridad, capacidad de integración y documentación necesaria para
operarlo.

\begin{longtable}[]{@{}
  >{\raggedright\arraybackslash}p{(\columnwidth - 6\tabcolsep) * \real{0.2500}}
  >{\raggedright\arraybackslash}p{(\columnwidth - 6\tabcolsep) * \real{0.2500}}
  >{\raggedright\arraybackslash}p{(\columnwidth - 6\tabcolsep) * \real{0.2500}}
  >{\raggedright\arraybackslash}p{(\columnwidth - 6\tabcolsep) * \real{0.2500}}@{}}
\caption{Criterios de recepción de adquisiciones y trabajos de terceros.}\label{tab:6-8}\\
\toprule\noalign{}
\begin{minipage}[b]{\linewidth}\raggedright
\textbf{Tipo de dependencia}
\end{minipage} & \begin{minipage}[b]{\linewidth}\raggedright
\textbf{Responsable Only Simple Solutions}
\end{minipage} & \begin{minipage}[b]{\linewidth}\raggedright
\textbf{Criterio de recepción}
\end{minipage} & \begin{minipage}[b]{\linewidth}\raggedright
\textbf{Evidencia principal}
\end{minipage} \\
\midrule\noalign{}
\endfirsthead
\toprule\noalign{}
\begin{minipage}[b]{\linewidth}\raggedright
\textbf{Tipo de dependencia}
\end{minipage} & \begin{minipage}[b]{\linewidth}\raggedright
\textbf{Responsable Only Simple Solutions}
\end{minipage} & \begin{minipage}[b]{\linewidth}\raggedright
\textbf{Criterio de recepción}
\end{minipage} & \begin{minipage}[b]{\linewidth}\raggedright
\textbf{Evidencia principal}
\end{minipage} \\
\midrule\noalign{}
\endhead
\bottomrule\noalign{}
\multicolumn{4}{@{}l}{\rule{0pt}{3.2ex}\textbf{Fuente:} Elaboración propia.} \\
\endlastfoot
Servicios Microsoft Azure & Arquitecto de Solución y Líder de Operación
& Arquitectura, configuración, disponibilidad, seguridad, observabilidad
y reversibilidad verificadas & Configuración versionada e informes de
prueba \\
Seguridad y conectividad & Arquitecto de Solución y Encargado de
Seguridad de la Información & Segmentación, conectividad, controles de
seguridad y monitoreo verificados & Configuraciones y resultados de
pruebas \\
Hardware e infraestructura local & Arquitecto de Solución y Líder de
Implantación y Gestión del Cambio & Correspondencia con especificación,
dimensionamiento, compatibilidad y funcionamiento & Inventario técnico,
certificados y pruebas \\
Obras y habilitaciones a cargo de Ancoa & Líder de Implantación y
Gestión del Cambio & Sitio preparado para instalar y operar los
componentes previstos & Acta de recepción técnica y resultados de
verificación \\
Integraciones con proveedores de Ancoa & Arquitecto de Solución &
Contrato de interfaz, autenticación, trazabilidad, manejo de errores y
funcionamiento extremo a extremo verificados & Casos de prueba,
registros y conciliaciones \\
Trabajo de subcontratistas o aliados & Líder del frente correspondiente
y Líder de Calidad & Producto completo, probado, documentado y sin
observaciones bloqueantes & Entregable, pruebas y registro de
conformidad \\
\end{longtable}

En un servicio de Azure deberá comprobarse que la configuración
implementada corresponde a la arquitectura definida y que puede ser
monitoreada, administrada, respaldada y recuperada. Para los componentes
de seguridad y conectividad se verificará, entre otros aspectos, la
segmentación, los accesos, la comunicación entre componentes y la
generación de registros suficientes para supervisión.

En el caso de hardware e infraestructura local, la recepción deberá
demostrar que el componente corresponde a lo especificado y funciona
correctamente en el entorno objetivo. Para las obras y habilitaciones se
comprobará que el sitio proporciona las condiciones requeridas para
instalar y operar el equipamiento previsto. Las integraciones externas
deberán demostrar intercambio correcto de información, autenticación,
tratamiento de errores, trazabilidad, reintentos y conciliación cuando
corresponda.

Los productos elaborados por subcontratistas o aliados deberán demostrar
el mismo nivel de completitud que un trabajo ejecutado directamente por
Only Simple Solutions. Los certificados de fabricante o proveedor podrán
utilizarse como parte de la evidencia, pero no sustituirán las pruebas
que deban realizarse en el contexto específico de Ancoa.

La recepción técnica de una dependencia tampoco equivaldrá por sí sola a
la aceptación del entregable contractual en el cual participe. Cuando un
producto o servicio externo forme parte de un entregable, Only Simple
Solutions incorporará las evidencias de verificación, trazabilidad y
observaciones resueltas antes de presentarlo a la contraparte técnica de
Ancoa.

\paragraph{6.1.6.5 Gestión de incumplimientos y contingencias}\label{6-1-6-5-gestion-de-incumplimientos-y-contingencias}

Cuando una dependencia externa presente atrasos, defectos o indisponibilidad, el responsable técnico evaluará su impacto sobre las actividades posteriores y propondrá las acciones de recuperación disponibles. Dependiendo de la situación, podrán utilizarse componentes técnicamente equivalentes, alternativas temporales aprobadas, reordenamiento de actividades, anticipación de trabajos independientes o coexistencia temporal con el componente anterior.

Cuando la resolución dependa de un proveedor de Ancoa, el impedimento
será registrado y escalado oportunamente, pero no se considerará fuera
de la gestión de Only Simple Solutions. El responsable deberá mantener
el seguimiento, documentar el impacto y proponer todas las medidas de
mitigación que se encuentren bajo control del proyecto.

Para las dependencias más críticas se identificarán alternativas que
reduzcan la exposición a un único proveedor, servicio o componente
cuando ello sea técnicamente viable. Ninguna contingencia podrá
utilizarse para reducir las condiciones de seguridad, segregación,
calidad, continuidad o aceptación definidas para la solución. Si la
alternativa propuesta modifica una decisión arquitectónica, requisito o
línea base, deberá gestionarse mediante el control integrado de cambios.

\paragraph{6.1.6.6 Cierre y salida de terceros}\label{6-1-6-6-cierre-y-salida-de-terceros}

La participación de un proveedor, subcontratista o aliado se considerará cerrada cuando, además de finalizar su trabajo, se hayan completado las obligaciones técnicas y documentales asociadas. Esto comprende la entrega de documentación, configuraciones, código o scripts que formen parte de la solución, la incorporación de los elementos correspondientes al inventario, la resolución de observaciones, la transferencia de conocimiento y el cierre de incidencias que permanezcan abiertas.

También deberá verificarse la revocación de accesos y credenciales
temporales y la devolución o eliminación controlada de la información
que haya sido utilizada por el tercero.

En los servicios de nube y servicios administrados se conservarán además
los datos, configuraciones, respaldos, definiciones de infraestructura y
antecedentes necesarios para mantener la continuidad o realizar una
transición futura. El objetivo es evitar que un componente crítico quede
condicionado de manera exclusiva al conocimiento, credenciales o
documentación mantenidos por un tercero.

\paragraph{6.1.6.7 Responsabilidades}\label{6-1-6-7-responsabilidades}

El jefe de proyecto consolidará las adquisiciones y dependencias externas dentro de la planificación general y coordinará aquellas que puedan afectar hitos o requieran escalamiento. El arquitecto de solución tendrá a su cargo la evaluación arquitectónica de servicios de nube, infraestructura y componentes tecnológicos. El líder de integración coordinará las dependencias de interfaz con las plataformas y proveedores existentes, y el líder funcional verificará las reglas de negocio y condiciones funcionales asociadas.

El encargado de seguridad de la información verificará los controles de
seguridad, segregación, accesos, trazabilidad y protección de
información. El líder de datos participará cuando la dependencia
involucre intercambio, migración, transformación o conciliación de
información, prestando especial atención a la separación entre los
ámbitos comercial y financiero. El líder de calidad y pruebas verificará
la existencia de criterios y evidencias suficientes antes de incorporar
un resultado externo a un entregable, mientras que el líder de
operación/SRE revisará las condiciones de monitoreo, soporte,
continuidad y recuperación necesarias para mantener el servicio.

Por su parte, el líder de implantación y gestión del cambio coordinará
la preparación de los sitios, disponibilidad de equipamiento y
habilitaciones necesarias para los despliegues. Los líderes de cada
frente conservarán la responsabilidad sobre las dependencias externas
utilizadas en sus productos de trabajo, aunque su ejecución material
corresponda a otra organización.

Ancoa mantendrá las compras de terreno, obras civiles y gestiones
contractuales que le asigna el alcance contractual. Only Simple
Solutions conservará la provisión de componentes de solución bajo su
responsabilidad y la conectividad, seguridad y canalizaciones de DC-01 y DC-02 según la Tabla 6.8, además de especificación, coordinación, recepción y configuración
comprometidas. El inventario por sitio identificará responsable de
provisión y criterio de recepción; ningún componente se considerará
disponible solo por una comunicación de su comprador.

De esta forma, toda dependencia externa tendrá una responsabilidad
definida desde su planificación hasta su recepción, integración y
cierre, evitando que las actividades desarrolladas fuera del equipo
directo queden fuera del control del proyecto.

La Oferta Técnica describirá las características técnicas, cantidades
requeridas, capacidades, responsabilidades, dependencias y criterios de
aceptación necesarios para ejecutar el proyecto. No incluirá precios,
tarifas, valores unitarios ni antecedentes que permitan inferir la
valorización económica de los componentes o servicios considerados.

\subsubsection{6.1.7. Mecanismos de decisión y cadencias de gobierno del proyecto}\label{6-1-7-mecanismos-de-decision-y-cadencias-de-gobierno-del-proyecto}

Las instancias descritas en 6.1.2 funcionarán con una cadencia estable y con límites de autoridad explícitos. Las ceremonias ágiles, reuniones técnicas y coordinaciones de trabajo complementarán este gobierno, pero no crearán niveles paralelos de aprobación. Los equipos podrán resolver el detalle de ejecución dentro de las líneas base y decisiones vigentes; los asuntos que excedan sus atribuciones se escalarán a la instancia correspondiente. La Tabla 6.10 compara sus cadencias y ámbitos de decisión.

\begin{longtable}[]{@{}
  >{\raggedright\arraybackslash}p{(\columnwidth - 4\tabcolsep) * \real{0.3333}}
  >{\raggedright\arraybackslash}p{(\columnwidth - 4\tabcolsep) * \real{0.3333}}
  >{\raggedright\arraybackslash}p{(\columnwidth - 4\tabcolsep) * \real{0.3333}}@{}}
\caption{Cadencias y ámbitos de decisión del proyecto.}\label{tab:6-9}\\
\toprule\noalign{}
\begin{minipage}[b]{\linewidth}\raggedright
\textbf{Instancia}
\end{minipage} & \begin{minipage}[b]{\linewidth}\raggedright
\textbf{Cadencia}
\end{minipage} & \begin{minipage}[b]{\linewidth}\raggedright
\textbf{Ámbito principal}
\end{minipage} \\
\midrule\noalign{}
\endfirsthead
\toprule\noalign{}
\begin{minipage}[b]{\linewidth}\raggedright
\textbf{Instancia}
\end{minipage} & \begin{minipage}[b]{\linewidth}\raggedright
\textbf{Cadencia}
\end{minipage} & \begin{minipage}[b]{\linewidth}\raggedright
\textbf{Ámbito principal}
\end{minipage} \\
\midrule\noalign{}
\endhead
\bottomrule\noalign{}
\multicolumn{3}{@{}l}{\rule{0pt}{3.2ex}\textbf{Fuente:} Elaboración propia.} \\
\endlastfoot
Comité Ejecutivo & Mensual & Decisiones estratégicas, asuntos escalados,
cambios de alcance y riesgos mayores \\
Comité de Proyecto & Quincenal & Cronograma, entregables, desviaciones,
riesgos, dependencias y decisiones operativas \\
Comité de Arquitectura & Mensual & Decisiones de arquitectura, deuda
técnica, riesgos y controversias técnicas \\
Comité de Operación & Mensual desde el mes 13 & Niveles de servicio,
incidentes, problemas, continuidad y mejora continua \\
Reunión semanal de seguimiento & Semanal & Impedimentos, dependencias y
compromisos de ejecución \\
Ceremonias Scrum & Según el ciclo de dos semanas & Planificación,
coordinación, revisión y mejora del trabajo del sprint \\
\end{longtable}

Las decisiones se registrarán con responsable y plazo. Los asuntos que
no puedan resolverse en una instancia se elevarán a la siguiente sin
asumir aprobación por silencio. Cuando una decisión afecte una línea
base, un entregable o un compromiso aprobado, se iniciará el control
integrado de cambios definido en 6.1.3.6 antes de ejecutar la
modificación.

\paragraph{6.1.7.1. Gobierno por etapas y transición a operación}\label{6-1-7-1-gobierno-por-etapas-y-transicion-a-operacion}

El énfasis de las instancias de gobierno evolucionará con el proyecto. Durante los primeros meses predominarán el levantamiento, la arquitectura, el desarrollo, la integración y la preparación de infraestructura. A medida que se acerque la marcha blanca de la Etapa 1, aumentará la participación de operación e implantación y el Comité de Operación comenzará a revisar soporte, observabilidad, continuidad, procedimientos y preparación de los equipos.

Antes de la producción de la Etapa 1 en el mes 16 y de la Etapa 2 en el
mes 21, el jefe de proyecto reunirá las evaluaciones de calidad,
seguridad, datos, integración, implantación y operación. El Comité de
Proyecto revisará la preparación y la contraparte técnica formalizará la
conformidad que corresponda. Si existe una condición bloqueante, se
aplicará el plan de recuperación o reversión necesario sin sustituir el
procedimiento de aceptación.

Desde el mes 21, con ambas etapas en producción, el foco del gobierno se
desplazará progresivamente hacia niveles de servicio, incidentes,
problemas, capacidad, continuidad, cambios evolutivos y transferencia de
conocimiento. El Comité Ejecutivo mantendrá la visión estratégica y
contractual; el Comité de Proyecto continuará coordinando los
entregables y cambios transversales que correspondan; y el Comité de
Operación asumirá el seguimiento ordinario del servicio hasta el cierre
del mes 56.

\subsection{6.2 Metodología de Desarrollo Software}\label{6-2-metodologia-de-desarrollo-software}

La metodología de desarrollo de software describe cómo Only Simple
Solutions construirá, integrará, probará y entregará la solución de
Multitiendas Ancoa. Se basa en equipos Scrum estables, organizados por
responsabilidades de negocio y plataforma, con sprints de dos semanas.
Para coordinar dependencias se propone una revisión conjunta cada cuatro
sprints, equivalente a ocho semanas; esa cadencia se relacionará con el
calendario contractual sin tratarla como dos meses exactos. Se adoptan
prácticas seleccionadas de coordinación inspiradas en SAFe, sin crear
órganos paralelos de aprobación (Schwaber y Sutherland, 2020; Scaled
Agile, Inc., 2025).

El desarrollo incorporará DevSecOps, integración y entrega continuas,
infraestructura como código y pruebas automatizadas. Cada incremento
deberá aportar evidencia trazable de calidad y seguridad. La producción
quedará sujeta a las conformidades contractuales, ventanas comerciales y
separación entre Retail y Emisor, con independencia del ritmo de los
sprints.

La sección se organiza en siete partes. La sección 6.2.1 presenta el
enfoque de desarrollo y sus implicancias; las secciones 6.2.2 y 6.2.3
explican cómo ese enfoque gestiona los requerimientos, la arquitectura
evolutiva, la refactorización y la deuda técnica; las secciones 6.2.4 y
6.2.5 describen las prácticas de DevSecOps, infraestructura como código,
entrega continua y automatización de pruebas; la sección 6.2.6 define la
organización de los equipos con sus ceremonias, artefactos, cadencias y
mecanismos de decisión; y la sección 6.2.7 explica cómo se coordinan las
dos etapas para cumplir el tiempo de salida al mercado. El Formulario
T-10 recogerá el detalle operativo de estas prácticas.

\subsubsection{6.2.1. Enfoque de desarrollo y ciclos de entrega}\label{6-2-1-enfoque-de-desarrollo-y-ciclos-de-entrega}

Only Simple Solutions adoptará un enfoque híbrido para el desarrollo de software, combinando la planificación y el control formal de los entregables con la construcción incremental mediante Scrum. Esta elección responde a dos condiciones del proyecto. El cronograma contractual establece fechas obligatorias que no pueden modificarse, mientras que el entorno tecnológico actual presenta incertidumbres que deberán resolverse durante el levantamiento y la construcción de la solución.

La operación de Ancoa se sostiene sobre nueve plataformas de seis
proveedores, conectadas mediante catorce interfaces punto a punto. La
información disponible sobre sus reglas, intercambios de datos y
dependencias deberá verificarse para definir cómo realizar la transición
sin interrumpir la operación. Por ello, la planificación establecerá los
hitos, las dependencias y los criterios de aceptación, mientras que las
iteraciones permitirán comprobar supuestos, validar reglas de negocio y
ajustar el detalle de implementación a los hallazgos.

Los resultados de estas validaciones se incorporarán al backlog y a la
planificación del trabajo siguiente. Los ajustes se realizarán dentro
del alcance aprobado, preservando las fechas de producción de la Etapa 1
en el mes 16 y de la Etapa 2 en el mes 21. Cuando un hallazgo amenace
esos compromisos, se evaluará su impacto y se establecerán acciones de
recuperación, conforme al control integrado de cambios definido en la
sección 6.1.3.6.

La construcción se organizará en sprints de dos semanas. Al comienzo de cada intervalo de planificación, de cuatro sprints, los equipos y dueños
de producto acordarán objetivos, dependencias y contribución a hitos; en
cada sprint seleccionarán trabajo según prioridad, capacidad disponible
y condiciones de entrada. Las actividades de planificación, revisión y
mejora se incluirán en la capacidad, en vez de suponer que todas las
horas nominales son construcción. La cadencia de coordinación será de
cuatro sprints, equivalentes a ocho semanas, y se programará por fechas
y recursos manteniendo los hitos contractuales.

Durante cada sprint se comprobarán los supuestos técnicos y las reglas
de negocio relacionados con el trabajo seleccionado. Al finalizar, el
equipo presentará el incremento probado a los representantes de Ancoa
que correspondan a los procesos afectados. Sus observaciones y los
hallazgos técnicos permitirán precisar el trabajo de los siguientes
sprints dentro del alcance aprobado, y la retrospectiva permitirá
acordar mejoras en la forma de trabajar del equipo.

Cada incremento deberá cumplir la definición de terminado, con código
revisado, pruebas y controles aplicables satisfactorios, documentación
actualizada y evidencia trazable a los requerimientos. La entrega de un
incremento no implicará su puesta inmediata en producción. Las versiones
se liberarán conforme a los criterios de aceptación, las aprobaciones y
las ventanas de congelamiento establecidas para el proyecto.

Este enfoque tiene implicancias directas en cinco ámbitos, la gestión de
requerimientos, la arquitectura evolutiva, la refactorización, la deuda
técnica y el tiempo de salida al mercado. La Figura 6.6 resume la
condición del proyecto que hace exigente cada ámbito y la respuesta que
la metodología desarrolla en las secciones siguientes.

\begin{center}
\begin{minipage}{\linewidth}
\centering
\ossFigCaption{Implicancias del enfoque de desarrollo.}\label{fig:6-6}

\includegraphics[width=\linewidth,height=0.66\textheight,keepaspectratio]{figuras/recepcion-20261025/dc7d6b56c8a3.png}

\raggedright\textbf{Fuente:} elaboración propia.
\end{minipage}
\end{center}

La figura relaciona cada condición con su mecanismo de control. El
alcance se detalla en historias, la arquitectura se gobierna mediante
decisiones y pruebas, el código se refactoriza con regresiones, la deuda
se registra y la liberación respeta ventanas. El orden de construcción
podrá adaptarse dentro del alcance y las líneas base, sin flexibilizar
requisitos obligatorios, calidad ni fechas.

\subsubsection{6.2.2. Gestión de requerimientos y trazabilidad}\label{6-2-2-gestion-de-requerimientos-y-trazabilidad}

La gestión partirá del catálogo v3.0 y del alcance de la solución. Su hoja principal contiene 223 requerimientos funcionales, 75 no funcionales y 9 obligaciones del proponente, que suman 307 elementos. Los 374 requisitos transversales se gestionarán por su identificador y carácter obligatorio, deseable o según corresponda, sin sumarlos indiscriminadamente como un único catálogo de comportamientos.

El catálogo corresponde a la versión de Only Simple Solutions (2026).
Toda actualización conservará equivalencias de identificadores,
historial de cambios y trazabilidad de la oferta.

La aprobación de H1 formalizará la línea base del levantamiento y su
trazabilidad. Los equipos precisarán cómo cumplir cada elemento
aplicable y tramitarán formalmente los cambios del alcance aprobado. El
Formulario T-12 identificará la declaración de cumplimiento, el
componente o práctica, la ubicación de su desarrollo y la evidencia
prevista.

Cada requisito se relacionará con el producto o servicio responsable y
con el trabajo que lo satisface, conservando su identificador oficial en
la cadena. La Figura 6.7 distingue la trazabilidad declarada en la
oferta de los registros y resultados que se producirán durante
ejecución. Las obligaciones del proponente se vincularán con entregables
documentales o análisis, no necesariamente con historias funcionales de
software.

\begin{center}
\begin{minipage}{\linewidth}
\centering
\ossFigCaption{Trazabilidad del requisito a la evidencia de aceptación.}\label{fig:6-7}

\includegraphics[width=\linewidth,height=0.66\textheight,keepaspectratio]{figuras/recepcion-20261025/43633f99f021.png}

\raggedright\textbf{Fuente:} Elaboración propia.
\end{minipage}
\end{center}

La trazabilidad funcionará en ambos sentidos. Hacia adelante comprobará
que cada elemento aplicable tenga trabajo y método de verificación;
hacia atrás permitirá llegar desde una evidencia al requisito y versión
que la originaron. Jira conservará relaciones de gestión; el repositorio
vinculará cambios con su trabajo; y el pipeline publicará artefacto,
versión, ambiente y resultado de los controles. La integración de esas
herramientas se habilitará y comprobará con casos completos, vínculos
ausentes y pruebas fallidas, para que la matriz advierta incumplimientos
en vez de certificar únicamente la existencia de enlaces.

La matriz de oferta describirá cumplimiento y evidencia futura; durante
ejecución se conservarán resultados verificables, observaciones y actas.
SharePoint mantendrá las versiones aprobadas y sus paquetes de
aceptación con permisos por dominio, en correspondencia con el
Formulario T-12.

Una historia solo entrará a un sprint cuando cumpla la definición de
preparado. La historia está vinculada a su requisito, tiene criterios de
aceptación escritos en formato Dado, Cuando, Entonces y acordados con el
product owner, identifica sus dependencias con otros equipos o terceros,
cuenta con datos de prueba y declara si opera en el contexto Retail, en
el Emisor o en la frontera entre ambos. Esta última condición responde a
la separación operativa entre Retail y Emisor. Toda historia que cruce
información entre los dos negocios requerirá, antes de planificarse, la
evaluación documentada de datos, seguridad y la contraparte de cumplimiento designada por Ancoa, y no se liberará sin la ficha de intercambio aprobada según la ruta única de la sección 6.1.2.

Los criterios de aceptación se escriben con la estructura Dado, Cuando,
Entonces, propia del desarrollo guiado por el comportamiento (Cucumber,
s.f.-a). Cada parte cumple una función. \textbf{Dado} describe el
contexto o el estado inicial del que parte el escenario; \textbf{Cuando}
indica la acción o el evento que ejecuta el usuario o el sistema; y
\textbf{Entonces} establece el resultado observable que debe producirse
para dar por cumplido el criterio. Escritos así, los criterios usan el
  lenguaje del negocio y permiten acordar resultados observables. Para
ejecutarlos como pruebas, los developers implementarán los pasos, datos
y aserciones correspondientes; el formato escrito no genera por sí solo
una prueba automática (Cucumber, s.f.-a, s.f.-b). Por ejemplo, para la
evaluación crediticia en la caja.

\textbf{Dado} que existe una solicitud de evaluación crediticia completa en el
contexto del Emisor y los datos necesarios para evaluarla,
\textbf{Cuando} se solicita su evaluación desde el punto de venta,
\textbf{Entonces} el servicio devuelve una decisión y el cupo que
corresponda, conservando la trazabilidad de la evaluación. La prueba de
desempeño verificará un P95 no superior a 8 segundos bajo la carga
objetivo. La aceptación de un crédito seguirá exigiendo información
precontractual entregada y evidencia del consentimiento; no se presupone
que una autorización de compra con cupo vigente sea la misma operación.

Cada historia reúne varios escenarios de este tipo, incluidos los de
rechazo o de error, y el product owner los acuerda con el equipo antes
de que la historia entre al sprint. En el ejemplo, el resultado esperado
incorpora además la meta de desempeño del requisito no funcional. Su
comprobación utilizará la campaña de carga en un ambiente representativo
y publicará el resultado por versión; una regresión rápida del pipeline
no sustituirá esa evidencia.

Para comprobar estos requerimientos de forma continua se usarán funciones de aptitud. Una función de aptitud es una prueba automatizada que mide si la solución mantiene una característica de arquitectura comprometida, como un tiempo de respuesta, la separación de datos entre Retail y Emisor o la compatibilidad de un contrato de interfaz, y entrega un resultado verificable por versión (Ford et al., 2017). Las funciones rápidas se ejecutarán en cada integración del pipeline; las de desempeño, autonomía y resiliencia, en las campañas de los ambientes representativos.

Los requerimientos no funcionales se comprobarán progresivamente. Las
regresiones rápidas del pipeline protegerán contratos y reglas
estructurales; las campañas de carga, autonomía, resiliencia,
recuperación y aceptación se ejecutarán en ambientes representativos
según el plan de pruebas. El desempeño se medirá de extremo a extremo en
P95 bajo el peak declarado, incluidos evaluación crediticia de hasta 8
segundos, confirmación del pedido durante el evento de hasta 3 segundos
y propagación del precio de hasta 5 minutos (Bases Técnicas
Transversales, 2026, núm. 1.7 y 9.1; Bases Técnicas del Caso 09, 2026,
RT-09.01).

Los componentes on-premise demostrarán al menos 24 horas continuas de
autonomía degradada, conforme al piso contractual. Además se verificará
la sincronización automática en no más de 30 minutos tras el escenario
de 8 horas sin enlace, sin pérdida de ventas o documentos y con
resolución determinista. El alcance mantendrá excluida la apertura de
tarjetas y ampliación de cupos sin conexión; la compra con cupo vigente
requerirá topes, vigencia local, controles y factibilidad autorizados
por la filial, o se declarará no disponible con el procedimiento
alternativo aplicable (Bases Administrativas, 2026, art. 16.4; Bases
Técnicas Transversales, 2026, RT-03.10 a RT-03.13; Bases Técnicas del
Caso 09, 2026, cap. 15).

El orden de construcción lo decidirán los dueños de producto de cada
gerencia dentro del alcance aprobado. La priorización considerará, en
este orden, el hito contractual al que contribuye la historia, las
dependencias técnicas que habilita, el riesgo que reduce y el valor para
la operación. Una ampliación real del alcance se tramitará mediante el
control integrado de cambios de la sección 6.1.3.6 antes de incorporarse
como trabajo comprometido. Si el hallazgo corresponde a una obligación
contractual ya exigible que el catálogo omitió, se completarán su
registro y trazabilidad sin tratarla como un alcance opcional o un
servicio adicional.

\subsubsection{6.2.3. Evolución de la arquitectura y gestión de deuda técnica}\label{6-2-3-evolucion-de-la-arquitectura-y-gestion-de-deuda-tecnica}

La arquitectura de la solución no se diseñará completa al inicio ni se dejará emerger sin control. Se definirá una arquitectura base, que se aprobará en el hito H2, en el mes 4, junto con el plan de seguridad y el modelo de datos; desde ahí evolucionará de forma controlada a medida que los equipos integren las plataformas conservadas, sustituyan las capacidades previstas y construyan los nuevos servicios. Esta sección explica cómo se gobierna esa evolución, cómo se refactoriza el código y cómo se controla la deuda técnica.

\paragraph{6.2.3.1 Arquitectura evolutiva}\label{6-2-3-1-arquitectura-evolutiva}

La arquitectura base del Subdocumento 4 delimita servicios y autoridades de Retail y Emisor y el control de cruces X-01, con repositorios, identidades y claves separados. Los servicios se identifican con R:M para Mercadería, R:V para Venta, R:CL para Relación con el cliente y F:C para Crédito. Los habilitadores de identidad, integración, observabilidad, ambientes y entrega se construirán antes de los servicios que dependan de ellos. X-01 gobierna finalidades y políticas aplicadas en ambos extremos del intercambio; no constituye un repositorio unificado de clientes ni una vía informal de acceso a los datos financieros.

La nomenclatura y el emplazamiento se mantienen coherentes con el
alcance y la arquitectura de la propuesta. Las modificaciones se
registrarán y aprobarán por los mecanismos de 6.1.3.6, manteniendo la
coherencia del alcance y la arquitectura entre documentos.

La transición distinguirá el destino de cada plataforma. El sistema
empresarial y emisor tributario, el WMS principal y el marketplace
permanecen integrados; el núcleo Retail, el crédito de 2011 y el POS se
renuevan por etapas según el alcance. Comercio electrónico y
fidelización se mantienen e integran en el alcance declarado. Una
sustitución exigiría evidencia y una solicitud formal de cambio
aprobada; no se incorpora como decisión discrecional del equipo durante
el desarrollo.

Cuando se sustituya una capacidad, se aplicará modernización gradual
mediante el patrón estrangulador (Fowler, 2004). Cada nuevo servicio se
habilitará con pruebas de equivalencia, integridad, conciliación,
continuidad y retorno. Las catorce interfaces actuales se inventariarán
y recibirán un tratamiento individual de sustitución, encapsulación o
adaptación; una interfaz se retirará solo tras comprobar el flujo que la
reemplaza y sus dependencias.

Toda decisión que afecte a más de un equipo, a la frontera
Retail--Emisor o a un requerimiento no funcional se registrará como una
decisión de arquitectura (ADR), con su contexto, las alternativas
evaluadas, la decisión y sus consecuencias (Nygard, 2011), y será aprobada por el Comité de Arquitectura. Cuando la decisión habilite un cruce entre Retail y Emisor, el Comité solo aprobará su implementación técnica, después de la aprobación del cruce por la ruta única de la sección 6.1.2. Las funciones de aptitud de la
sección 6.2.2 protegen esas decisiones. Además de las metas de
desempeño, el pipeline verificará reglas estructurales, como que ningún
servicio del retail acceda directamente a datos del Emisor o que un
cambio en el contrato de una interfaz no rompa a los servicios que la
consumen.

La aprobación técnica de una ADR no reemplazará una solicitud de cambio
cuando afecte alcance, línea base u obligación contractual. Los
contratos y esquemas tendrán versiones y pruebas de compatibilidad,
además de procedimiento de retorno. Las decisiones sustituidas se
conservarán con su relación a la nueva decisión (Nygard, 2011).

\paragraph{6.2.3.2 Refactorización}\label{6-2-3-2-refactorizacion}

La refactorización será parte del trabajo diario y no una fase aparte. Cada historia incluye mejorar el código que modifica, y la definición de terminado exige justificar el efecto del cambio sobre duplicación y complejidad, comprobar las pruebas y cumplir las puertas acordadas. Un aumento localizado se evaluará por su necesidad y efecto global, sin autorizar incumplimientos de seguridad o del contrato. Las refactorizaciones mayores, como dividir un servicio o cambiar el modelo de un evento, se planificarán como historias habilitadoras al preparar cada intervalo de planificación y requerirán una ADR.

La refactorización segura se apoyará en pruebas automatizadas de
regresión y de contrato. El flujo de integración bloqueará cambios que
incumplan la cobertura de al menos 70 \% de la lógica de negocio. Como
objetivo adicional de calidad se aplicará al menos 80 \% sobre el código
nuevo de lógica de negocio, medido por separado. La herramienta y la
regla de cobertura identificarán líneas o ramas ejecutables, alcance de
medición y exclusiones justificadas; un buen resultado del código nuevo
no ocultará incumplimientos del mínimo contractual (Bases Técnicas
Transversales, 2026, RT-04.11).

\paragraph{6.2.3.3 Deuda técnica}\label{6-2-3-3-deuda-tecnica}

La deuda técnica se gestionará como un pasivo visible. Todo atajo, defecto de diseño o dependencia obsoleta que se detecte se registrará en el backlog con una etiqueta de deuda, su origen, su impacto y una fecha estimada de pago. Para clasificarla se usarán los cuadrantes de Fowler (2009), que distinguen si la deuda se asumió de forma deliberada o inadvertida, y si fue prudente o imprudente. La Figura 6.8 resume el tratamiento de cada clase.

\begin{center}
\begin{minipage}{\linewidth}
\centering
\ossFigCaption{Cuadrantes de la deuda técnica.}\label{fig:6-8}

\includegraphics[width=\linewidth,height=0.66\textheight,keepaspectratio]{figuras/recepcion-20261025/8bd9e2a8c5c5.png}

\raggedright\textbf{Fuente:} Elaboración propia a partir de Fowler (2009).
\end{minipage}
\end{center}

La Figura 6.8 distingue intención y prudencia, sin equiparar antigüedad
y deuda. Los ejemplos siguientes ilustran ambas clases de deuda.
Omitir pruebas para alcanzar una fecha es un atajo deliberado e
imprudente; duplicar una regla de cálculo en dos servicios es un defecto
inadvertido que debe corregirse según su riesgo y las pruebas de
regresión. La coexistencia planificada con una plataforma conservada se
registra como transición y no se presume deuda.

La política conservará tres reglas de control. Primero, se reservará 15 \% de capacidad para deuda priorizada en la construcción ordinaria. Las reservas se calculan sobre el 100 \% de la capacidad nominal de cada equipo en el sprint. En construcción ordinaria, el 15 \% se destina a deuda técnica priorizada y el 85 \% restante a construcción, pruebas, calidad, seguridad, coordinación y soporte. En los equipos involucrados en una marcha blanca, el 30 \% se destina a estabilización, que incluye el pago de deuda técnica, y el 70 \% restante a las demás actividades. Las dos reservas no se suman: la reserva máxima de un equipo en un sprint es 30 \%, y los equipos que no participan en la marcha blanca mantienen la de 15 \%.

Segundo, se aplicarán umbrales de razón de deuda técnica no superior a 5
\% y duplicación no superior a 3 \%. La razón de deuda se medirá como
esfuerzo estimado de remediación respecto de esfuerzo de desarrollo del
código evaluado, con las reglas de la herramienta documentadas; la
duplicación identificará su unidad, población y exclusiones. El plan de
calidad fijará el método y evidencia antes de aplicar el bloqueo. Estos
porcentajes constituyen política interna de Only Simple Solutions y se
medirán con el método definido en el plan de calidad.

Tercero, no se liberará un cambio con incumplimientos obligatorios de
seguridad, segregación o protección de datos. Las condiciones de
aceptación mantendrán ausencia de hallazgos críticos y altos abiertos.
El registro de deuda distinguirá los defectos de diseño de la solución
de las tareas de transición; la antigüedad de una plataforma o su
coexistencia temporal no la convierten automáticamente en deuda que deba
retirarse (Fowler, 2009).

La Figura 6.9 resume las tres reglas y sus distintos tipos de evidencia.

\begin{center}
\begin{minipage}{\linewidth}
\centering
\ossFigCaption{Reglas de la política de deuda técnica.}\label{fig:6-9}

\includegraphics[width=\linewidth,height=0.66\textheight,keepaspectratio]{figuras/recepcion-20261025/cf22bae951d1.png}

\raggedright\textbf{Fuente:} Elaboración propia.
\end{minipage}
\end{center}

La Figura 6.9 separa capacidad, medición y condiciones obligatorias. Los
porcentajes propios requieren nivelación y un método de cálculo; la
regla de seguridad y segregación se aplica a cada liberación. El estado
de la deuda se revisará al cierre de cada intervalo de planificación y se
informará al Comité de Arquitectura, conservando responsable y fecha de
tratamiento.

\subsubsection{6.2.4. DevSecOps e infraestructura como código}\label{6-2-4-devsecops-e-infraestructura-como-codigo}

Only Simple Solutions aplicará DevSecOps para que desarrollo, seguridad y operación compartan un flujo automatizado desde la integración del cambio hasta su funcionamiento. El despliegue seguirá el Subdocumento 4, con servicios nuevos de Retail y Emisor segregados en Azure y componentes locales para continuidad e integración. La recuperación y el tratamiento de información financiera se implementarán conforme a la arquitectura aprobada, preservando las condiciones de residencia, custodia, segregación y conmutación aplicables.

La Figura 6.10 muestra el flujo de entrega. Cada cambio pasa por cinco
etapas, código, construcción, QA, preproducción y producción, y cada
etapa termina en una puerta de calidad que la versión debe superar para
avanzar. La compilación y los controles de integración precederán a las
promociones; ninguna persona escribirá directamente sobre la rama
principal.

\begin{center}
\begin{minipage}{\linewidth}
\centering
\ossFigCaption{Flujo de entrega DevSecOps.}\label{fig:6-10}

\includegraphics[width=\linewidth,height=0.66\textheight,keepaspectratio]{figuras/recepcion-20261025/ccd38c5c8b72.png}

\raggedright\textbf{Fuente:} Elaboración propia.
\end{minipage}
\end{center}

La seguridad se incorporará desde el diseño. El modelado de amenazas
abarcará cada componente e integración externa y se actualizará ante
cambios de arquitectura, con revisión reforzada de la frontera. La
integración ejecutará SAST, análisis de composición, detección de
secretos y escaneo de imágenes; QA añadirá verificación dinámica e
integración. Cada versión liberada tendrá SBOM CycloneDX o SPDX y
artefactos firmados con procedencia verificada conforme a SLSA nivel 3 o
superior. Nuevas dependencias requerirán aprobación registrada de
licencia, mantenimiento, vulnerabilidades y procedencia (Bases Técnicas
Transversales, 2026, RT-04.05 y RT-11.22 a RT-11.27).

Las pruebas de intrusión se realizarán por un tercero independiente
antes de cada paso a producción y anualmente en operación. Los hallazgos
críticos y altos abiertos impedirán cumplir la condición de salida; los
plazos de remediación no son una dispensa para liberarlos en producción.
Calidad y seguridad conservarán informes, resultados y cierre de
observaciones (Bases Técnicas Transversales, 2026, RT-11.20 y núm.
20.1).

El dominio financiero tendrá repositorios, identidades de servicio,
permisos y secretos segregados de Retail. La custodia de claves se
resolverá con los mecanismos aprobados por clasificación, residencia y
separación de funciones, utilizando HSM cuando corresponda y sin
embebidos en código o imágenes. Los ambientes no productivos utilizarán
datos sintéticos o anonimización/seudonimización verificable; la sola
etiqueta de enmascaramiento no acreditará ese control. Ningún
desarrollador tendrá acceso interactivo ordinario a producción;
cualquier acceso excepcional será aprobado, temporal, registrado y con
sesión grabada. La recuperación conservará esos controles en su
emplazamiento secundario.

La infraestructura se definirá como código revisado y versionado, con
Terraform para los recursos seleccionados de Azure según el Subdocumento
4. Los mecanismos de configuración y despliegue local se elegirán por
compatibilidad, seguridad y soporte; la selección se documentará por
componente en la arquitectura y en el Formulario T-11, con su función y
condiciones de soporte. Las tiendas tendrán configuraciones
reproducibles y controles de desviación; el transporte GitOps que
corresponda sincronizará configuración, mientras la replicación y
conciliación de datos se comprobarán por sus procedimientos propios.

La configuración por ambiente permanecerá externalizada. El pipeline
construirá una vez el artefacto, registrará su huella, firma y
resultados y promoverá ese mismo artefacto sin recompilarlo. Los
secretos residirán en gestores con rotación automática y auditoría de
acceso; ninguna credencial se incorporará a código, imágenes o archivos
de configuración. Las migraciones de esquema serán versionadas,
compatibles con coexistencia de versiones y reversibles de forma
automatizada; sus efectos sobre datos se ensayarán antes de producción.

La infraestructura como código permitirá reconstruir los ambientes. H3, en el mes 6, entregará Desarrollo, QA, Preproducción, Producción y Recuperación ante
Desastres habilitados y operativos. Preproducción reproducirá topología,
configuración, versiones y volumen representativo de producción.
Recuperación mantendrá la replicación y el procedimiento de conmutación
y retorno probados antes de producción y al menos semestralmente durante
operación; no es un quinto paso ordinario de promoción de código. La
Figura 6.11 muestra los cinco ambientes y distingue su función (Bases
Técnicas Transversales, 2026, RT-04.01 y núm. 4.1).

\begin{center}
\begin{minipage}{\linewidth}
\centering
\ossFigCaption{Ambientes del proyecto.}\label{fig:6-11}

\includegraphics[width=\linewidth,height=0.66\textheight,keepaspectratio]{figuras/recepcion-20261025/5664879c53b5.png}

\raggedright\textbf{Fuente:} Elaboración propia.
\end{minipage}
\end{center}

Los controles automáticos determinarán la preparación técnica de una
versión. Antes de Preproducción o Producción se comprobarán las
autorizaciones que correspondan y los criterios de salida; el paso
productivo exigirá conformidad contractual y ausencia de congelamiento.
La autorización humana se registrará y la ejecución técnica será
automatizada y reproducible. Una aprobación no eliminará los resultados
obligatorios de las puertas.

\subsubsection{6.2.5. Entrega continua y automatización de pruebas}\label{6-2-5-entrega-continua-y-automatizacion-de-pruebas}

La entrega continua busca que el software esté siempre en condiciones de ser liberado, aunque se libere solo cuando el negocio lo decide (Humble, s.f.). Para lograrlo, los equipos trabajarán con desarrollo basado en tronco, con ramas de corta duración, de uno o dos días, que se integran a la rama principal después de una revisión por pares y de superar el pipeline. Las funcionalidades que aún no deben estar disponibles para los usuarios se desplegarán desactivadas mediante interruptores de funcionalidad, de modo que el despliegue técnico queda separado de la liberación al negocio.

Esta separación permite verificar la Etapa 2 en ambientes no productivos
mientras la Etapa 1 está en marcha blanca. Los interruptores de
funcionalidad mantendrán inactivas capacidades que no deban usarse
todavía, sin conservar ramas paralelas. Su activación con efecto sobre
transacciones reales se evaluará como cambio productivo, con controles
de datos, seguridad, conformidad y calendario; la validación funcional
del dueño de producto no autorizará por sí sola el paso ni una excepción
al congelamiento.

Las pruebas combinarán una base de verificaciones unitarias, pruebas de
integración y contratos y verificaciones de extremo a extremo de los
flujos críticos. Estos incluyen venta con crédito, reservas, precio,
conteo, operación sin enlace y conciliación, control de cruces X-01 y
migración de cartera. Se mantendrá el bloqueo de cobertura de lógica de
negocio de 70 \% y se medirá separadamente el objetivo de 80 \% sobre
código nuevo.

Las campañas de carga se realizarán a 1,5 veces el peak declarado y
comprobarán los umbrales de experiencia del usuario; se ejecutarán
ensayos de resiliencia, recuperación, accesibilidad y seguridad antes de
cada producción, junto con las pruebas de aceptación y los dos ensayos
previos de migración. El Plan de Pruebas y Validación del Subdocumento 9
y T-13 reunirán niveles, datos, ambientes, calendario y criterios de
salida; cada resultado quedará relacionado con requisito y versión
(Bases Técnicas Transversales, 2026, núm. 20.1).

La Figura 6.12 resume las puertas de calidad del flujo de entrega y el
criterio que una versión debe cumplir para avanzar a la etapa siguiente.

\begin{center}
\begin{minipage}{\linewidth}
\centering
\ossFigCaption{Criterios de las puertas de calidad.}\label{fig:6-12}

\includegraphics[width=\linewidth,height=0.66\textheight,keepaspectratio]{figuras/recepcion-20261025/2e57db6c8171.png}

\raggedright\textbf{Fuente:} Elaboración propia.
\end{minipage}
\end{center}

Las puertas convierten la definición de terminado en controles
repetibles con resultados por versión. La evidencia del pipeline
respaldará QA y certificación, pero Calidad consolidará además los
informes, resultados independientes, trazabilidad, observaciones
resueltas y conformidades que corresponden a cada hito. La aceptación
requerirá el acta de la contraparte técnica y no se producirá
automáticamente porque la ejecución CI sea satisfactoria.

La liberación a producción será progresiva. Los servicios en nube
utilizarán despliegue azul-verde; los componentes de tienda pasarán por
el piloto controlado de tres tiendas definido en el alcance y después
por olas, con criterios de avance y detención medidos. Preproducción
ensayará la reversión automatizada de aplicación, configuración y
esquemas compatibles, incluyendo recuperación y conciliación de datos
cuando el cambio los afecte. No se liberará un cambio que carezca de
reversibilidad demostrada.

Si el canario o una ola incumplen sus umbrales, se detendrá el avance y
se ejecutará el retorno previsto, conservando transacciones y
evidencias. Una recuperación extraordinaria no reemplazará la
automatización comprometida ni dispensará las condiciones del contrato.
La estrategia por sitio y sus tiempos quedarán documentados en el plan
de implantación y T-18.

\subsubsection{6.2.6. Organización y decisiones del equipo de desarrollo}\label{6-2-6-organizacion-y-decisiones-del-equipo-de-desarrollo}

Los equipos de desarrollo serán estables durante la implementación y se organizarán según los dominios de la arquitectura, para que cada uno sea dueño de un conjunto de módulos de principio a fin. Lo diseña, lo construye, lo prueba y acompaña su operación. Esta organización reduce las dependencias entre equipos y hace coincidir la estructura de los equipos con la de la solución.

La organización distinguirá ámbitos de Retail, Emisor y plataforma.
Retail cubrirá los nueve servicios de Mercadería, Venta y Relación con
el cliente; Emisor, F:C-01 originación, F:C-02 cartera y F:C-03
evidencia financiera; plataforma habilitará integración, identidad,
observabilidad, ambientes y control de cruces X-01 con las validaciones
de datos y cumplimiento. La composición y dedicación de los equipos se
documentarán con sus productos, dependencias y capacidad en el
Subdocumento 12 y en los Formularios T-8 y T-15.

Cada producto tendrá un dueño de producto identificado por Ancoa para
ordenar su backlog dentro del alcance aprobado. Los servicios técnicos
se relacionarán con esos productos, sin crear automáticamente un dueño
diferente para cada microservicio. Cada equipo Scrum incluirá un product
owner, un Scrum Master y developers y tendrá capacidad para producir un
incremento conforme a la definición de terminado. Los equipos que
trabajen sobre el mismo producto compartirán objetivo del producto,
backlog y product owner. Cualquier facilitación compartida tendrá
dedicación explícita y no duplicará horas. Los líderes de desarrollo,
arquitectura, calidad, seguridad, datos e integración conservarán sus
atribuciones de coordinación y control; las solicitudes de habilitación
e incidencias podrán seguir un flujo Kanban separado.

\paragraph{6.2.6.1 Ceremonias, artefactos y cadencias}\label{6-2-6-1-ceremonias-artefactos-y-cadencias}

Los equipos seguirán los eventos de Scrum y sus compromisos de producto, sprint y terminado (Schwaber y Sutherland, 2020). La coordinación conjunta de ocho semanas utilizará planificación de dependencias y revisión de resultados; se distinguen esas prácticas seleccionadas de los eventos propios de Scrum y de los comités contractuales. La Figura 6.13 representa cuatro sprints y la interacción de ambos niveles, sin equipararlos a dos meses calendario.

\begin{center}
\begin{minipage}{\linewidth}
\centering
\ossFigCaption{Cadencia de un intervalo de planificación.}\label{fig:6-13}

\includegraphics[width=\linewidth,height=0.66\textheight,keepaspectratio]{figuras/recepcion-20261025/579a3cae8227.png}

\raggedright\textbf{Fuente:} Elaboración propia.
\end{minipage}
\end{center}

La Figura 6.13 muestra cuatro sprints consecutivos y coordinaciones
adicionales dentro de ocho semanas. Sus símbolos de gobierno son
ilustrativos; las fechas mensuales y quincenales se fijarán en el
calendario contractual, sin deducirlas de una equivalencia entre semanas
y meses.

La Figura 6.14 detalla cada ceremonia, su cadencia, sus participantes y
el resultado que produce.

\begin{center}
\begin{minipage}{\linewidth}
\centering
\ossFigCaption{Ceremonias y cadencias del desarrollo.}\label{fig:6-14}

\includegraphics[width=\linewidth,height=0.66\textheight,keepaspectratio]{figuras/recepcion-20261025/722588ff4876.png}

\raggedright\textbf{Fuente:} Elaboración propia.
\end{minipage}
\end{center}

Las ceremonias de cada equipo sostienen el ritmo de entrega cada dos
semanas; las del intervalo de planificación aseguran que los equipos avancen hacia los
mismos objetivos y que las dependencias, en especial las de la
plataforma y la frontera de datos, se resuelvan antes de bloquear el
trabajo. La revisión del sprint no reemplaza la aceptación contractual,
que se formaliza en las actas de los hitos conforme a la sección 6.1.2.
Entre ceremonias, la coordinación cotidiana de los equipos se realizará
por Slack; todo acuerdo que afecte tareas, responsabilidades,
dependencias, plazos o decisiones se formalizará en Jira o SharePoint,
según corresponda, para que no quede solo registrado en una
conversación.

Los artefactos del desarrollo serán el backlog del producto, el backlog
del sprint y el incremento, con objetivo del producto, objetivo del
sprint y definición de terminado. Se añadirán objetivos del intervalo de planificación,
tablero de dependencias, ADR, registro de deuda y
trazabilidad. Jira mantendrá trabajo y relaciones; el repositorio,
código y artefactos técnicos; SharePoint, actas y documentación
controlada. Las configuraciones y resultados de esas relaciones se
comprobarán durante la habilitación del flujo de entrega.

\paragraph{6.2.6.2 Mecanismos de decisión}\label{6-2-6-2-mecanismos-de-decision}

Las decisiones del desarrollo se tomarán en el nivel que disponga de información y atribuciones para resolverlas. La Figura 6.15 distingue aprobación del cambio, validación de datos y seguridad, conformidad productiva y autonomía técnica. La decisión quedará registrada y se escalará cuando exceda la autoridad; ninguna aprobación de detalle sustituirá el control integrado de cambios.

\begin{center}
\begin{minipage}{\linewidth}
\centering
\ossFigCaption{Mecanismos de decisión del desarrollo.}\label{fig:6-15}

\includegraphics[width=\linewidth,height=0.66\textheight,keepaspectratio]{figuras/recepcion-20261025/89d8043bc9be.png}

\raggedright\textbf{Fuente:} Elaboración propia.
\end{minipage}
\end{center}

El diagrama distingue entre las decisiones técnicas, que los equipos
toman con autonomía, y las que afectan compromisos contractuales, la
frontera entre negocios o la operación de Ancoa, que requieren una
instancia formal. Ningún rol ágil, ni el product owner ni el Scrum
Master, tiene atribución para modificar el alcance o las fechas del
contrato.

La Figura 6.15 identifica la instancia principal para resolver cada
materia; sus rutas no excluyen otros controles aplicables al mismo
cambio. Una modificación puede afectar a la vez el alcance, el cruce de
datos y la producción. En ese caso acumulará la aprobación del cambio,
las validaciones de datos, seguridad y cumplimiento y la conformidad
productiva antes de ejecutarse.

\subsubsection{6.2.7. Coordinación entre etapas y tiempo de entrega}\label{6-2-7-coordinacion-entre-etapas-y-tiempo-de-entrega}

En este proyecto, el tiempo de salida al mercado depende de preparar y poner cada capacidad en operación dentro del calendario contractual y comercial. La Etapa 1 pasa a producción en el mes 16 y la Etapa 2 en el 21. El plan incorporará las cinco ventanas de congelamiento, la fecha variable del evento anual y los dos períodos simultáneos de construcción, marcha blanca y producción.

El calendario de trabajo utilizará fechas, no una suma simple de
duraciones comerciales. La unión de ventanas evitará contar dos veces
eventos contenidos en una misma campaña y comprobará los cruces entre
períodos protegidos. El cronograma integrado relacionará los meses
contractuales con fechas calendario desde el acta de inicio, sin alterar
los hitos obligatorios. La Figura 6.16 representa las fases y
convivencias del contrato; el cronograma integrado incorporará las
ventanas comerciales vigentes.

La Figura 6.16 sintetiza las fases por meses y sus convivencias. La
ubicación por fechas de los intervalos de planificación, hitos E-25 y
ventanas comerciales se desarrollará en el cronograma integrado; la
figura no representa ese detalle.

\begin{center}
\begin{minipage}{\linewidth}
\centering
\ossFigCaption{Calendario de entregas del desarrollo.}\label{fig:6-16}

\includegraphics[width=\linewidth,height=0.66\textheight,keepaspectratio]{figuras/recepcion-20261025/dd39960c0ee2.png}

\raggedright\textbf{Fuente:} Elaboración propia.
\end{minipage}
\end{center}

El calendario exige programar hacia atrás desde cada paso contractual y
las ventanas disponibles. Las versiones, pruebas, documentación y
conformidades estarán preparadas con la holgura justificada en el
cronograma integrado. La fecha de Preproducción se determinará en el
cronograma integrado a partir de dependencias, pruebas y ventanas
disponibles. Durante los congelamientos no se programarán cambios
productivos ni activaciones que los eludan; el diagnóstico, desarrollo y
ensayo en ambientes aislados se coordinarán sin afectar la operación
protegida. Una urgencia no creará por sí misma autorización para
intervenir sistemas.

La coordinación entre etapas se resolverá en la programación de
dependencias y en las revisiones conjuntas. Entre los meses 13 y 15 se
sostendrá la marcha blanca de Etapa 1 en paralelo al desarrollo de Etapa
2. Entre 19 y 20, Etapa 1 estará en producción mientras Etapa 2 cumple
su marcha blanca; se controlarán contratos, esquemas y versiones
compatibles, autoridad única de datos compartidos, conciliación,
incidentes y retorno. El paso del mes 21 preservará disponibilidad,
desempeño e integridad de Etapa 1.

Los recursos se asignarán por frente y función, diferenciando
estabilización, construcción, pruebas y atención productiva. La
aprobación de detalle y línea base de Etapa 2 seguirá H8 y las demás
condiciones del Formulario E-25, sin inferir fechas de una equivalencia
entre sprints y meses. Las siguientes cuatro prácticas apoyarán esa
coordinación.

\begin{itemize}
\item
  \textbf{Entrega continua hasta Preproducción.} El software se integra
  y se prueba todos los días, aunque se libere a producción solo en las
  ventanas disponibles; así, cuando se abre una ventana, la versión ya
  está verificada.
\item
  \textbf{Despliegue separado de la liberación.} Los interruptores
  mantendrán inactivas funciones verificadas hasta contar con
  validación, autorización y ventana. Su activación con efectos reales
  respetará el mismo gobierno productivo.
\item
  \textbf{Capacidad reservada para la estabilización.} Se reservará 30 \% de capacidad para los equipos involucrados en cada marcha blanca. Esa reserva incluye el pago de deuda técnica y no se suma al 15 \% de la construcción ordinaria (sección 6.2.3.3); el 70 \% restante cubre construcción, pruebas, calidad, seguridad, coordinación y soporte. El T-15 y el cronograma reflejarán el reparto sin desplazar hitos.
\item
  \textbf{Lotes pequeños.} Las historias se dimensionan para completarse
  dentro de un sprint, lo que reduce el riesgo de cada liberación y
  facilita revertirla.
\end{itemize}

El desempeño de entrega se medirá por servicio con cuatro indicadores:
frecuencia de despliegue, tiempo de entrega, tasa de cambios fallidos y
tiempo de restauración. La frecuencia contabilizará despliegues reales
en Producción por período; el tiempo de entrega transcurrirá desde el
commit hasta el despliegue productivo, incluidas esperas; la tasa de
cambios fallidos relacionará despliegues que requieren intervención con
despliegues totales; y la restauración medirá el retorno del servicio
tras un cambio fallido. La Figura 6.17 distingue esos resultados de las
métricas internas de QA y Preproducción (DORA, s.f.; Bases Técnicas
Transversales, 2026, RT-04.12).

Los objetivos de frecuencia y tiempo hasta Producción se formularán por
servicio en el calendario de liberaciones autorizadas, con su período y
forma de cómputo; una meta no habilitará cambios durante congelamientos.
Se aplicarán como objetivos internos una tasa de cambios fallidos no
superior a 15 \% y una restauración de cambios fallidos en hasta una
hora. Esta restauración no sustituye el RTO/RPO de recuperación ante
desastres. Como indicadores complementarios se conservarán la
integración a QA al menos diaria por equipo y el objetivo interno de
hasta cinco días hábiles para verificar un cambio en Preproducción.

\begin{center}
\begin{minipage}{\linewidth}
\centering
\ossFigCaption{Indicadores de tiempo de entrega.}\label{fig:6-17}

\includegraphics[width=\linewidth,height=0.66\textheight,keepaspectratio]{figuras/recepcion-20261025/baf289d8c137.png}

\raggedright\textbf{Fuente:} Elaboración propia.
\end{minipage}
\end{center}

Los indicadores se revisarán por intervalo de planificación y en el
informe mensual, y continuarán midiéndose durante la Operación desde el
mes 21 al 56. Sus resultados conservarán versión, servicio, ventana,
numerador, denominador y fuente. Cuando no existan despliegues o fallos
en un período se registrará la ausencia de población, sin inventar tasas
o tiempos. Las desviaciones sostenidas activarán análisis y medidas con
responsable y plazo; las que afecten obligaciones, hitos o línea base
seguirán el escalamiento y control de cambios.

\section{Referencias}\label{referencias}

Las siguientes fuentes respaldan los procesos, controles y restricciones
citados en este capítulo.

\noindent\hangindent=1.5em\hangafter=1 Bases Administrativas para la Preparación de la Propuesta. (2026). \emph{Licitación N.º TFEP-01/2026}, versión 1.0.

\noindent\hangindent=1.5em\hangafter=1 Bases Técnicas del Caso 09 Cadena Multitienda. (2026). \emph{Licitación N.º TFEP-01/2026}, versión 1.0.

\noindent\hangindent=1.5em\hangafter=1 Bases Técnicas Transversales. (2026). \emph{Licitación N.º TFEP-01/2026}, versión 1.0.

\noindent\hangindent=1.5em\hangafter=1 Cucumber. (s.f.-a). \emph{Gherkin reference}. \url{https://cucumber.io/docs/gherkin/reference/}

\noindent\hangindent=1.5em\hangafter=1 Cucumber. (s.f.-b). \emph{Step definitions}. \url{https://cucumber.io/docs/cucumber/step-definitions/}

\noindent\hangindent=1.5em\hangafter=1 DORA. (s.f.). \emph{DORA's software delivery performance metrics}. \url{https://dora.dev/guides/dora-metrics/}

\noindent\hangindent=1.5em\hangafter=1 Ford, N., Parsons, R., y Kua, P. (2017). \emph{Building evolutionary architectures: Support constant change}. O'Reilly Media.

\noindent\hangindent=1.5em\hangafter=1 Fowler, M. (2004, 29 de junio). \emph{Original strangler fig application}. \url{https://martinfowler.com/bliki/OriginalStranglerFigApplication.html}

\noindent\hangindent=1.5em\hangafter=1 Fowler, M. (2009, 14 de octubre). \emph{Technical debt quadrant}. \url{https://martinfowler.com/bliki/TechnicalDebtQuadrant.html}

\noindent\hangindent=1.5em\hangafter=1 Humble, J. (s.f.). \emph{What is continuous delivery?} Continuous Delivery. \url{https://continuousdelivery.com/}

\noindent\hangindent=1.5em\hangafter=1 Kanban Guides. (2025, mayo). \emph{The Kanban Guide} (v2025.5). \url{https://kanbanguides.org/the-kanban-guide/2025.5/}

\noindent\hangindent=1.5em\hangafter=1 Ley N.º 21.719. (2024). \emph{Regula la protección y el tratamiento de los datos personales y crea la Agencia de Protección de Datos Personales}. Biblioteca del Congreso Nacional de Chile. \url{https://www.bcn.cl/leychile/Navegar?idNorma=1209272}

\noindent\hangindent=1.5em\hangafter=1 Nygard, M. (2011, 15 de noviembre). \emph{Documenting architecture decisions}. Cognitect. \url{https://www.cognitect.com/blog/2011/11/15/documenting-architecture-decisions}

\noindent\hangindent=1.5em\hangafter=1 Only Simple Solutions. (2026). \emph{Catálogo de requerimientos atomizados y depurados} (v3.0). RequerimientosAtomizados\_Depuracion\_Alcance.xlsx.

\noindent\hangindent=1.5em\hangafter=1 Prasad, V., Rajkumar, P., y Rastogi, S. C. (2006). \emph{Managing firmed fixed price projects using EVM: A case-study}. Project Management Institute. \url{https://www.pmi.org/learning/library/managing-fixed-price-projects-evm-8142}

\noindent\hangindent=1.5em\hangafter=1 Project Management Institute {[}PMI{]}. (2017). \emph{A guide to the project management body of knowledge (PMBOK® guide)} (6.ª ed.). Project Management Institute.

\noindent\hangindent=1.5em\hangafter=1 Prosci. (s.f.). \emph{The Prosci ADKAR model}. \url{https://www.prosci.com/methodology/adkar}

\noindent\hangindent=1.5em\hangafter=1 Scaled Agile, Inc.~(2025, 25 de febrero). \emph{Planning interval}. \url{https://framework.scaledagile.com/planning-interval}

\noindent\hangindent=1.5em\hangafter=1 Schwaber, K., y Sutherland, J. (2020, noviembre). \emph{The Scrum Guide}. Scrum Guides. \url{https://scrumguides.org/scrum-guide.html}

\section{DECLARACIÓN USO IA}\label{declaracion-de-uso-de-ia}

El equipo declara que el uso de herramientas de inteligencia artificial generativa en este subdocumento se limitó a las actividades indicadas en la siguiente tabla. Todo contenido, dato, cálculo, decisión de diseño, diagrama, conclusión y referencia fue revisado y verificado por las personas identificadas, quienes asumen como propio el resultado final y la responsabilidad íntegra sobre su exactitud, originalidad y coherencia técnica.

\begin{longtable}[]{@{}
  >{\raggedright\arraybackslash}p{(\columnwidth - 12\tabcolsep) * \real{0.1100}}
  >{\raggedright\arraybackslash}p{(\columnwidth - 12\tabcolsep) * \real{0.1700}}
  >{\raggedright\arraybackslash}p{(\columnwidth - 12\tabcolsep) * \real{0.3000}}
  >{\raggedright\arraybackslash}p{(\columnwidth - 12\tabcolsep) * \real{0.1300}}
  >{\raggedright\arraybackslash}p{(\columnwidth - 12\tabcolsep) * \real{0.1400}}
  >{\raggedright\arraybackslash}p{(\columnwidth - 12\tabcolsep) * \real{0.1500}}@{}}
\caption{Declaración de uso de IA}\label{tab:declaracion-uso-ia}\\
\toprule\noalign{}
\begin{minipage}[b]{\linewidth}\raggedright
Sección
\end{minipage} & \begin{minipage}[b]{\linewidth}\raggedright
Herramienta
\end{minipage} & \begin{minipage}[b]{\linewidth}\raggedright
Finalidad de uso
\end{minipage} & \begin{minipage}[b]{\linewidth}\raggedright
Nivel en texto
\end{minipage} & \begin{minipage}[b]{\linewidth}\raggedright
Nivel en diagramas
\end{minipage} & \begin{minipage}[b]{\linewidth}\raggedright
Revisión humana
\end{minipage} \\
\midrule\noalign{}
\endfirsthead
\toprule\noalign{}
\begin{minipage}[b]{\linewidth}\raggedright
Sección
\end{minipage} & \begin{minipage}[b]{\linewidth}\raggedright
Herramienta
\end{minipage} & \begin{minipage}[b]{\linewidth}\raggedright
Finalidad de uso
\end{minipage} & \begin{minipage}[b]{\linewidth}\raggedright
Nivel en texto
\end{minipage} & \begin{minipage}[b]{\linewidth}\raggedright
Nivel en diagramas
\end{minipage} & \begin{minipage}[b]{\linewidth}\raggedright
Revisión humana
\end{minipage} \\
\midrule\noalign{}
\endhead
\bottomrule\noalign{}
\endlastfoot
6.1 & Codex GPT-6 Sol Alto / Claude Sonnet 5.5 & Se utilizó para organizar los subpuntos de la sección, generar figuras para complementar el texto, identificar referencias, identificar con qué otros subdocumentos y anexos se relaciona en la introducción, generar la redacción final y verificar inconsistencias con los otros subdocumentos de la licitación. & Alto & Alto & Daniel Cepeda \\
6.2 & Codex GPT-6 Sol Alto / Claude Sonnet 5.5 & Se utilizó para organizar los subpuntos de la sección, generar figuras para complementar el texto, identificar referencias, identificar con qué otros subdocumentos y anexos se relaciona en la introducción, generar la redacción final y verificar inconsistencias con los otros subdocumentos de la licitación. & Alto & Alto & Samira Becerra Saa \\
Anexo & Claude Opus 5.5 & Se utilizó para trasladar al anexo las planillas de la sección 6.1.4, darles el diseño de los anexos del Subdocumento 2 y alinear sus campos con el registro de interesados. & Alto & Ninguno & Samira Becerra Saa \\
Formulario T-09 & Claude Opus 5.5 & Se utilizó para generar el formulario en base a la sección 6.1 y para corregir las inconsistencias con el subdocumento y los otros formularios. & Alto & Ninguno & Samira Becerra Saa \\
Formulario T-10 & Claude Sonnet 5.5 & Se utilizó para generar el archivo en base al subdocumento 6 y las instrucciones del usuario. También, para buscar referencias. & Alto & Ninguno & Samira Becerra Saa \\
\end{longtable}

A continuación, se detalla la escala de niveles para completar la tabla anterior.

\begin{tabular}{@{}p{0.069\linewidth}p{0.429\linewidth}p{0.429\linewidth}@{}}
\toprule Nivel & Texto & Diagramas \\
\midrule
Ninguno & Sin uso de IA. & Sin uso de IA. \\
Bajo & Corrección ortográfica, de estilo o reformulación de frases escritas por el grupo. & Sugerencias de formato o disposición sobre un diagrama elaborado por el grupo. \\
Medio & Borradores o síntesis de partes que el grupo reescribió y verificó. & Diagrama generado con asistencia (por ejemplo, código Mermaid o PlantUML) a partir de un modelo definido por el grupo. \\
Alto & Texto generado o modificado sustancialmente por IA; la revisión humana se registra en su propia columna. & Diagrama generado sustancialmente por IA. \\
\bottomrule
\end{tabular}

\ossFinalPage

\end{document}
